\documentclass[a4paper,11pt]{article}
\usepackage[english]{babel}
\usepackage{amsmath,amssymb,graphicx,enumerate,bbm,stmaryrd,mathrsfs,latexsym,theorem}
%\usepackage{stmaryrd,bbm}
\usepackage[dvipsnames]{xcolor}
\usepackage[left= 2.5cm, bottom = 4cm, top = 3.5cm, right= 2.5cm]{geometry}
\usepackage[citecolor=blue,colorlinks=true]{hyperref}
%\usepackage[notref,notcite]{showkeys}
\usepackage{authblk} % for affiliations
\usepackage{tikz}

%\usepackage{refcheck}


%
%\newtheorem{theorem}{Theorem}[section]
%\newtheorem{lemma}[theorem]{Lemma}
%\newtheorem{prop}[theorem]{Proposition}
%\newtheorem{cor}[theorem]{Corollary}
%
%\theoremstyle{definition} 
%\newtheorem{definition}[theorem]{Definition}
%%\theoremstyle{remark}
%\newtheorem{rem}[theorem]{Remark}



%%%%%%%%%% Start TeXmacs macros
\renewcommand{\leq}{\leqslant}
\renewcommand{\geq}{\geqslant}
\newcommand{\assign}{:=}
\newcommand{\backassign}{=:}
\newcommand{\precprec}{\prec\!\!\!\prec}
\newcommand{\mathd}{\mathrm{d}}
\newcommand{\tmcolor}[2]{{\color{#1}{#2}}}
\newcommand{\tmop}[1]{\ensuremath{\operatorname{#1}}}
\newcommand{\tmrsup}[1]{\textsuperscript{#1}}
\newenvironment{enumeratealpha}{\begin{enumerate}[a{\textup{)}}] }{\end{enumerate}}
\newenvironment{proof}{\noindent\textbf{Proof\ }}{\hspace*{\fill}$\Box$\medskip}
\newtheorem{theorem}{Theorem}[section]
\newtheorem{lemma}[theorem]{Lemma}
\newtheorem{proposition}[theorem]{Proposition}
\newtheorem{corollary}[theorem]{Corollary}
{\theorembodyfont{\rmfamily}\newtheorem{remark}[theorem]{Remark}}

\numberwithin{equation}{section}
%%%%%%%%%% End TeXmacs macros

\iffalse
% old trees
\newcommand{\tthree}[1]{#1\!\tmrsup{\resizebox{.6em}{!}{\includegraphics{trees/3_tree.eps}}}}
\newcommand{\ttwor}[1]{#1\!\tmrsup{\resizebox{.6em}{!}{\includegraphics{trees/2r_tree_res.eps}}}}
\newcommand{\tthreer}[1]{#1\!\tmrsup{\!\resizebox{.6em}{!}{\includegraphics{trees/3r_tree_res.eps}}}}
\newcommand{\tthreeone}[1]{#1\!^{\!\resizebox{.6em}{!}{\includegraphics{trees/31_tree.eps}}}}
\newcommand{\ttwo}[1]{#1\!\tmrsup{\!\resizebox{.7em}{!}{\includegraphics{trees/2_tree.eps}}}}
\newcommand{\ttwothree}[1]{#1\!\tmrsup{\resizebox{.6em}{!}{\includegraphics{trees/23_tree.eps}}}}
\newcommand{\tthreethree}[1]{#1\!\tmrsup{\resizebox{.9em}{!}{\includegraphics{trees/33_tree.eps}}}}
\newcommand{\tthreetwo}[1]{#1\!\tmrsup{\resizebox{.9em}{!}{\includegraphics{trees/32_tree.eps}}}}
\newcommand{\tthreethreer}[1]{#1\!^{\resizebox{.9em}{!}{\includegraphics{trees/33_tree_res.eps}}}}
\newcommand{\tthreethreebr}[1]{#1\!\tmrsup{\resizebox{.9em}{!}{\includegraphics{trees/2a1_tree_res.eps}}}}
\newcommand{\ttwoone}[1]{#1\!^{\resizebox{.6em}{!}{\includegraphics{trees/21_tree.eps}}}}
\newcommand{\ttwothreer}[1]{#1\!^{\resizebox{.9em}{!}{\includegraphics{trees/23_tree_res.eps}}}}
\newcommand{\tthreetwor}[1]{#1\!^{\resizebox{.9em}{!}{\includegraphics{trees/32_tree_res.eps}}}}
\newcommand{\tone}[1]{#1\!\tmrsup{\resizebox{.4em}{!}{\includegraphics{trees/1_tree_mod.eps}}}}
\newcommand{\tzero}[1]{#1\!\tmrsup{{\varnothing}}}
\fi

% new trees
\newcommand{\tone}[1]{#1^{\resizebox{0.4em}{!}{\input{trees-tikz/1_tree_mod.tikz}}}}
\newcommand{\tzero}[1]{#1^{\varnothing}}
\newcommand{\tthree}[1]{#1^{\resizebox{0.6em}{!}{\input{trees-tikz/3_tree.tikz}}}}
\newcommand{\ttwor}[1]{#1\tmrsup{\resizebox{0.6em}{!}{\input{trees-tikz/2r_tree_res.tikz}}}}
\newcommand{\tthreer}[1]{#1\tmrsup{\!\resizebox{0.6em}{!}{\input{trees-tikz/3r_tree_res.tikz}}}}
\newcommand{\ttwo}[1]{#1^{\!\resizebox{0.6em}{!}{\input{trees-tikz/2_tree.tikz}}}}

\newcommand{\tthreeone}[1]{#1^{\!\resizebox{0.6em}{!}{\input{trees-tikz/31_tree.tikz}}}}
\newcommand{\ttwoone}[1]{#1^{\!\resizebox{0.6em}{!}{\input{trees-tikz/21_tree_mod.tikz}}}}

\newcommand{\tthreethree}[1]{#1\tmrsup{\!\resizebox{!}{.8em}{\input{trees-tikz/33_tree.tikz}}}}
\newcommand{\tthreetwo}[1]{#1\tmrsup{\!\resizebox{!}{.8em}{\input{trees-tikz/32_tree.tikz}}}}
\newcommand{\tthreethreer}[1]{#1^{\!\resizebox{!}{.8em}{\input{trees-tikz/33_tree_res.tikz}}}}
\newcommand{\tthreethreerprime}[1]{#1^{\prime\!\resizebox{!}{.8em}{\input{trees-tikz/33_tree_res.tikz}}}}
\newcommand{\ttwothreer}[1]{#1^{\!\resizebox{!}{.8em}{\input{trees-tikz/23_tree_res.tikz}}}}
\newcommand{\tthreetwor}[1]{#1^{\!\resizebox{!}{.8em}{\input{trees-tikz/32_tree_res.tikz}}}}




\newcommand{\CC}{\mathscr{C} \hspace{.1em}}
\newcommand{\CD}{D{\hspace{.1em}}}
\newcommand{\CF}{F{\hspace{.1em}}}
\newcommand{\CB}{B{\hspace{.1em}}}
\newcommand{\CA}{A{\hspace{.1em}}}
\newcommand{\DD}{D{\hspace{.1em}}}
\newcommand{\LL}{\mathscr{L}}
\newcommand{\UU}{\mathscr{U}}
\newcommand{\VV}{\mathscr{V}}
\newcommand{\Q}{\mathscr{Q} \hspace{.2em}}
\newcommand{\SSS}{S{\hspace{.1em}}}
\newcommand{\PP}{P{\hspace{.1em}}}
\newcommand{\XX}{X}

\newcommand{\rmk}[1]{\textcolor{red}{\framebox{#1}}}
\newcommand{\rmbb}[1]{\textcolor{black}{#1}}
\newcommand{\rmb}[1]{\textcolor{blue}{#1}}
\newcommand{\rmm}[1]{\textcolor{magenta}{#1}}
\newcommand{\rmg}[1]{\textcolor{Bittersweet}{#1}}

\newcommand{\R}{\mathbb{R}}
\newcommand{\Z}{\mathbb{Z}}
\newcommand{\N}{\mathbb{N}}
\newcommand{\Td}{{\mathbb{T}^{d}}}
\newcommand{\Prec}{\prec\!\!\!\prec}
\newcommand{\8}{\infty}
\newcommand{\dd}{\mathrm{d}}
\DeclareMathOperator{\esssup}{esssup}
%\newcommand{\Preceq}{\preccurlyeq\!\!\!\prec}
%

%


\newcommand{\axiomtext}{Assumption}


\setcounter{tocdepth}{2}

\title{Global solutions to elliptic\\ and parabolic $\Phi^4$ models in
Euclidean space.}


\author[1]{Massimiliano Gubinelli}
\author[2]{Martina Hofmanov\'a}
\affil[1]{\small Hausdorff Center for Mathematics\\ \& Institute for Applied Mathematics\\ University of Bonn\\
Endenicher Allee 60\\
53115 Bonn, Germany.  \href{mailto:gubinelli@iam.uni-bonn.de}{gubinelli@iam.uni-bonn.de} }
%\email
%
\affil[2]{\small Fakult\"at f\"ur Mathematik, Universit\"at Bielefeld, Postfach 10 01 31, 33501 Bielefeld, Germany. \href{mailto:hofmanova@math.uni-bielefeld.de}{hofmanova@math.uni-bielefeld.de}}
%\email{hofmanov@math.tu-berlin.de}


\begin{document}
\maketitle

\begin{abstract}
We prove existence of global  solutions to singular SPDEs on $\mathbb{R}^d$ with cubic nonlinearities and additive white noise perturbation, both in the elliptic setting in dimensions $d=4,5$ and in the parabolic setting for $d=2,3$.  We prove uniqueness and coming down from infinity for the parabolic equations. A  motivation for considering  these equations is the construction of scalar interacting Euclidean quantum field theories. The parabolic equations are related to the $\Phi^4_d$ Euclidean quantum field theory via  Parisi--Wu stochastic quantization, while the elliptic equations are linked to the $\Phi^4_{d-2}$ Euclidean quantum field theory via the Parisi--Sourlas dimensional reduction mechanism.  \\[.4em]
  
\noindent  \textbf{Keywords:} singular SPDEs, paracontrolled distributions, global solutions, stochastic quantization, dimensional reduction, semilinear elliptic equations.
\end{abstract}


\tableofcontents

\section{Introduction}


This paper is concerned with elliptic and parabolic partial differential equations related to the $\Phi^{4}$ Euclidean quantum field theory on the full space.
More precisely, we consider the following semilinear elliptic partial differential equation on $\R^{d}$ for $d=4,5$,
\begin{equation}\label{eq:phi4}
(-\Delta+\mu)\varphi+\varphi^{3}=\xi,
\end{equation}
where $\xi$ is a space white noise on $\R^{d}$ and $\mu>0$.
We also consider the Cauchy problem for the  semilinear parabolic partial differential equation on $\R_+\times\R^{d}$ with $d=2,3$, given by
\begin{equation}\label{eq:phi4p}
(\partial_{t}-\Delta+\mu)\varphi+\varphi^{3}=\xi,
\end{equation}
where $\xi$ is a space-time white noise on $\R_+\times\R^{d}$ and $\mu\in \R$.

Both equations fall in the category of the so-called \emph{singular} SPDEs, a loose term which means that they are classically ill-posed due to the very irregular nature of the noise $\xi$. Indeed, solutions are expected to take values only in spaces of distributions of negative regularity and the non-linear terms appearing in the equations cannot be given a canonical meaning.   Recent progresses by Hairer~\cite{hairer_theory_2014} and others~\cite{GIP,kupiainen_renormalization_2016,otto_quasilinear_2016} have provided various existence theories for local solutions of the above parabolic equations in a periodic spatial domain. The key idea is to identify suitable subspaces of distributions large enough to contain the candidate solutions and structured enough to allow for the definition of the non-linear terms.  These theories define solutions for the above equations once the non-linear term is \emph{renormalized}, which formally can be understood as a subtraction of an (infinite) correction term:
$$
\varphi^3 \mapsto \varphi^3 - \infty \varphi.
$$
More rigorously, and as we hinted above, this formal expression has to be understood in the sense that even though both terms separately are not well defined, certain combination has a well-defined meaning for a restricted class of distributions $\varphi$. The byproduct of the renormalization is that additional data (in the form of polynomials of the driving noise) have to be considered in order to identify canonically the result of the renormalization. 
It is not the main aim of this paper to discuss the features  of the local solution theory for singular SPDEs as this has been done extensively in the references cited above. 

Our aim  here is to develop a simple global solution theory for equations~\eqref{eq:phi4} and~\eqref{eq:phi4p}. Global solutions rely on specific properties of the equations, in particular here on the \emph{right} sign of the cubic non-linearity. The existence of  global in time solutions of the parabolic equation~\eqref{eq:phi4p} is relevant to the problem of \emph{stochastic quantization} of the $\Phi^4_d$ Euclidean  field theory, that is the measure $\nu$ on distributions over the $d$-dimensional periodic domain $\Lambda=\mathbb{T}^d$ formally given by the Euclidean path integral
\begin{equation}
\label{eq:phi4measure}
\nu(\mathd \phi)= \exp\left[-\int_{\Lambda} \left(\frac12|\nabla \phi|^2+\frac{\mu}{2}\phi^2+\frac14\phi^4\right) \right] \mathd \phi,
\end{equation}
where $\mathbb{T}=\R/2\pi \mathbb{Z}$.
Global in space solutions, that is solutions defined over all $\R^d$ correspond to the infinite volume limit of such a measure. Existence and uniqueness of global space-time solutions for the parabolic model in $d=2$ has been proved by Mourrat and Weber~\cite{MW17}. More recently the same authors have proven existence and uniqueness of global solutions in time on $\mathbb{T}^3$ in~\cite{MWcomedown}. In this last paper they also prove the stronger property, namely, that the solutions \emph{come down from infinity}, meaning that after a finite time the solution belongs to a compact set of the state space uniformly in the initial condition, a very strong property which is entirely due to the presence of the  cubic drift. These results show that singular SPDEs can be used to implement rigorously  the stochastic quantization approach first suggested by Parisi and Wu~\cite{parisi_perturbation_1981} and \emph{construct} random fields sampled according to the measure~\eqref{eq:phi4measure}. Another recent interesting approach which uses the SPDE to construct the measure $\nu$ is that of Albeverio and Kusuoka~\cite{albeverio_invariant_2017} which uses the invariance of approximations and uniform energy estimates on the SPDE to deduce tightness and existence of the limiting $\Phi^4_d$ measure~\eqref{eq:phi4measure}. 

In the present work we complete the picture by proving the global space-time existence and uniqueness for eq.~\eqref{eq:phi4p} in $\R^3$ with an associated \emph{coming down from infinity} property. This will be essentially a byproduct of the technique we develop to analyze the elliptic model~\eqref{eq:phi4} on $\R^d$ with $d=4,5$. The choice of dimensions has a two-fold origin: first it corresponds to the dimensions where the singularities of the elliptic equation match those of the parabolic one for $d=2,3$. Second (and partially related reason) is that there exists a very interesting conjecture of \emph{dimensional reduction} formulated first by Parisi and Sourlas~\cite{parisi_random_1979} which links the behavior of certain SPDEs in $d$ dimensions to that of Euclidean field theories in $d-2$ dimensions. In particular, it is conjectured that the trace on a codimension 2 hyperplane  of solutions to eq.~\eqref{eq:phi4} in $\R^d$ should have the law of the (parabolic) $\Phi^4_{d-2}$ model in $\R^{d-2}$, at least for $d=3,4,5$. This conjecture has been partially validated by rigorous arguments of Klein et al.~\cite{klein_supersymmetry_1983, klein_supersymmetry_1984} in the context of a regularized version of the models. Our study of the singular equation is another step to the full rigorous verification of the dimensional reduction phenomenon. The existence theory of the $d=3$ elliptic model is relatively straightforward and we will not consider it here.

Given the importance of these models in the mathematical physics literature and the open interesting conjectures they are related to, we found essential to devise streamlined arguments to treat global solutions of these equations. The main technical problem with globalization in the solution theory of singular SPDEs is given by the fact that the noise grows at infinity requiring the use of weighted spaces. This in turn requires to exploit fine properties of the equations in order to close the estimates. Witness of the important technical difficulties involved in the global analysis is the \emph{tour de force} that Mourrat and Weber~\cite{MWcomedown} had to put in place to solve the parabolic model on $\mathbb{T}^3$. One of the aim of the present paper is to provide also a simpler proof of their result, proof which is more in line with standard arguments of functional analysis/PDE theory. 
In order to do so we developed a new \emph{localization} technique which allows to split distributions belonging to weighted spaces into an irregular component which behaves nicely at the spatial infinity and a smooth component which grows in space. The localization technique allows to split singular SPDEs in two equations:
\begin{itemize}
\item[-]
 one containing the irregular terms but linear (or almost linear) and not requiring any particular care in the handling of the weighted spaces;
 \item[-] the other containing all the more regular terms and all the non-linearities which can be analyzed using standard PDE arguments, in particular pointwise maximum principle and  pointwise coercive estimates whose weighted version are easy to establish. This avoids the use of weighted $L^p$ spaces and related energy estimates which complicate the analysis of Mourrat and Weber~\cite{MWcomedown} and also of  Albeverio and Kusuoka~\cite{albeverio_invariant_2017}.
\end{itemize}
Other two improvements which we realize in this paper are the following:
\begin{itemize}
\item[a)] we use a direct $L^2$ energy estimate to establish uniqueness for the parabolic model, simplifying the proof and taking full advantage of our $L^\infty$ a priori estimates;
\item[b)] we use a time dependent weight to prove the coming down from infinity, going around the painful induction present in Mourrat and Weber paper and following quite closely the strategy one would adopt for classical driven reaction diffusion equations.
\end{itemize}
A problem which still remains open is that of the global uniqueness in the elliptic setting. Probably uniqueness does not hold or holds only for large masses. This is suggested by the behavior of the corresponding $\Phi^4_{d-2}$ model which is expected to undergo a phase transition at small temperature, corresponding here to a small mass.

\paragraph{Organization of the paper.} In Section \ref{s:prel} we introduce the basic notation and recall  various preliminary results concerning weighted Besov spaces. Then we present interpolation results  and construct the above mentioned localization operators, which are essential in the main body of the paper. As the next step, we establish Schauder and coercive estimates in weighted Besov spaces in both elliptic and parabolic setting and finally we discuss the basic results of the paracontrolled calculus.

In Section \ref{s:proba}, we recall the results of probabilistic analysis connected to the construction of the  stochastic objects needed in the sequel.

Sections \ref{sec:44}, \ref{sec:45} are devoted to the existence for the elliptic $\Phi^{4}$ model in dimension 4 and 5, respectively. More precisely, in the first step, we decompose the equations into  systems of two equations, one irregular  and the other one regular and containing the cubic nonlinearity. The next step is the cornerstone of our  analysis: we derive new a priori estimates for the unknowns of the decomposed system, which are then employed in order to establish existence of solutions. Here we first solve the equations on a large torus using a combination of a variational approach together with the Schaefer's fixed point theorem. Then we let the size of the torus converge to infinity and use compactness.

The a priori estimates from Sections \ref{sec:44}, \ref{sec:45} play the key role in the parabolic setting as well. Namely, in Sections \ref{sec:42}, \ref{sec:43} we study the parabolic $\Phi^{4}$ model in dimension 2 and 3, respectively. We follow a similar decomposition into a system of equations (only with a slight modification in dimension 3) and derive parabolic a priori estimates in analogy to the elliptic situation. These bounds are then used in the proof of existence. However, we proceed differently than in the elliptic setting: we work directly on the full space and mollify the noise, which leads to existence of smooth approximate solutions. The uniform estimates together with a compactness argument allow us to pass to the limit.



In Section \ref{s:uniq} we establish uniqueness of solutions in the parabolic setting. Unlike in the previous sections, it is not enough to work in the $L^{\infty}$-scale of weighted Besov spaces with polynomial weights. In particular, to compensate for the loss of weight in our estimates we employ exponential weights, requiring a different definition of the associated Besov spaces. This is discussed in Section \ref{ss:b}. The  proof of uniqueness then uses solely energy-type estimates in the $L^{2}$-scale of Besov spaces which  takes the full advantage of the well-chosen space-time weight.

Section \ref{s:d} is then concerned with the coming down from infinity property. Here we work with an additional weight in time which vanishes at zero and therefore allows to obtain bounds independent of the initial condition. Such a weight requires careful Schauder and coercive estimates that are established in Sections \ref{ss:d1}, \ref{ss:d2}. The proof of the coming down from infinity  then relies on our approach to a priori estimates from Section \ref{sec:44}, \ref{sec:45} together with a delicate control of the behavior at zero.

Finally, in Appendix \ref{s:a} we collect certain auxiliary results concerning existence for elliptic and parabolic variants of our problem in the smooth setting. Appendix \ref{s:sch} is then devoted to a refined Schauder estimate needed in Section \ref{s:d}.




We point out that for didactic reasons and in order not to blur our arguments, we chose to include in  Section \ref{s:prel}  only the results needed for the existence in Sections ~\ref{sec:44}, \ref{sec:45}, \ref{sec:42}, \ref{sec:43}. Further generalizations are  needed for uniqueness in Section \ref{s:uniq} and for the coming down from infinity in Section \ref{s:d}. The corresponding preliminaries are then discussed directly in the respective sections.


\paragraph{Acknowledgement.} MG is partially supported by the German Research Foundation (DFG) via CRC 1060.
 



\section{Preliminaries}
\label{s:prel}

\subsection{Weighted Besov spaces}
\label{ssec:besov}


As the first step, we introduce weighted Besov spaces which will be used in the sequel. Recall that the collection of  admissible weight functions is the collection of all positive $C^\infty(\mathbb{R}^d)$ functions $\rho$ with the following properties:
\begin{enumerate}
\item For all $\gamma\in\mathbb{N}_0^d$ there is a positive constant $c_\gamma$ with
$$
|D^\gamma\rho(x)|\leqslant c_\gamma \rho(x),\qquad\text{for all }x\in\mathbb{R}^d.
$$
\item There are two constants $c>0$ and $b\geq 0$ such that
$$
0<\rho(x)\leqslant c\rho(y)(1+|x-y|^2)^{b/2},\qquad\text{for all }x,y\in\mathbb{R}^d.
$$
\end{enumerate}



The space of Schwartz functions on $\R^{d}$ is denoted by $\mathcal{S}(\R^{d})$ and its dual, the space of tempered distributions is $\mathcal{S}'(\R^{d})$. The Fourier transform of $u\in \mathcal{S}'(\R^{d})$ is given by
$$
\mathcal{F}u(z)=\int_{\R^{d}}u(x)e^{-i z\cdot x}\,\dd x,
$$
so that the inverse Fourier transform is given by $\mathcal{F}^{-1}u(x)=(2\pi)^{-d}\mathcal{F}u(-x).$
 By $(\Delta_{i})_{i\geq -1}$ we denote the Littlewood--Paley blocks corresponding to a dyadic partition of unity.
If $\rho$ is an admissible weight and $\alpha\in\mathbb{R}$, we define the weighted Besov space $B^\alpha_{\infty,\infty}(\rho)\backassign \CC^\alpha(\rho)$ as the collection of all $f\in\mathcal{S}'(\R^{d})$ with finite norm
\[ \| f \|_{\CC^{\alpha} (\rho)} = \sup_{i\geq -1} 2^{i \alpha} \| \Delta_i f
   \|_{L^{\infty} (\rho)} = \sup_{i\geq -1} 2^{i \alpha} \| \rho \Delta_i f
   \|_{L^{\infty}} . \]
More details can be found  e.g. in \cite{T06}. 
Particularly,  due to \cite[Theorem 6.5]{T06}, it holds true that
\begin{equation}\label{eq:1}
   \|f\|_{\CC^\alpha(\rho)}\sim\|\rho f\|_{\CC^\alpha}
\end{equation}
   in the sense of equivalence of norms, where the latter denotes the norm in the classical (unweighted) Besov space $\CC^\alpha=B^\alpha_{\infty,\infty}(\R^{d})$. Moreover, it was shown in \cite[Theorem 6.9]{T06} that for $\alpha\in(0,M)$ with $M\in\N$, the weighted space $\CC^\alpha(\rho)$ admits an equivalent norm given by
\begin{equation}\label{eq:2h}
\|f\|_{L^\infty(\rho)}+\sup_{0<|h|\leqslant 1}|h|^{-\alpha}\|\Delta^M_h f\|_{L^\infty(\rho)},
\end{equation}
where $\Delta^M_h$ is the $M^{th}$-order finite difference operator  defined inductively by
$$
(\Delta^1_hf)(x)=f(x+h)-f(x),\qquad (\Delta^{\ell+1}_h)f(x)=\Delta^1_h(\Delta^{\ell}_hf)(x),\qquad\ell\in\N,
$$
Introduce a partition of unity  $\sum_{m\in\mathbb{Z}^{d}} \Lambda_m = 1$, where $\Lambda_{m}(x):=\Lambda(x-m)$ for a compactly supported $C^{\8}$-function $\Lambda$ on $\R^{d}$ and $m\in\mathbb{Z}^{d}$. Then the following localization principle for weighted Besov spaces follows from \eqref{eq:1} and \cite[Theorem 2.4.7]{T92}: let $\alpha\in\R$ then
\begin{equation}\label{eq:lem2.1}
   \|f\|_{\CC^\alpha(\rho)}\sim\sup_{m\in\mathbb{Z}^{d}}\|\Lambda_{m} f\|_{\CC^\alpha(\rho)}
\end{equation}
holds true in the sense of equivalence of norms. For most of our purposes, the following result in the case $\alpha>0$ will be sufficient. Let  $\sum_{k\geq-1} w_k = 1$ be a smooth partition of unity in spherical dyadic slices
where $w_{-1}$ is supported in a ball containing zero and there exists an annulus $\mathcal{A}=\{x\in\R^{d};a\leq|x|\leq b\}$ for some $0<a<b$ such that each $w_k$ for $k\geq0$ is supported in the annulus $2^k\mathcal{A}$. Set
$\tilde{w}_k = \sum_{\substack{i\geq-1\\i \sim k}} w_i$, where we write $i\sim k$ provided $\mathrm{supp}\,  w_{i}\cap \mathrm{supp}\, w_{k}\neq \emptyset.$


\begin{lemma}\label{lem:2.1}
It holds true that
\[ \| f \|_{L^{\infty} (\rho)} \leqslant \sup_{k\geq-1} \| \tilde{w}_k f
   \|_{L^{\infty} (\rho)}, \]
and if $\alpha>0$ then also
   $$
   \| f \|_{\CC^{\alpha} (\rho)} \lesssim\sup_{k\geq-1} \| \tilde{w}_k f
   \|_{\CC^{\alpha} (\rho)}.
   $$
\end{lemma}

\begin{proof}
Due to the construction of $(\tilde w_k)_{k\geq -1}$, for every $x\in\R^d$ there exists $k\geq -1$ such that $f(y)=\tilde{w}_k(y)f(y)$ for all $y\in \R^d$ with $|x-y|<1$. Consequently, the first claim follows. To show the second one, let $M\in\N$ be the smallest integer such that $\alpha<M$.
Then, it can be observed that, in addition to \eqref{eq:2h}, also 
$$
\| f \|_{L^{\infty} (\rho)} +\sup_{0<|h|< \frac{1}{M}}h^{-\alpha}\|\Delta^M_h f\|_{L^\infty(\rho)}
$$
defines an equivalent norm on $\CC^\alpha(\rho)$. The first summand is estimated as in the previous step. For the second summand, consider $|h|<\frac{1}{M}$. Since $(\Delta^M_hf)(x)$ depends only on  values  $f(y)$ for  $|y-x|\leqslant M|h|<1$, we deduce that for every $x\in \R^d$ there exists $k\in\N_0$ such that $f(y)=\tilde w_k(y)f(y)$ whenever $|y-x|<1$ and consequently also $(\Delta^M_hf)(y)=(\Delta^M_h(\tilde w_kf))(y)$. Thus
$$
\sup_{0<|h|< \frac{1}{M}}h^{-\alpha}\|\Delta^M_h f\|_{L^\infty(\rho)}\leqslant \sup_{k\geq-1}\sup_{0<|h|< \frac{1}{M}}h^{-\alpha}\|\Delta^M_h (\tilde w_kf)\|_{L^\infty(\rho)}
$$
and the second claim follows.
\end{proof}




Throughout this paper $\rho$ stands for a weight which is admissible and either constant or decreasing at infinity. It depends only on the space variable in the case of elliptic problems  or on both space and time for parabolic equations. We will not repeat the word ``admissible'' in the sequel. Moreover,  we will often work with  polynomial weights of the form $\rho(x)=\langle x\rangle^{-\nu}$ where  $\langle x\rangle=(1+|x|^{2})^{1/2}$ and $\nu\geq 0$. In the same spirit we will consider space-time dependent polynomial weights or $\rho(t,x)=\langle (t,x)\rangle^{-\nu}=(1+|(t,x)|^{2})^{-\nu/2}$ for $\nu\geq 0$. In addition, certain non-admissible weights will be needed in Section \ref{s:uniq} and Section \ref{s:d}. Namely, the proof of uniqueness in Section \ref{s:uniq} employs a weight that vanishes exponentially at infinity and consequently the definition of the associated Besov spaces cannot be based on Schwartz functions but rather on the so-called Gevrey classes as discussed in \cite{MW17}. The coming down from infinity property in Section~\ref{s:d} then requires a weight in time that vanishes in zero and is therefore also not an admissible weight in the sense of the above definition. The necessary results for these particular weights are discussed in Section~\ref{ss:b} and Sections \ref{ss:d0},  \ref{ss:d1}, \ref{ss:d2}.

Let $\rho$ be a polynomial space-dependent weight. Then the following embedding holds true
\begin{equation}\label{eq:emb}
\CC^{\beta_1}(\rho^{\gamma_1})\subset \CC^{\beta_2}(\rho^{\gamma_2}) \quad\text{provided}\quad \beta_1\geq\beta_2,\ \gamma_1\leqslant \gamma_2,
\end{equation}
and, according to \cite[Theorem 6.31]{T06}, the embedding in \eqref{eq:emb} is compact provided  $ \beta_1>\beta_2$ and $\gamma_1< \gamma_2$.



For parabolic equations, we will also need weighted   function spaces of space-time dependent functions/distributions. Let $\rho$ be a polynomial space-time weight and $\alpha\in\R$ and denote $\rho_{t}(\cdot)=\rho(t,\cdot)$, $t\in[0,\8)$.  Then $C\CC^{\alpha}(\rho)$ is the space of space-time distributions $f$ that are continuous in time, satisfy  $f(t)\in\CC^{\alpha}(\rho_{t})$ for every $t\in[0,\infty)$, and have finite norm
\[
\|f\|_{C\CC^{\alpha}(\rho)}:=\sup_{t\geq 0} \|(\rho f)(t)\|_{\CC^{\alpha}}.
\]
If a mapping $f:[0,\infty)\to  \CC^{\alpha}(\rho_{0}) $ is only bounded but not continuous, we write $f\in L^{\infty}\CC^{\alpha}(\rho)$ with the norm
\[
\|f\|_{L^{\infty}\CC^{\alpha}(\rho)}:=\esssup_{t\geq 0} \|(\rho f)(t)\|_{\CC^{\alpha}}<\8.
\]
Time regularity will be measured in terms of classical H\"older norms. In particular, for $\alpha\in(0,1)$ and $\beta\in\R$ we denote by $C^{\alpha}\CC^{\beta}(\rho)$  the space of mappings $f:[0,\infty)\to  \CC^{\beta}(\rho_{0}) $ with finite norm
$$
\|f\|_{C^{\alpha}\CC^{\beta}(\rho)}:=\sup_{t\geq 0} \|(\rho f)(t)\|_{\CC^{\beta}}+\sup_{s,t\geq 0,s\neq t} \frac{\|(\rho f)(t)-(\rho f)(s)\|_{\CC^{\beta}}}{|t-s|^{\alpha}}.
$$
It can be seen (cf. \eqref{eq:2h}) that since $\rho$ is a polynomial weight, this norm is equivalent to
\begin{equation}\label{eq:teq}
\|f\|_{C^{\alpha}\CC^{\beta}(\rho)}\sim\sup_{t\geq 0} \|(\rho f)(t)\|_{\CC^{\beta}}+\sup_{s,t\geq 0,s\neq t} \frac{\|\rho_{t}(f(t)- f(s))\|_{\CC^{\beta}}}{|t-s|^{\alpha}}.
\end{equation}
Similarly, we define the space $C^{\alpha}L^{\8}(\rho)$. 


In the case we consider only a finite time interval $[0,T]$, for some $T>0$, and a time-independent weight $\rho$, we write $f\in C_{T}\CC^{\alpha}(\rho)$, $f\in L^{\infty}_{T}\CC^{\alpha}(\rho)$, $C_{T}^{\alpha}\CC^{\beta}(\rho)$ and $C_{T}^{\alpha}L^{\8}(\rho)$ with straightforward modifications in the corresponding norms.




\subsection{Interpolation}

We present a simple interpolation result for  weighted Besov spaces.


\begin{lemma}
  \label{lemma:interp}
 Let $\kappa\in (0,1)$ and let $\rho$ be a space-time weight.  We have, for any $\alpha \in [0, 2 + \kappa]$
  \[ \| \psi \|_{\CC^{\alpha} (\rho^{1 + \alpha})} \lesssim \| \psi
     \|_{L^{\infty} (\rho)}^{1 - \alpha / (2 + \kappa)} \| \psi \|_{\CC^{2 +
     \kappa} (\rho^{3 + \kappa})}^{\alpha / (2 + \kappa)} .\]
\end{lemma}

\begin{proof}
 It holds
  \[ \|  \Delta_k \psi \|_{L^{\infty} (\rho^{1 +
     \alpha})} \lesssim \|  \rho^{1 + \alpha} \Delta_k \psi
     \|_{L^{\infty}} \lesssim \|  \rho \Delta_k \psi
     \|_{L^{\infty}}^{1 - \alpha / (2 + \kappa)} \|  \rho^{3 +
     \kappa} \Delta_k \psi \|_{L^{\infty}}^{\alpha / (2 + \kappa)} \]
  \[ \lesssim \| \psi \|_{L^{\infty} (\rho)}^{1 - \alpha /
     (2 + \kappa)} \| \Delta_k \psi \|_{L^{\infty}
     (\rho^{3 + \kappa})}^{\alpha / (2 + \kappa)} \lesssim 2^{- \alpha k} \|
\psi \|_{ L^{\infty} (\rho)}^{1 - \alpha / (2 + \kappa)} \|
 \psi \|_{\CC^{2 + \kappa} (\rho^{3 +
     \kappa})}^{\alpha / (2 + \kappa)} \]
  which proves the  claim.
\end{proof}


We will also need the following version  adapted to time-dependent problems.



\begin{lemma}
  \label{lemma:interp2}
Let $\kappa\in (0,1)$ and let $\rho$ be a space-time weight.  We have, for any $\alpha \in [0, 2 + \kappa]$
  \[ \| \psi \|_{C\CC^{\alpha} (\rho^{1 + \alpha})} \lesssim \| \psi
     \|_{L^{\8}L^{\infty} (\rho)}^{1 - \alpha / (2 + \kappa)} \| \psi \|_{C\CC^{2 +
     \kappa} (\rho^{3 + \kappa})}^{\alpha / (2 + \kappa)} .\]
Moreover, if  $\alpha/2\notin\N_{0}$ then
       \[ \| \psi \|_{C^{\alpha/2}L^{\8} (\rho^{1 + \alpha})} \lesssim \| \psi
     \|_{L^{\8}L^{\infty} (\rho)}^{1 - \alpha / (2 + \kappa)} \| \psi \|_{C^{(2+\kappa)/2}_{b}L^{\8} (\rho^{3 + \kappa})}^{\alpha / (2 + \kappa)} .\]
\end{lemma}

\begin{proof}
The first claim is a straightforward modification of Lemma \ref{lemma:interp}.
The second one can be obtained by the same approach since for $\alpha/2\notin\N_{0}$ the H\"older space $C^{\alpha/2}$ can be identified with the Besov space $B^{\alpha/2}_{\8,\8}$ and functions in $C^{\alpha/2}L^{\8}(\rho)$ can be naturally extended to be defined on the full space $\R\times\R^{d}$ while preserving the same norm.
\end{proof}





\subsection{Localization operators}
\label{ssec:local}



Here we construct localization operators $\UU_{>},
\UU_{\leqslant}$ which play the key role in our analysis. These localizers allow to decompose  a distribution $f$ into a sum of two components: one belongs to a (weighted) Besov space of higher regularity whereas the other one is less regular. To this end, let  $\sum_{k\geq-1} w_k = 1$ be a smooth dyadic  partition of unity on $\R^{d}$
where $w_{-1}$ is supported in a ball containing zero and there exists an annulus $\mathcal{A}=\{x\in\R^{d};a\leq|x|\leq b\}$ for some $0<a<b$ such that each $w_k$ for $k\geq0$ is supported in the annulus $2^k\mathcal{A}$.
%
Let $(L_{k})_{k\geq -1}\subset [-1,\infty)$ be a sequence of real numbers and let $f\in\mathcal{S}'(\R^{d})$. We define the localization operators by
\[ \UU_{>} f = \sum_k w_k \Delta_{> L_k} f, \qquad
   \UU_{\leqslant} f = \sum_k w_k \Delta_{\leqslant L_k} f, \]
where $\Delta_{>L_{k}}=\sum_{j:j>L_{k}}\Delta_{j}$ and $\Delta_{\leq L_{k}}=\sum_{j:j\leq L_{k}}\Delta_{j}$. We point out that in the sequel, we will use various localizing sequence $(L_{k})_{k\geq -1}$, depending on the context. However, for notational simplicity, we
will not denote these operators by different symbols.


\begin{lemma}\label{lem:local}
\rmbb{Let $L>0$ be given. There exists a  choice of parameters $(L_k)_{k\geq -1}$ such that for all $\alpha,\delta,\gamma>0$ and $a,b\in\R$ it holds true
      \[ \| \UU_{>} f \|_{\CC^{-\alpha - \delta} (\rho^{- a})} \lesssim 2^{- \delta L} \| f \|_{\CC^{-\alpha}
     (\rho^{-a+\delta})}, \qquad \| \UU_{\leqslant} f \|_{\CC^{-\alpha+\gamma}
     (\rho^{b})} \lesssim 2^{\gamma	
     L} \| f \|_{\CC^{-\alpha} (\rho^{b-\gamma})}, \]
          where the proportionality constant depends on $\alpha,\delta,\gamma,a,b$ but is independent of $f$.}
\end{lemma}

\begin{proof}
  Denote $c_k = - \log_2 \|  w_k \|_{L^{\infty} (\rho)}$ and let $\beta>\alpha+\delta$. Then we have
  \[ \|  w_k \|_{\CC^{\beta} (\rho^{\gamma})} \simeq \|  w_k \|_{L^{\infty}
     (\rho)}^{\gamma} = 2^{- \gamma c_k}, \]
  \[ \|  w_k \|_{\CC^{\beta} (\rho^{- \delta})} \simeq \|  w_k
     \|_{L^{\infty} (\rho)}^{- \delta} = 2^{\delta c_k} . \]
According to \eqref{eq:lem2.1}    and since there exists $M\in\N$ such that for every $m\in\mathbb{Z}^{d}$  the support of $\Lambda_{m}$ intersects the support of $w_{k}$ only for $k\in A_{m}$, where $A_{m}$ in a set of cardinality at most $M$, it holds 
  \[ \| \UU_{>} f \|_{\CC^{-\alpha - \delta} (\rho^{- a})} \lesssim
     \sup_{m\in\mathbb{Z}^{d}}\|\Lambda_{m}\UU_{>}f\|_{\CC^{ -\alpha - \delta} (\rho^{- a})}\lesssim \sup_{m\in\mathbb{Z}^{d}}\|\Lambda_{m}\sum_{k\in A_{m}}w_{k}\Delta_{>L_{k}}f\|_{\CC^{-\alpha - \delta} (\rho^{- a})}\]
 \[  \lesssim_{M} \sup_{k}\|  w_k \Delta_{> L_k} f \|_{\CC^{ -\alpha - \delta} (\rho^{- a})}
     \lesssim \sup_k \|  w_k \|_{\CC^{\beta} (\rho^{-\delta})}
     \| \Delta_{> L_k} f \|_{\CC^{ -\alpha - \delta} (\rho^{-a+\delta})} \]
  \[ \lesssim \sup_k 2^{\delta c_k - \delta L_k} \| f \|_{\CC^{- \alpha}
     (\rho^{-a+\delta})} \lesssim 2^{- \delta L} \| f \|_{\CC^{- \alpha}
     (\rho^{-a+\delta})}, \]
where we set $L_k =  c_k+ L$.
  On the other hand, the same argument implies
  \[ \| \UU_{\leqslant} f \|_{\CC^{-\alpha+\gamma} (\rho^b)} \lesssim \sup_k
     \|  w_k \Delta_{\leqslant L_k} f \|_{\CC^{-\alpha+\gamma} (\rho^b)} \lesssim
     \sup_k \|  w_k \|_{\CC^{\beta} (\rho^{\gamma})} \| \Delta_{\leqslant
     L_k} f \|_{\CC^{-\alpha+\gamma} (\rho^{b-\gamma})} \]
  \[ \lesssim \sup_k 2^{\gamma L_k - \gamma c_k} \| f
     \|_{\CC^{-\alpha} (\rho^{b-\gamma})} \lesssim 2^{\gamma L} \| f \|_{\CC^{- \alpha}
     (\rho^{b-\gamma})}.\]
\end{proof}

\begin{remark}
     Note that the sequence $(L_k)_{k\geq -1}$ in Lemma \ref{lem:local} does not depend on any of the parameters $\alpha,\delta,\kappa,a,b$ nor on the function $f$.

\end{remark}





We will also need the following version adapted to time-dependent problems. Let $(v_{\ell})_{\ell\geq -1}$ be a smooth dyadic partition of unity on $[0,\8)$ such that  $v_{-1}$ is supported in a ball containing zero and there exists an annulus $\mathcal{A}=\{t\in[0,\infty);a\leq t\leq b\}$ for some $0<a<b$ such that each $v_\ell$ for $\ell\geq0$ is supported in the annulus $2^k\mathcal{A}$. 
Let $\tilde v_{\ell}=\sum_{i:i\sim\ell}v_{i}$.
 For a given sequence $(L_{k,\ell})_{k,\ell\geq -1}$  we define localization operators $\VV_{>},
\VV_{\leqslant}$  by
\begin{equation}\label{eq:locd}
\VV_{>} f = \sum_{k,\ell} v_{\ell}w_k \Delta_{> L_{k,\ell}} f, \qquad
   \VV_{\leqslant} f = \sum_{k,\ell} v_{\ell}w_k \Delta_{\leqslant L_{k,\ell}} f.
   \end{equation}


\begin{lemma}\label{lem:local2}
Let $L>0$ be given and  let $\rho$ be a space-time weight. There exists a  choice of parameters $(L_{k,\ell})_{k,\ell\geq -1}$ such that for all $\alpha,\delta,\gamma>0$ and $a,b\in\R$ it holds true
      \[ \| \VV_{>} f \|_{C\CC^{- \alpha - \delta} (\rho^{- a})} \lesssim 2^{- \delta L} \| f \|_{C\CC^{- \alpha}
     (\rho^{-a+\delta})}, \qquad \| \VV_{\leqslant} f \|_{C\CC^{-\alpha+\gamma}
     (\rho^{b})} \lesssim 2^{\gamma
     L} \| f \|_{C\CC^{- \alpha} (\rho^{b-\gamma})}, \]
          where the proportionality constant depends on $\alpha,\delta,\kappa,a,b$ but is independent of $f$.
\end{lemma}

\begin{proof}
Similarly to the proof of Lemma \ref{lemma:interp2} we denote $c_{k,\ell} = - \log_2 \|\tilde v_{\ell}  w_k \|_{CL^{\infty} (\rho)}$ and let $\beta>\alpha+\delta$.
  Then we have
  \[ \|\tilde v_{\ell}  w_k \|_{C\CC^{\beta} (\rho^{\gamma})} \simeq \|\tilde v_{\ell}  w_k \|_{CL^{\infty}
     (\rho)}^{\gamma} = 2^{- \gamma c_{k,\ell}}, \]
  \[ \|\tilde v_{\ell}  w_k \|_{C\CC^{\beta} (\rho^{- \delta})} \simeq \|\tilde v_{\ell}  w_k
     \|_{CL^{\infty} (\rho)}^{- \delta} = 2^{\delta c_{k,\ell}} . \]
In view of  \eqref{eq:lem2.1} and Lemma \ref{lem:2.1}, we deduce (similarly to the proof of Lemma \ref{lem:local}) that
  \[ \| \VV_{>} f \|_{C\CC^{- \alpha - \delta} (\rho^{- a})} \lesssim
     \sup_{k,\ell} \|\tilde v_{\ell}  w_k \Delta_{> L_{k,\ell}} f \|_{C\CC^{- \alpha - \delta} (\rho^{- a})}
    \]
    \[ \lesssim \sup_{k,\ell} \|\tilde v_{\ell}  w_k \|_{C\CC^{\beta} (\rho^{-\delta})}
     \| \Delta_{> L_{k,\ell}} f \|_{C\CC^{- \alpha - \delta} (\rho^{-a+\delta})} \]
  \[ \lesssim \sup_{k,\ell} 2^{\delta c_{k,\ell} - \delta L_{k,\ell}} \| f \|_{C\CC^{- \alpha}
     (\rho^{-a+\delta})} \lesssim 2^{- \delta L} \| f \|_{C\CC^{- \alpha}
     (\rho^{-a+\delta})}, \]
where we set $L_{k,\ell} =  c_{k,\ell}+ L$.
  On the other hand, it holds
  \[ \| \VV_{\leqslant} f \|_{C\CC^{-\alpha+\gamma} (\rho^b)} \lesssim \sup_{k,\ell}
     \|\tilde v_{\ell}  w_k \Delta_{\leqslant L_{k,\ell}} f \|_{C\CC^{-\alpha+\gamma} (\rho^b)} \lesssim
     \sup_{k,\ell} \|\tilde v_{\ell}  w_k \|_{C\CC^{\beta} (\rho^{\gamma})} \| \Delta_{\leqslant
     L_{k,\ell}} f \|_{C\CC^{-\alpha+\gamma} (\rho^{b-\gamma})} \]
  \[ \lesssim \sup_{k,\ell} 2^{\gamma L_{k,\ell} - \gamma c_{k,\ell}} \| f
     \|_{C\CC^{- \alpha} (\rho^{b-\gamma})} \lesssim 2^{\gamma L} \| f \|_{C\CC^{- \alpha}
     (\rho^{b-\gamma})}.\]
\end{proof}











\subsection{Elliptic Schauder estimates}
\label{s:ell1}


We proceed with Schauder estimates valid for elliptic partial differential equations with cubic nonlinearities. Throughout the paper, we denote $\Q=-\Delta+\mu$.

\begin{lemma}
  \label{lemma:schauder-ellptic}Fix $\kappa > 0$ and let $\psi \in \CC^{2 +
  \kappa} (\rho^{3 + \kappa}) \cap L^{\infty} (\rho)$ be a classical solution
  to
  \[ \Q \psi + \psi^3 = \Psi \]
  then
  \[ \| \psi \|_{\CC^{2 + \kappa} (\rho^{3 + \kappa})} \lesssim_{\rho,\mu}  \|
     \Psi \|_{\CC^{\kappa} (\rho^{3 + \kappa})} + \| \psi \|_{L^{\infty}
     (\rho)}^{3 + \kappa}. \]
\end{lemma}



\begin{proof}
In view of \cite[Theorem 6.5]{T06} it holds
$$
\|\Q f\|_{\CC^\alpha(\rho)}\simeq_{\mu} \|f\|_{\CC^{2+\alpha}(\rho)}
$$
in the sense of equivalence of norms. Hence
$$
\|\psi\|_{\CC^{2+\kappa}(\rho^{3+\kappa})}\lesssim \|\Q\psi\|_{\CC^\kappa(\rho^{3+\kappa})}\leqslant \|\Psi\|_{\CC^\kappa(\rho^{3+\kappa})}+\|\psi^3\|_{\CC^\kappa(\rho^{3+\kappa})}
$$
and we estimate  using Lemma \ref{lemma:interp} and weighted Young inequality 
  \[ \| \psi^3 \|_{\CC^{\kappa} (\rho^{3 + \kappa})} \lesssim \| \psi
     \|_{L^{\infty} (\rho)}^2 \| \psi \|_{\CC^{\kappa} (\rho^{1 + \kappa})}\lesssim \| \psi
     \|_{L^{\infty} (\rho)}^2 \| \psi \|_{L^{\infty} (\rho)}^{1 - \kappa / (2
     + \kappa)} \| \psi \|_{\CC^{2 + \kappa} (\rho^{3 + \kappa})}^{\kappa / (2
     + \kappa)}
    \]
    \[
    \leq
    c\left( \|\psi\|_{L^{\infty} (\rho)}^2 \| \psi \|_{L^{\infty} (\rho)}^{1 - \kappa / (2
     + \kappa)}\right)^{(2 + \kappa) / 2} + \frac{1}{2}\left( \| \psi \|_{\CC^{2 +
     \kappa} (\rho^{3 + \kappa})}^{\kappa / (2 + \kappa)} \right)^{(2 +
     \kappa) / \kappa}
     \]
     \[
     \leqslant c\| \psi
     \|_{L^{\infty} (\rho)}^{3 + \kappa} +
     \frac{1}{2} \| \psi \|_{\CC^{2 + \kappa} (\rho^{3 + \kappa})} .
    \]
Thus we finally deduce the claim.
%\[ \| \psi \|_{\CC^{2 + \kappa} (\rho^{3 + \kappa})} \lesssim \| \Psi
%     \|_{\CC^{\kappa} (\rho^{3 + \kappa})} + \| \psi \|_{L^{\infty} (\rho)}^{3
%     + \kappa} . \]
\end{proof}





\subsection{Elliptic coercive estimates}
\label{s:ell2}

An essential result in our analysis is the following maximum principle in the weighted setting.

\begin{lemma}
  \label{lemma:apriori-elliptic}
  Fix $\kappa > 0$ and let $\psi \in \CC^{2 +
  \kappa} (\rho^{3 + \kappa}) \cap L^{\infty} (\rho)$ be a classical solution
  to
  \[ \Q \psi + \psi^3 = \Psi. \]
  Then the following a priori estimate holds
  \[ \| \psi \|_{L^{\infty} (\rho)} \lesssim_{\rho,\mu} 1 + \| \Psi \|_{L^{\infty}
     (\rho^3)}^{1 / 3}. \]
\end{lemma}



\begin{proof}
Let $\rho > 0$ be the weight from the statement of the Lemma and let $\bar{\psi} = \rho \psi$. Due to the assumption, $\bar{\psi}$ is bounded and locally belongs  to $\CC^{2+\kappa}$. 
  Assume for a moment that $\bar{\psi}$ has a global maximum and let $\hat{x}$
  be a global maximum point of $\bar{\psi}$. Then at $\hat{x}$ we have
  \[ 0 = \nabla \bar{\psi} = \rho \nabla \psi + \psi \nabla \rho, \qquad 0
     \leqslant - \Delta \bar{\psi} = - \rho \Delta \psi - (\Delta \rho) \psi -
     2 \nabla \rho \nabla \psi = - \rho \Delta \psi - \left[ (\Delta \rho) - 2
     \frac{| \nabla \rho |^2}{\rho} \right] \psi, \]
  so always at $\hat{x}$ we also have
  \begin{equation*}%\label{eq:2.16}
   \psi^3 + \mu \psi \leqslant \Psi - \left[ \frac{(\Delta \rho)}{\rho} - 2
     \frac{| \nabla \rho |^2}{\rho^2} \right] \psi 
  \end{equation*}
and  multiplying by $\rho^3$ leads to
  \[ (\bar{\psi})^3 \leqslant \rho^3 \Psi - \rho^2 (\mu + [(\Delta \rho) /
     \rho - 2 (\nabla \rho / \rho)^2]) \bar{\psi}. \]
If $\bar{\psi}(\hat{ x})\geq 0$     then
  $(\bar{\psi}^{})_+^3 \leqslant \| \rho^3 \Psi
  \|_{L^{\infty}} + C_{\rho,\mu} \| \rho^2 \bar{\psi} \|_{L^{\infty}}$. A similar
  reasoning at minima gives  $(\bar{\psi}^{})_-^3 \leqslant \| \rho^3 \Psi
  \|_{L^{\infty}} + C_{\rho,\mu} \| \rho^2 \bar{\psi} \|_{L^{\infty}}$, hence
  \[ \| \psi \|_{L^{\infty} (\rho)} \leqslant \| \Psi \|_{L^{\infty}
     (\rho^3)}^{1 / 3} + C_{\rho,\mu} \| \psi \|_{L^{\infty} (\rho)}^{1 / 3} .
  \]
Using weighted Young inequality we can absorb the second term of the r.h.s. into the l.h.s. and conclude that
  \[ \| \psi \|_{L^{\infty} (\rho)} \lesssim_{\rho,\mu} 1 + \| \Psi \|_{L^{\infty}
     (\rho^3)}^{1 / 3} . \]

Next, we consider the situation when $\psi\rho$ does not attain its global maximum. Since $\psi\rho$ is smooth and  bounded on $\R^{d}$ due to the assumption, it follows that  $\psi\rho^{1+\delta}$ vanishes at infinity for every $\delta\in (0,1)$. Consequently, it has a global maximum point and the previous part of the proof applies with $\rho$ replaced by $\rho^{1+\delta}$. The conclusion then follows by sending $\delta\to 0$ since the corresponding constant $c_{\rho^{1+\delta},\mu}$ is bounded uniformly in $\delta\in (0,1).$
%
%
\end{proof}



\subsection{Parabolic Schauder estimates}

As the next step, we derive a parabolic analog of Section \ref{s:ell1}.
To this end,  we first observe that the following Schauder estimates hold true in the weighted Besov spaces. They can be proved similarly to  \cite[Lemma A.9]{GIP}, see also \cite[Section 3.2]{MW17}.

\begin{remark}\label{rem:mu}
We note that the Schauder estimates below are formulated for a positive mass $\mu>0$. However, it can be observed that for the parabolic $\Phi^{4}$ model studied in Sections \ref{sec:42}, \ref{sec:43}, \ref{s:uniq}, \ref{s:d} this does not bring any loss of generality. Indeed, we may always add a linear term with positive mass to both sides of the equation and consider the original massive term as a right hand side. This is not true for the elliptic $\Phi^{4}$ model where the positivity of the mass seems to be essential. For notational simplicity we therefore adopt the convention that $\mu>0$ throughout the paper, that is,  for both elliptic and parabolic equations.
\end{remark}

Recall that we denoted $\Q=-\Delta+\mu$ and let $\LL=\partial_{t}+\Q$. This notation will be used throughout the paper.

\begin{lemma}\label{lemma:sch}
Let $\mu>0$, $\alpha\in\R$ and let $\rho$ be a  space-time weight. Let $v$ and $w$ solve, respectively,
$$
\LL v=f,\quad v(0)=0,\quad\qquad \LL  w=0,\quad w(0)=w_{0}.
$$ 
Then  it holds uniformly over $t\geq 0$
\begin{equation}\label{eq:1vw}
\|v(t)\|_{\CC^{2+\alpha}(\rho_{t})}\lesssim \|f\|_{L^{\infty}\CC^{\alpha}(\rho)},\qquad
\|w(t)\|_{\CC^{2+\alpha}(\rho_{t})}\lesssim \|w_{0}\|_{\CC^{2+\alpha}(\rho_{0})},
\end{equation}
if $0\leqslant 2+\alpha < 2$ then
$$
\|v\|_{C^{(2+\alpha)/{2}}L^{\infty}(\rho)}\lesssim \|f\|_{L^{\infty}\CC^{\alpha}(\rho)},\qquad\|v\|_{C^{1}L^{\infty}(\rho)}\lesssim \|f\|_{C\CC^{0}(\rho)},
$$

$$
 \|w\|_{C^{(2+\alpha)/2}L^{\infty}(\rho)}\lesssim \|w_{0}\|_{\CC^{2+\alpha}(\rho_{0})}.
$$
\end{lemma}

\begin{proof}
Denote $P_{t}=e^{t(\Delta-\mu)}$ be the semigroup of operators generated by $\Delta-\mu$ and recall that  $\mu>0$. Consider a time independent weight $\rho$ and  observe that similarly to \cite[Lemma A.7, Lemma A.8]{GIP} it holds true  uniformly over $t\geq 0$
\begin{align*}
\|P_{t}g\|_{\CC^{\alpha+\delta}(\rho)}\lesssim e^{-\mu t}t^{-\delta/2}\|g\|_{\CC^{\alpha}(\rho)},\qquad \|P_{t}g\|_{\CC^{\delta}(\rho)}\lesssim e^{-\mu t} t^{-\delta/2}\|g\|_{L^{\8}(\rho)}
\end{align*}
and if $0\leq\alpha \leqslant 2$
$$
\|(P_{t}-\mathrm{Id})g\|_{L^{\8}(\rho)}\leqslant |e^{-\mu t}-1|\|e^{t\Delta}g \|_{L^{\8}(\rho)}+e^{-\mu t}\|(e^{t\Delta}-\mathrm{Id})g \|_{L^{\8}(\rho)} \lesssim t^{\alpha/2}\|g\|_{\CC^{\alpha}(\rho)}.
$$
For a space-time weight, we obtain by the same  argument 
\begin{align}\label{eq:sch10}
\|P_{t}g\|_{\CC^{\alpha+\delta}(\rho_{t})}\lesssim e^{-\mu t}t^{-\delta/2}\|g\|_{\CC^{\alpha}(\rho_{t})},\qquad \|P_{t}g\|_{\CC^{\delta}(\rho_{t})}\lesssim e^{-\mu t} t^{-\delta/2}\|g\|_{L^{\8}(\rho_{t})},
\end{align}
$$
\|(P_{t}-\mathrm{Id})g\|_{L^{\8}(\rho_{t})} \lesssim t^{\alpha/2}\|g\|_{\CC^{\alpha}(\rho_{t})}.
$$
Then, if $2^{-2k}>t$  it follows from the fact that the weight is nonincreasing in time that
\begin{align*}
\left\|\int_{0}^{t}P_{t-s}\Delta_{k}f_{s}\mathrm{d}s\right\|_{L^{\infty}(\rho_{t})}&\lesssim t 2^{-k\alpha}\|f\|_{L^{\infty}\CC^{\alpha}(\rho)}\lesssim 2^{-k(2+\alpha)}\|f\|_{L^{\infty}\CC^{\alpha}(\rho)}.
\end{align*}
If $2^{-2k}\leqslant t$ then we split the integral into two parts
\begin{align*}
\left\|\int_{t-2^{-2k}}^{t}P_{t-s}\Delta_{k}f_{s}\mathrm{d}s\right\|_{L^{\infty}(\rho_{t})}&\lesssim 2^{-k\alpha}\int_{t-2^{-2k}}^{t}e^{-\mu(t-s)}\mathrm{d}s \left\|f\right\|_{\CC^{\alpha}(\rho_{t})}\lesssim 2^{-k(2+\alpha)}\|f\|_{L^{\infty}\CC^{\alpha}(\rho)}
\end{align*}
and
\begin{align*}
&\left\|\int_{0}^{t-2^{-2k}}P_{t-s}\Delta_{k}f_{s}\mathrm{d}s\right\|_{L^{\infty}(\rho_{t})}\lesssim e^{-\mu t} \int_{0}^{t-2^{-2k}}e^{\mu s} (t-s)^{-1-\varepsilon}\dd s \, 2^{-k(\alpha+2(1+\varepsilon))}\|f\|_{L^{\infty}\CC^{\alpha}(\rho)}\\
&\qquad\lesssim t^{-\varepsilon}e^{-\mu t} \int_{0}^{1-2^{-2k}/t}e^{\mu ts} (1-s)^{-1-\varepsilon}\dd s \, 2^{-k(\alpha+2(1+\varepsilon))}\|f\|_{L^{\infty}\CC^{\alpha}(\rho)}\\
&\qquad\lesssim 2^{-k(\alpha+2)}\|f\|_{L^{\infty}\CC^{\alpha}(\rho)}= 2^{-k(\alpha+2)}\|f\|_{L^{\infty}\CC^{\alpha}(\rho)}.
\end{align*}
Note that all the above inequalities are uniform over $t\geq 0$. Hence the first bound in \eqref{eq:1vw} follows. The second one is obtained as (recall that the weight is nonincreasing in time)
$$
\|w(t)\|_{\CC^{2+\alpha}(\rho_{t})}\lesssim e^{-\mu t}\|w_{0}\|_{\CC^{2+\alpha}(\rho_{t})} \lesssim \|w_{0}\|_{\CC^{2+\alpha}(\rho_{0})}.
$$

The time regularity of $w$ follows from
\begin{align*}
\|w(t)-w(s)\|_{L^{\infty}(\rho_{t})}&=\|P_{s}(P_{t-s}-\mathrm{Id})w_{0}\|_{L^{\infty}(\rho_{t})}\lesssim |t-s|^{(2+\alpha)/2}\|w_{0}\|_{\CC^{2+\alpha}(\rho_{0})}
\end{align*}
and due to
$$
v(t)-v(s)=(P_{t-s}-\mathrm{Id})v({s})+\int_{s}^{t} P_{t-r}f({r})\dd r,\quad s<t,
$$
we obtain
\begin{align*}
\|v(t)-v(s)\|_{L^{\infty}(\rho_{t})}&\lesssim |t-s|^{(2+\alpha)/2}\|v({s})\|_{\CC^{2+\alpha}(\rho_{t})}+ |t-s|^{(2+\alpha)/2}\|f \|_{L^{\infty}\CC^{\alpha}(\rho)}\\
&\lesssim |t-s|^{(2+\alpha)/2}\|f \|_{L^{\infty}\CC^{\alpha}(\rho)}.
\end{align*}
The proof is complete.
\end{proof}


Next, we derive a Schauder estimate for  parabolic equations including a cubic nonlinearity.

\begin{lemma}
  \label{lemma:schauder-par}
Let $\mu>0$ and let $\rho$ be a space-time weight.  Fix $\kappa > 0$ and let $\psi \in C\CC^{2 +
  \kappa} (\rho^{3 + \kappa}) \cap C^{1}L^{\infty} (\rho^{3+\kappa})\cap L^{\infty}L^{\infty}(\rho)$ be a classical solution
  to
  \[\LL  \psi + \psi^3 = \Psi ,\qquad \psi(0)=\psi_{0}.\]
Then
  \[ \| \psi \|_{C\CC^{2 + \kappa} (\rho^{3 + \kappa})} +\|\psi\|_{C^{1}L^{\infty}(\rho^{3+\kappa})} \lesssim \|\psi_{0}\|_{\CC^{2+\kappa}(\rho_{0})}+  \|
     \Psi \|_{C\CC^{\kappa} (\rho^{3 + \kappa})} + \| \psi \|_{L^{\infty}L^{\infty}
     (\rho)}^{3 + \kappa}. \]
\end{lemma}


\begin{proof}
Due to Lemma \ref{lemma:sch} it holds
$$
\|\psi\|_{C\CC^{2+\kappa}(\rho^{3+\kappa})}+\|\psi\|_{C^{1}L^{\infty}(\rho^{3+\kappa})}\lesssim\|\psi_{0}\|_{\CC^{2+\kappa}(\rho_{0})}+  \|\Psi\|_{C\CC^\kappa(\rho^{3+\kappa})}+\|\psi^3\|_{C\CC^\kappa(\rho^{3+\kappa})}
$$
 and we estimate  pointwise in time using Lemma \ref{lemma:interp2} and the weighted Young inequality
  \[ \| \psi^3 \|_{C\CC^{\kappa} (\rho^{3 + \kappa})} \lesssim \| \psi
     \|_{L^{\infty}L^{\infty} (\rho)}^2 \| \psi \|_{C\CC^{\kappa} (\rho^{1 + \kappa})}\lesssim \| \psi
     \|_{L^{\infty}L^{\infty} (\rho)}^2 \| \psi \|_{L^{\infty}L^{\infty} (\rho)}^{1 - \kappa / (2
     + \kappa)} \| \psi \|_{C\CC^{2 + \kappa} (\rho^{3 + \kappa})}^{\kappa / (2
     + \kappa)}
    \]
    $$
    \leq
    c\left( \|\psi\|_{L^{\infty}L^{\infty} (\rho)}^2 \| \psi \|_{L^{\infty}L^{\infty} (\rho)}^{1 - \kappa / (2
     + \kappa)}\right)^{(2 + \kappa) / 2} + \frac{1}{2}\left( \| \psi \|_{C\CC^{2 +
     \kappa} (\rho^{3 + \kappa})}^{\kappa / (2 + \kappa)} \right)^{(2 +
     \kappa) / \kappa}
    $$
    $$ \leqslant c\| \psi
     \|_{L^{\infty} L^{\infty}(\rho)}^{3 + \kappa} +
     \frac{1}{2} \| \psi \|_{C\CC^{2 + \kappa} (\rho^{3 + \kappa})} .
    $$
    Hence
    $$
     \| \psi^3 \|_{C\CC^{\kappa} (\rho^{3 + \kappa})} \leqslant c\| \psi
     \|_{L^{\infty}L^{\infty} (\rho)}^{3 + \kappa} +
     \frac{1}{2} \| \psi \|_{C\CC^{2 + \kappa} (\rho^{3 + \kappa})} 
    $$
and the claim follows.
\end{proof}




\subsection{Parabolic coercive estimates}

Similarly to Section \ref{s:ell2} we obtain the following maximum principle for parabolic equations.


\begin{lemma}
  \label{lemma:apriori-parabolic}
 Let $\mu\in\R$ and let $\rho$ be a  space-time weight. Fix $\kappa > 0$ and let $\psi \in C\CC^{2+\kappa}(\rho^{3+\kappa})\cap C^{1}L^{\infty}(\rho^{3+\kappa})\cap L^{\infty}L^{\infty}(\rho)$ be a classical
  solution to
  \[\LL  \psi + \psi^3 = \Psi,\qquad \psi(0)=\psi_{0} .\]
 Then the following a priori estimate holds
  \[ \| \psi  \|_{L^{\infty}L^{\infty} (\rho)} \lesssim 1+ \|\psi_{0}\|_{L^{\infty}(\rho_{0})}+ \|\Psi\|^{1/3}_{L^{\infty}L^{\infty}(\rho^{3})},
  \]
  where $\rho_{0}=\rho(0,\cdot)$.
\end{lemma}

\begin{proof}
Let $\bar\psi=\psi\rho$ and assume for the moment that $\bar\psi$ attains its (global) maximum $M=\bar\psi(t^{*},x^{*})$ at the point $(t^{*},x^{*})$. If $M\leq0$, then it is necessary to investigate the minimum point (or alternatively the maximum of $-\bar\psi$), which we discuss below. Let us therefore assume  that $M>0$. If $t_{*}=0$ then
$$
\bar\psi\leqslant \|\psi_{0}\|_{L^{\infty}(\rho_{0})}
$$
Assume that $t^{*}>0$.
Then
$$
\rho^{2}\partial_{t}\bar\psi+\rho^{2}(-\Delta+\mu)\bar\psi +\bar\psi^{3} = \rho^{3}\Psi +\rho\partial_{t}\rho\bar\psi - \rho^{2}(\Delta \rho) \psi - 2\rho^{2}\nabla\rho\nabla\psi.
$$
and
$$
\partial_{t}\bar\psi(t^{*},x^{*})= 0,\qquad\nabla\bar\psi(t^{*},x^{*})=0,\qquad\Delta\bar\psi(t^{*},x^{*})\leq0
$$
hence $\rho\nabla\psi=-\psi\nabla\rho$. Consequently $-\rho^{2}\Delta\bar\psi(t^{*},x^{*})\geq0$ and also $\rho\partial_{t}\rho\bar\psi(t^{*},x^{*})\leqslant 0$ since $\partial_{t}\rho
\leqslant 0$. Hence
$$
M^{3}\leqslant \left[\rho^{3}\Psi -\mu\rho^{2}\bar\psi - \rho^{2}(\Delta \rho) \psi - 2\rho^{2}\nabla\rho\nabla\psi\right]_{|_{(t^{*},x^{*})}}
$$
$$
\leqslant \|\Psi\|_{L^{\infty}L^{\infty}(\rho^{3})} +\rho^{2}(t^{*},x^{*})\left[|\mu|+\left\|\frac{\nabla\rho}{\rho}\right\|^{2}_{L^{\infty}}+\left\|\frac{\Delta\rho}{\rho}\right\|_{L^{\infty}}\right]\|\bar\psi\|_{L^{\infty}L^{\infty}}
$$
$$
\leqslant \|\Psi\|_{L^{\infty}L^{\infty}(\rho^{3})} +c_{\rho,\mu}\|\bar\psi\|_{L^{\infty}L^{\infty}}.
$$
Therefore we deduce that
$$
\bar\psi\lesssim \|\psi_{0}\|_{L^{\infty}(\rho_{0})}+\|\Psi\|^{1/3}_{L^{\infty}L^{\infty}(\rho^{3})} +c_{\rho,\mu}\|\bar\psi\|^{1/3}_{L^{\infty}L^{\infty}}.
$$
The same argument applied to $-\bar\psi$ yields
$$
-\bar\psi\lesssim \|\psi_{0}\|_{L^{\infty}(\rho_{0})}+\|\Psi\|^{1/3}_{L^{\infty}L^{\infty}(\rho^{3})} +c_{\rho,\mu}\|\bar\psi\|^{1/3}_{L^{\infty}L^{\infty}}
$$
hence, taking supremum over $(t,x)\in [0,\8)\times\R^{d}$ and applying the weighted Young inequality yields the claim.


Next, we consider the situation when $\psi\rho$ does not attain its global maximum. Since $\psi\rho$ is smooth and  bounded on $[0,\infty)\times\R^{d}$ due to the assumption, it follows that  $\psi\rho^{1+\delta}$ vanishes at infinity for every $\delta\in (0,1)$. Consequently, it has a global maximum point and the previous part of the proof applies with $\rho$ replaced by $\rho^{\delta}$. The conclusion then follows by sending $\delta\to 0$ since the corresponding constant $c_{\rho^{1+\delta},\mu}$ is bounded uniformly in $\delta\in (0,1).$
\end{proof}





\subsection{Paracontrolled calculus}
\label{ssec:para}


The foundations of paracontrolled calculus were laid down in the seminal work \cite{GIP} of Gubinelli, Imkeller and Perkowski, to which we  refer the reader for a number of facts used here. We refer to the book \cite{BCD} of Bahouri, Chemin and Danchin for a gentle introduction to the use of paradifferential calculus in the study of nonlinear PDEs. We shall then freely use the decomposition  $fg=f\prec g+f\circ g+f\succ g$, where $f\prec g=g\succ f$ and $f\circ g$, respectively, stands for the paraproduct of $f$ by $g$ and the corresponding resonant term, defined in terms of Littlewood--Paley decomposition.


The following basic results are obtained similarly to the unweighted setting.

\begin{lemma}\label{lem:bcd}
Let $\rho$ be an admissible weight.
\begin{enumerate}
\item Let $\mathcal{A}$ be an annulus, let $\alpha\in\R$ and let $(u_j)_{j\geq-1}$ be a sequence of smooth functions such that $\mathcal{F}u_j$ is supported in $2^j\mathcal{A}$ and  $\|u_j\|_{L^\infty(\rho)}\lesssim 2^{-j\alpha}$ for all $j\geq-1$. Then
$$
u=\sum_{j\geq-1}u_j\in \CC^\alpha(\rho)\qquad\text{and}\qquad\|u\|_{\CC^\alpha(\rho)}\lesssim \sup_{j\geq -1}\{2^{j\alpha}\|u_j\|_{L^\infty(\rho)}\}.
$$
\item Let $\mathcal{B}$ be a ball, let $\alpha>0$ and let $(u_j)_{j\geq-1}$ be a sequence of smooth functions such that $\mathcal{F}u_j$ is supported in $2^j\mathcal{B}$ and  $\|u_j\|_{L^\infty(\rho)}\lesssim 2^{-j\alpha}$ for all $j\geq-1$. Then
$$
u=\sum_{j\geq-1}u_j\in \CC^\alpha(\rho)\qquad\text{and}\qquad\|u\|_{\CC^\alpha(\rho)}\lesssim \sup_{j\geq -1}\{2^{j\alpha}\|u_j\|_{L^\infty(\rho)}\}.
$$
\end{enumerate}
\end{lemma}

\begin{proof}
The proof follows the lines of \cite[Lemma A.3]{GIP}.
\end{proof}





\begin{lemma}[Paraproduct estimates]\label{lem:para}
Let $\rho_{1},\rho_{2}$ be admissible weights  and $\beta\in\R$. Then  it holds
\begin{equation*}%\label{eq:para1}
\|f\prec g\|_{\CC^\beta(\rho_{1}\rho_{2})}\lesssim_\beta\|f\|_{L^\infty(\rho_{1})}\|g\|_{\CC^{\beta}(\rho_{2})},
\end{equation*}
and if $\alpha<0$ then
\begin{equation*}%\label{eq:para2}
\|f\prec g\|_{\CC^{\alpha+\beta}(\rho_{1}\rho_{2})}\lesssim_{\alpha,\beta}\|f\|_{\CC^{\alpha}(\rho_{1})}\|g\|_{\CC^{\beta}(\rho_{2})}.
\end{equation*}
If  $\alpha+\beta>0$ then it holds
\begin{equation*}%\label{eq:para3}
\|f\circ g\|_{\CC^{\alpha+\beta}(\rho_{1}\rho_{2})}\lesssim_{\alpha,\beta}\|f\|_{\CC^{\alpha}(\rho_{1})}\|g\|_{\CC^{\beta}(\rho_{2})}.
\end{equation*}
\end{lemma}

\begin{proof}
The proof follows the lines of \cite[Lemma 2.1]{GIP} and uses  Lemma \ref{lem:bcd} instead of \cite[Lemma A.3]{GIP}
\end{proof}

We also obtain the following weighted analog of \cite[Lemma 2.2, Lemma 2.3]{GIP}, which is proved analogously.


\begin{lemma}\label{lem:2.2}
Let $\rho_{1},\rho_{2}$ be admissible weights   and $\alpha\in(0,1)$ and $\beta\in\R$. For all $j\geq -1$ it holds
$$
\|\Delta_j(fg)-f\Delta_jg\|_{L^\infty(\rho_{1}\rho_{2})}\lesssim 2^{-j\alpha}\|f\|_{\CC^{\alpha}(\rho_{1})}\|g\|_{L^\infty(\rho_{2})}.
$$
$$
\|\Delta_j(f\prec g)-f\Delta_j g\|_{L^\infty(\rho_{1}\rho_{2})}\lesssim 2^{-j(\alpha+\beta)}\|f\|_{\CC^\alpha(\rho_{1})}\|g\|_{\CC^\beta(\rho_{2})}.
$$
\end{lemma}

With this in hand, we derive a weighted commutator estimate.

\begin{lemma}[Commutator lemma]\label{lem:com}
Let $\rho_{1}, \rho_{2}, \rho_{3}$ be admissible weights and let $\alpha\in (0,1)$ and $\beta,\gamma\in \R$ such that $\alpha+\beta+\gamma>0$ and $\beta+\gamma<0$. Then there exist a trilinear bounded operator $\mathrm{com}$ satisfying
$$
\|\mathrm{com}(f,g,h)\|_{\CC^{\alpha+\beta+\gamma}(\rho_{1}\rho_{2}\rho_{3})}\lesssim \|f\|_{\CC^\alpha(\rho_{1})}\|g\|_{\CC^\beta(\rho_{2})}\|h\|_{\CC^\gamma(\rho_{3})}
$$
and for smooth functions $f,g,h$
$$
\mathrm{com}(f,g,h)=(f\prec g)\circ h - f(g\circ h).
$$
\end{lemma}

\begin{proof}
In view of Lemma \ref{lem:bcd}, Lemma \ref{lem:para}, Lemma \ref{lem:2.2}, the proof is the same as the proof of \cite[Lemma 2.4]{GIP}.
\end{proof}


%We will also need the following paralinearization lemma for the square.
%
%\begin{lemma}[Paralinearization]\label{lem:paralin}
%Let $\rho$ be an admissible weight and $\alpha\in (0,1)$. Then
%$$
%\|f^2-2f\prec f\|_{\CC^{2\alpha}(\rho)}\lesssim \|f\|_{\CC^\alpha(\rho^\frac{1}{2})}^2.
%$$
%\end{lemma}
%
%\begin{proof}
%The proof is similar to \cite[Lemma~2.6]{GIP}.
%\end{proof}



Moreover, we will make use of the time-mollified paraproducts as introduced in \cite[Section 5]{GIP}. Let $Q:\R\to\R_{+}$ be a smooth function, supported in $[-1,1]$ and $\int_{\R}Q(s)\mathrm{d}s=1$, and for $i\geq -1$ define the operator $Q_{i}:C\CC^{\alpha}(\rho)\to C\CC^{\alpha}(\rho)$ by
$$
Q_{i}f(t)=\int_{\R}2^{2i}Q(2^{2i}(t-s))f(s\vee 0)\mathrm{d} s.
$$
Finally, we define the modified paraproduct of $f,g\in C\CC^{\alpha}(\rho)$ by
$$
f\Prec g := \sum_{i\geq -1}(S_{i-1}Q_{i}f)\Delta_{i} g.
$$
Setting  $\LL=\partial_{t} +(-\Delta+\mu)$, the following useful properties of this paraproduct in weighted Besov spaces  can be shown similarly to \cite[Lemma 5.1]{GIP}. Here we denote by $[\LL,f\Prec]$ the commutator between $\LL$ and $f\Prec$, that is, $[\LL,f\Prec]g=\LL (f\Prec g) -f\Prec(\LL g)$.

\begin{lemma}\label{lem:5.1}
Let $\rho_{1},\rho_{2}$ be admissible space-time  weights. Let $\alpha\in (0,1),$ $\beta\in \R$, and let $f\in C\CC^{\alpha}(\rho_{1})\cap C^{\alpha/2}L^{\infty}(\rho_{1})$ and $g\in C\CC^{\beta}(\rho_{2})$. Then
$$
\big\|[\LL,f\Prec] g\big\|_{C\CC^{\alpha+\beta-2}(\rho_{1}\rho_{2})}\lesssim \big( \|f\|_{C^{\alpha/2}L^{\infty}(\rho_{1})}+\|f\|_{C\CC^{\alpha}(\rho_{1})} \big)\|g\|_{C\CC^{\beta}(\rho_{2})},
$$
and
$$
\|f\prec g-f\Prec g\|_{C\CC^{\alpha+\beta}(\rho_{1}\rho_{2})}\lesssim \|f\|_{C^{\alpha/2}L^{\infty}(\rho_{1})}\|g\|_{C\CC^{\beta}(\rho_{2})}.
$$
\end{lemma}







\section{Probabilistic analysis}
\label{s:proba}

\subsection{Space white noise}
\label{ssec:renorm-el}

Let $\xi$ be a space white noise on $\R^{d}$, that is, a family of centered Gaussian random variables $\{\xi(h);\,h\in L^{2}(\R^{d}) \}$ such that
$$
\mathbb{E}[\xi(h)^{2}]=\|h\|^{2}_{L^{2}}.
$$
Let $\xi_{M}$ denote its periodization on $\mathbb{T}^{d}_{M}=(M\mathbb{T})^{d}=\left[-\frac{M}{2},\frac{M}{2}\right]^{d}$ given by
$$
\xi_{M}(h):=\xi(h_{M}),\qquad\text{where }\ h_{M}(x)=\mathbf{1}_{[-\frac{M}{2},\frac{M}{2}]^{d}}(x)\sum_{y\in M\mathbb{Z}^{d}}h(x+y).
$$
Let 
\begin{equation*}
 \Q X = \xi, \qquad  \Q X_{M} = \xi_{M},
\end{equation*}
and denote by $  \llbracket X^2 \rrbracket, \llbracket X^3 \rrbracket $ and $  \llbracket X_{M}^2 \rrbracket, \llbracket X_{M}^3 \rrbracket $ the corresponding Wick powers. They can be constructed by using a suitable mollification $\xi_{\varepsilon}=\xi*\eta_{\varepsilon}$ and $\xi_{M,\varepsilon}=\xi_{M}*\eta_{\varepsilon}$ (where $\eta_{\varepsilon}$ stands for a smoothing kernel) and setting
$$
 \Q X_{\varepsilon} = \xi_{\varepsilon}, \qquad  \Q X_{M,\varepsilon} = \xi_{M,\varepsilon},
$$
$$
 \llbracket X^2 \rrbracket  \assign \lim_{\varepsilon\to 0}  \llbracket X_{\varepsilon}^2 \rrbracket  \assign \lim_{\varepsilon\to 0}  X_{\varepsilon}^2 - a_{\varepsilon}, \qquad   \llbracket X^3 \rrbracket\assign\lim_{\varepsilon\to 0}  \llbracket X_{\varepsilon}^3 \rrbracket\assign \lim_{\varepsilon\to 0} X_{\varepsilon}^3 - 3 a_{\varepsilon} X_{\varepsilon},
 $$
 $$
 \llbracket X_{M}^2 \rrbracket  \assign \lim_{\varepsilon\to 0}  \llbracket X_{M,\varepsilon}^2 \rrbracket  \assign \lim_{\varepsilon\to 0}  X_{M,\varepsilon}^2 - a_{M,\varepsilon}, \qquad   \llbracket X_{M}^3 \rrbracket\assign\lim_{\varepsilon\to 0} \llbracket X_{M,\varepsilon}^3 \rrbracket\assign\lim_{\varepsilon\to 0} X_{M,\varepsilon}^3 - 3 a_{M,\varepsilon} X_{M,\varepsilon},
 $$
 where $a_{\varepsilon}= \mathbbm{E}[X_{\varepsilon}^2(0)]$  and $a_{M,\varepsilon}= \mathbbm{E}[X_{M,\varepsilon}^2(0)]$ are  constants diverging  as $\varepsilon\to0$ and the limits are understood in a suitable Besov space a.s. More precisely, the following result holds.
 
 
 \begin{theorem}\label{thm:renorm}
Let $d=4$. Let $\rho(x)=\langle x\rangle^{-\nu}$ for some $\nu>0$. Then there exist random distributions $X,\llbracket X^{2}\rrbracket,\llbracket X^{3}\rrbracket$ and $X_{M},\llbracket X_{M}^{2}\rrbracket,\llbracket X_{M}^{3}\rrbracket$ given by the formulas above, such that for every $\kappa,\sigma>0$  it holds 
$$
\|X\|_{\CC^{-\kappa}(\rho^{\sigma})},\|\llbracket X^{2}\rrbracket\|_{\CC^{-\kappa}(\rho^{\sigma})}
,\|\llbracket X^{3}\rrbracket\|_{\CC^{-\kappa}(\rho^{\sigma})}\lesssim 1,
$$
$$
\|X_{M}\|_{\CC^{-\kappa}(\mathbb{T}^{4}_{M})},\|\llbracket X_{M}^{2}\rrbracket\|_{\CC^{-\kappa}(\mathbb{T}^{4}_{M})}
,\|\llbracket X_{M}^{3}\rrbracket\|_{\CC^{-\kappa}(\mathbb{T}^{4}_{M})}\lesssim 1,
$$
and in addition $X_{M}\to X$, $\llbracket X_{M}^{2}\rrbracket\to \llbracket X^{2}\rrbracket$, $\llbracket X_{M}^{3}\rrbracket\to \llbracket X^{3}\rrbracket$ in $\CC^{-\kappa}(\rho^{\sigma})$ a.s. as $M\to\infty$. \end{theorem}
 
 \begin{proof} We give a sketch of the proof since similar arguments are already present in the literature on parabolic $\Phi^4_d$ models, and in particular in the work of Mourrat and Weber~\cite{MW17}. Following the approach of Gubinelli and Perkowski~\cite{gubinelli_kpz_2017} we represent the random fields $X$ and $X_M$ as Wiener integrals over a white noise $W$ on $\R^4$. As a consequence we can write 
 \[ X (x) = \int_{\mathbbm{R}^4} e^{2 \pi i \theta \cdot x} \frac{W (\mathd
   \theta)}{\mu + | \theta |^2}, \qquad X_M (x) = \int_{\mathbbm{R}^4} e^{2
   \pi i [\theta]_M \cdot x} \frac{W (\mathd \theta)}{\mu + | [\theta]_M |^2},
\]
where $([\theta]_M)^i = M^{- 1} \lfloor M \theta^i - 1 / 2 \rfloor$, $i = 1,
\ldots, d$ is the discretization of $\theta\in\R^4$ on a grid of size $M^{-1}$. The reader can check that this gives a periodic random field with the correct covariance. Wick powers of $X$ (or $X_M$) can then be expressed as multiple Wiener integrals over $W$. We present the details for $\llbracket X^3 \rrbracket $:
\[ \llbracket X^3 \rrbracket (x) = \int_{(\mathbbm{R}^4)^3} e^{2 \pi i
   (\theta_1 + \theta_2 + \theta_3) \cdot x} \frac{W (\mathd \theta_1 \mathd
   \theta_2 \mathd \theta_3)}{\prod_{i = 1}^3 (\mu + | \theta_i |^2)}, \]
\[ \llbracket X^3_M \rrbracket (x) = \int_{(\mathbbm{R}^4)^3} e^{2 \pi i
   ([\theta_1]_M + [\theta_2]_M + [\theta_3]_M) \cdot x} \frac{W (\mathd
   \theta_1 \mathd \theta_2 \mathd \theta_3)}{\prod_{i = 1}^3 (\mu + |
   [\theta_i]_M |^2)} . \]
And $L^2$ bound on the Littlewood--Paley block of these quantities reads, for $k\geq -1$,
\[ \mathbbm{E} [| \Delta_k \llbracket X^3_M \rrbracket (x) |^2] =
   \int_{(\mathbbm{R}^4)^3} K_k ([\theta_1]_M + [\theta_2]_M +
   [\theta_3]_M)^2 \frac{\mathd \theta_1 \mathd \theta_2 \mathd
   \theta_3}{\prod_{i = 1}^3 (\mu + | [\theta_i]_M |^2)^2} \]
\[ \lesssim \int_{(\mathbbm{R}^4)^3} K_k (\theta_1 + \theta_2 + \theta_3)^2
   \frac{\mathd \theta_1 \mathd \theta_2 \mathd \theta_3}{\prod_{i = 1}^3 (\mu
   + | \theta_i |^2)^2} \lesssim 1, \]
where $K_{k}$ is the Fourier multiplier associated with $\Delta_k$. 
From this we deduce by hypercontractivity that $\mathbbm{E} [| \Delta_k
\llbracket X^3_M \rrbracket (x) |^p] \lesssim 1$ and therefore that
\[ \mathbbm{E} [\| \Delta_k \llbracket X^3_M \rrbracket \|_{L^p
   (\rho^{\sigma})}^p] \lesssim \int_{\mathbbm{R}^4} \mathbbm{E} [| \Delta_k
   \llbracket X^3_M \rrbracket (x) |^p] \rho (x)^{\sigma p} \mathd x \lesssim
   \int_{\mathbbm{R}^4} \rho (x)^{\sigma p} \mathd x \lesssim 1, \]
for $p$ sufficiently large so that the space integral is finite. As a
consequence of Bernstein inequality it follows that
\[ \| \Delta_k \llbracket X^3_M \rrbracket \|_{L^{\infty} (\rho^{\sigma})}
   \lesssim 2^{4 k / p} \| \Delta_k \llbracket X^3_M \rrbracket \|_{L^p
   (\rho^{\sigma})} \]
and therefore 
\[ \mathbbm{E} (\|  \llbracket X^3_M \rrbracket \|_{C^{- \kappa}
   (\rho^{\sigma})}^p) < \infty, \]
for $p$ large enough and $\kappa > 0$ small. Convergence of $\llbracket X^3_M
\rrbracket$ to $\llbracket X^3 \rrbracket$ can be handled by coupling,
observing that estimation of $\llbracket X^3_M \rrbracket - \llbracket X^3
\rrbracket$ involves computations similar to the above. Indeed, it holds
\[ \mathbbm{E} [| \Delta_k (\llbracket X^3_M \rrbracket - \llbracket X^3
   \rrbracket) (x) |^2] = \int_{(\mathbbm{R}^4)^3} \left( \frac{K_k
   ([\theta_1]_M + [\theta_2]_M + [\theta_3]_M)}{\prod_{i = 1}^3 (\mu + |
   [\theta_i]_M |^2)} - \frac{K_k (\theta_1 + \theta_2 +
   \theta_3)}{\prod_{i = 1}^3 (\mu + | \theta_i |^2)} \right)^2 \mathd
   \theta_1 \mathd \theta_2 \mathd \theta_3 \]
which by dominated convergence tends to zero as $M \rightarrow \infty$.
Therefore we can estimate
\[ \mathbbm{E} [\| \Delta_k (\llbracket X^3_M \rrbracket - \llbracket X^3
   \rrbracket) \|_{L^p (\rho^{\sigma})}^p] \lesssim \int_{\mathbbm{R}^4}
   \mathbbm{E} [| \Delta_k (\llbracket X^3_M \rrbracket - \llbracket X^3
   \rrbracket) (x) |^p] \rho (x)^{\sigma p} \mathd x \lesssim o_M (1) . \]
 \end{proof}
 
 This result will be used for the study of elliptic $\Phi^{4}$ model in dimension 4, see Section \ref{sec:44}. When $d=5$ then the space white noise becomes more irregular and our analysis requires additional probabilistic objects. More precisely, we let
\[ \Q \tthreeone{X} = \llbracket X^3 \rrbracket, \quad \Q
   \ttwoone{X} = \llbracket X^2 \rrbracket, \]
   \[ \Q \tthreeone{X}_{\varepsilon} = \llbracket X_{\varepsilon}^3 \rrbracket, \quad \Q
   \ttwoone{X}_{\varepsilon} = \llbracket X_{\varepsilon}^2 \rrbracket, \]
\[\tthreetwor{X}=\lim_{\varepsilon\to 0}\tthreeone{X}_{\varepsilon}\circ X_{\varepsilon},\qquad \ttwothreer{X} = \lim_{\varepsilon\to 0}\ttwoone{X}_{\varepsilon} \circ \llbracket X_{\varepsilon}^2 \rrbracket -\frac{b_{\varepsilon}}{3}, \qquad
   \tthreethreer{X} = \lim_{\varepsilon\to 0}\tthreeone{X}_{\varepsilon} \circ \llbracket X_{\varepsilon}^2 \rrbracket -  b_{\varepsilon}X_{\varepsilon}, \]
   where
   $$
   b_{\varepsilon} := 3 \mathbbm{E} [(\ttwoone{X}_{\varepsilon} \circ \llbracket X_{\varepsilon}^2 \rrbracket)(0)].
   $$
 Similarly, we define the periodic analogs.
 
  \renewcommand{\arraystretch}{1.5}
   \begin{table}
      \begin{center}
  \begin{tabular}{| r | c | c | c | c|c|c|c|c|}
  \hline
  $\tau$ & $X$ & $ \llbracket X^2 \rrbracket$ & $ \llbracket X^3 \rrbracket$ & $\tthreeone{X}$ & $\ttwoone{X}$ & $\tthreetwor{X}$  & $\ttwothreer{X}$ & $ \tthreethreer{X}$\\ \hline 
  $\alpha_{\tau}$ & $-\frac{1}{2}-\kappa$ & $-1-\kappa$ & $-\frac32-\kappa $& $\frac{1}{2}-\kappa$ & $1-\kappa$ & $-\kappa$ & $-\kappa$ & $-\frac{1}{2}-\kappa$ \\
  \hline
  \end{tabular}
  \caption{Space regularity of stochastic objects in the $d=5$ elliptic or $d=3$ parabolic case.}
  \label{t:reg}
  \end{center}
  \end{table}
 
 \begin{theorem}\label{thm:renorm2}
Let $d=5$. Let $\rho(x)=\langle x\rangle^{-\nu}$ for some $\nu>0$. Then there exist random distributions
\begin{equation}\label{eq:r45}
X,\llbracket X^{2}\rrbracket,\llbracket X^{3}\rrbracket, \tthreeone{X}, \ttwoone{X}, \tthreetwor{X}, \ttwothreer{X}, \tthreethreer{X}
\end{equation}
and their periodic versions
\begin{equation}\label{eq:r45M}
X_{M},\llbracket X_{M}^{2}\rrbracket,\llbracket X_{M}^{3}\rrbracket, \tthreeone{X}_{M}, \ttwoone{X}_{M}, \tthreetwor{X}_{M}, \ttwothreer{X}_{M}, \tthreethreer{X}_{M}
\end{equation}
given by the formulas above, such that if $\tau$ denotes one of the distributions in \eqref{eq:r45} and $\tau_{M}$ is the associated periodic version from \eqref{eq:r45M}, then $\tau\in \CC^{\alpha_{\tau}}(\rho^{\sigma})$ and $\tau_{M}\in \CC^{\alpha_{\tau}}(\mathbb{T}^{5}_{M})$ for $\alpha_{\tau}$  given by Table \ref{t:reg} and every $\kappa,\sigma>0$. Moreover,
$\tau_{M }\to \tau$ in $\CC^{\alpha_{\tau}}(\rho^{\sigma})$ a.s. as $M\to\infty$.
 \end{theorem}
 
 \begin{proof}
 Apart form the higher complexity of the terms involved in the $d=5$ case, the analysis proceeds like in Theorem~\ref{thm:renorm}. The various stochastic objects can be written as multiple iterated Wiener integrals and renormalizations accounts for cancellations of certain terms in the associated kernels. In the periodic and parabolic setting this analysis has already been performed several times with small variations, for example in~\cite{CC},~\cite{mourrat_construction_2016} and more recently in~\cite{furlan_weak_2017} and in~\cite{gubinelli_kpz_2017} for the KPZ equation. Estimation in weighted Besov spaces and convergence of the periodic to the non-periodic versions proceed like in the $\R^4$ case. 
 \end{proof}
 

\subsection{Space-time white noise}
\label{ssec:renorm-par}



If $\xi$ is a space white noise on $\R\times\R^{d}$, i.e. a family of centered Gaussian random variables $\{\xi(h);\,h\in L^{2}(\R\times\R^{d}) \}$ such that
$$
\mathbb{E}[\xi(h)^{2}]=\|h\|^{2}_{L^{2}},
$$
then we may define its periodization $\xi_{M}$ on $\mathbb{T}^{d}_{M}=\left[-\frac{M}{2},\frac{M}{2}\right]^{d}$  by
$$
\xi_{M}(h):=\xi(h_{M}),\qquad\text{where }\ h_{M}(t,x)=\mathbf{1}_{[-\frac{M}{2},\frac{M}{2}]^{d}}(x)\sum_{y\in M\mathbb{Z}^{d}}h(t,x+y).
$$
Our construction of solutions to the parabolic $\Phi^{4}$ model in Section \ref{sec:42} and Section \ref{sec:43} relies on  a smooth and space periodic approximation $\xi_{\varepsilon}$ of the driving space-time white noise $\xi$, defined on the torus of size $M=\frac 1\varepsilon$. To be more precise, let $\xi_{\varepsilon}$ be a periodic version of a space-time mollification of $\xi$ defined on $\R\times\mathbb{T}^{2}_{1/\varepsilon}$ and let $X,X_{\varepsilon}$ be stationary solutions to
\begin{equation*}
 \LL X = \xi, \qquad  \LL X_{\varepsilon} = \xi_{\varepsilon},
\end{equation*}
 and 
$$
 \llbracket X^2 \rrbracket  \assign \lim_{\varepsilon\to 0} X_{\varepsilon}^2 - a_{\varepsilon}, \qquad   \llbracket X^3 \rrbracket\assign\lim_{\varepsilon\to 0} X_{\varepsilon}^3 - 3 a_{\varepsilon} X_{\varepsilon},
 $$
 where again we can take $a_{\varepsilon}= \mathbbm{E}[X_{\varepsilon}^2(0,0)]$ is a diverging constant and the limits are understood in a suitable Besov space a.s. More precisely, the following result holds.
 
 
 \begin{theorem}\label{thm:renorm42}
Let $d=2$. Let $\rho(t,x)=\langle (t,x)\rangle^{-\nu}$ for some $\nu>0$. There exists a sequence of diverging constants $(a_{\varepsilon})_{\varepsilon\in (0,1)}$ and random distributions $X,\llbracket X^{2}\rrbracket,\llbracket X^{3}\rrbracket$  such that for all $\kappa,\sigma>0$  it holds
$$
\|X\|_{C\CC^{-\kappa}(\rho^{\sigma})},\|\llbracket X^{2}\rrbracket\|_{C\CC^{-\kappa}(\rho^{\sigma})}
,\|\llbracket X^{3}\rrbracket\|_{C\CC^{-\kappa}(\rho^{\sigma})}\lesssim 1,
$$
and
$$
 \llbracket X^2 \rrbracket  = \lim_{\varepsilon\to 0} X_{\varepsilon}^2 - a_{\varepsilon}, \qquad   \llbracket X^3 \rrbracket=\lim_{\varepsilon\to 0} X_{\varepsilon}^3 - 3 a_{\varepsilon} X_{\varepsilon},
 $$
where the limit is understood in $C\CC^{-\kappa}(\rho^{\sigma})$ a.s. as $\varepsilon\to\infty$.
\end{theorem}

 \begin{proof} The proof proceeds like in Theorem~\ref{thm:renorm} in the proof of which we also made reference to the relevant literature. We would just like to comment on how to obtain existence for all times within the claimed weighted space. Let $Y$ be one of the random fields considered in the theorem and $Y_\varepsilon$ the corresponding approximation. By standard estimates one obtains bounds of the form
 $$
 \mathbbm{E}[\|\rho^\sigma(t,\cdot) Y(t,\cdot)-\rho^\sigma(s,\cdot) Y(s,\cdot) \|_{\CC^{-\kappa}}^p]\lesssim |t-s|^{\delta p} \langle s \rangle^{-\beta p},
 $$ 
 for some small $\delta,\beta>0$ and large $p$, uniformly for $0\le s \le t \le s +1$.
 Standard Kolomogorov criterion can be applied to obtain that 
 $$
 \mathbbm{E}[\|\rho^\sigma Y \|_{C^{\delta'}([L,L+1];\CC^{-\kappa})}^p]\lesssim L^{-\beta p},
 $$ 
 for all $L\ge 1$ and some $\delta'>0$. Finally if $p$ is large enough this shows that the random variable
 $
 \sum_{L=0}^\infty \|\rho^\sigma Y \|_{C^{\delta'}([L,L+1];\CC^{-\kappa})}^p
 $ 
 has finite expectation. Finally a simple gluing argument implies that the random variable $\|\rho^\sigma Y \|_{C^{\delta'}(\R_+;\CC^{-\kappa})}$ has also finite $L^p$ moments. Weighted space convergence of the approximation $Y_\varepsilon$ to $Y$ can be handled similarly since we can establish that
 $$
\sup_L L^{\beta p}\mathbbm{E}[\|\rho^\sigma Y_\varepsilon - \rho^\sigma Y \|_{C^{\delta'}([L,L+1];\CC^{-\kappa})}^p]\lesssim  o_\varepsilon(1),
 $$ 
from which we obtain easily the convergence in the weighted norm $C^{\delta'}\CC^{-\kappa}(\rho^\sigma)$ as $\varepsilon \to 0$. 
 \end{proof}
 

Similarly to the elliptic 5 dimensional case, we define
\[ \LL \tthreeone{X} = \llbracket X^3 \rrbracket,\quad\tthreeone{X}(0)=0,   \qquad \LL
   \ttwoone{X} = \llbracket X^2 \rrbracket, \quad \ttwoone{X}(0)=0,\]
      \[ \LL \tthreeone{X}_{\varepsilon} = \llbracket X_{\varepsilon}^3 \rrbracket,\quad \tthreeone{X}_{\varepsilon}(0)=0, \qquad \LL
   \ttwoone{X}_{\varepsilon} = \llbracket X_{\varepsilon}^2 \rrbracket,\quad \ttwoone{X}_{\varepsilon}(0)=0, \]
\[\tthreetwor{X}=\lim_{\varepsilon\to 0}\tthreeone{X}_{\varepsilon}\circ X_{\varepsilon},\qquad \ttwothreer{X} = \lim_{\varepsilon\to 0}\ttwoone{X}_{\varepsilon} \circ \llbracket X_{\varepsilon}^2 \rrbracket -\frac{b_{\varepsilon}}{3}, \qquad
   \tthreethreer{X} = \lim_{\varepsilon\to 0}\tthreeone{X}_{\varepsilon} \circ \llbracket X_{\varepsilon}^2 \rrbracket -  b_{\varepsilon}X_{\varepsilon}, \]
   where $b_{\varepsilon}(t) = 3 \mathbbm{E}[(\ttwoone{X}_{\varepsilon} \circ \llbracket X_{\varepsilon}^2 \rrbracket )(t,0)]$ stands for a suitable renormalization constant which is $t$ dependent and such that $\sup_{t\ge 0} |b_\varepsilon(t)| \lesssim |\log \varepsilon|$. Moreover, it can be seen that, for each fixed $\varepsilon$, $b_{\varepsilon}$ is smooth and has bounded first derivative.
   

   
   
 \begin{theorem}\label{thm:renorm43}
Let $d=3$. Let $\rho(t,x)=\langle (t,x)\rangle^{-\nu}$ for some $\nu>0$. Then there exist random distributions
\begin{equation}\label{eq:r43}
X,\llbracket X^{2}\rrbracket, \tthreeone{X}, \ttwoone{X}, \tthreetwor{X}, \ttwothreer{X}, \tthreethreer{X}
\end{equation}
such that if $\tau$ denotes one of the distributions in \eqref{eq:r43} then $$\tau\in C\CC^{\alpha_{\tau}}(\rho^{\sigma})\cap C^{\delta/2}\CC^{\alpha_{\tau}-\gamma}(\rho^{\sigma})  $$ for $\alpha_{\tau}$  given by Table \ref{t:reg}, every $\kappa,\sigma>0$ and some $\delta,\gamma>0$. Moreover, if $\tau_{\varepsilon}$ is the smooth version of  $\tau$ then
$\tau_{\varepsilon }\to \tau$ in $ C\CC^{\alpha_{\tau}}(\rho^{\sigma})\cap C^{\delta/2}\CC^{\alpha_{\tau}-\gamma}(\rho^{\sigma})$ a.s. as $\varepsilon\to0$.
 \end{theorem}

 
 \begin{proof} The convergence and renormalization of the stochastic terms has been performed several times in the literature, see the proof of Theorem~\ref{thm:renorm2} for precise references. As for the convergence in the space-time weighted Besov--H\"older spaces arguments similar to those described in Theorem~\ref{thm:renorm42} can be applied to establish the claim.
 \end{proof}
 
\rmbb{\begin{remark}
We note that $\llbracket X^{3}\rrbracket$ can be only realized as a space-time random distribution and point evaluation for fixed times is not well defined. Thus, $\llbracket X^{3}\rrbracket$ was not included in the statement of Theorem \ref{thm:renorm43}. However, it is not needed in the subsequent analysis of the  $d=3$ parabolic case.
 \end{remark}}


\section{Elliptic $\Phi^4_4$ model}
\label{sec:44}

The goal of this section is  threefold. First, we derive a suitable decomposition of the elliptic $\Phi^{4}$ model \eqref{eq:phi4} in dimension 4. Second, we establish a priori estimates for the involved quantities. This will also serve as a basis for the investigation of the parabolic $\Phi^{4}$ model in dimension 2, see Section \ref{sec:42}. Finally, we employ Schaefer's fixed point theorem together with compactness arguments in order to construct solutions to the decomposed elliptic system. 


\subsection{Decomposition into simpler equations}
\label{ssec:decomp}

We  study the elliptic equation
\begin{equation*}%\label{eq:44}
(- \Delta + \mu) \varphi + \varphi^3 - 3 a \varphi - \xi = 0
\end{equation*}
in $\mathbbm{R}^4$ where $\xi$ is a space white noise and $a$ stands for a renormalization constant needed to define the stochastic objects below. We let $(- \Delta +
\mu) = \Q$ and introduce the ansatz
\[ \varphi = X + \phi + \psi, \]
with
\begin{equation*}%\label{eq:44a}
 \Q X = \xi, \qquad \llbracket X^3 \rrbracket \assign X^3 - 3 a X, \quad
   \llbracket X^2 \rrbracket \assign X^2 - a.
\end{equation*}
Consequently,
\[ 0 = \Q \varphi + \varphi^3 - 3 a \varphi - \xi = \Q \phi + \Q \psi +
   \llbracket X^3 \rrbracket + 3 (\phi + \psi) \llbracket X^2 \rrbracket + 3
   (\phi + \psi)^2 X + (\phi + \psi)^3 . \]
This equation will be decomposed into a system of equations, namely,
\begin{equation}\label{eq:44b}
\Q \phi + \Phi = 0, \qquad \Q \psi + \psi^3 + \Psi = 0,
\end{equation}
where in $\Phi$ we collect all the contributions of negative regularity and in
$\Psi$ all the others (belonging locally to $L^{\infty}$). In addition, with the help of  the operators $\UU_{\leqslant},
\UU_{>}$ defined in Section~\ref{ssec:local}, we localize all the irregular contributions. Namely, each irregular term depending on $\phi+\psi$ will be decomposed into two parts: one even more irregular but controlled by the $L^{\8}$-norm of $\phi+\psi$; and its regular counterpart, which will be included into $\Psi$. This step will be beneficial for the a priori estimates in Section \ref{ssec:apr} as it allows to estimate $\phi$ easily and therefore  eliminate various norms of $\phi$ from the estimates of $\psi$. In other words, thanks to the localizers $\UU_{\leqslant}, \UU_{>}$ we are able to decouple \eqref{eq:44b} and develop an efficient approach towards a priori estimates.

\rmbb{To be more precise,  recall that the operators $\UU_{\leqslant}, \UU_{>}$ depend on a given parameter $L>0$, which has to be chosen appropriately. Moreover, we will choose different values of $L$ for different stochastic objects while keeping in mind that $\UU_{>}$ and $\UU_{\leq}$ of one object shall be given by the same parameter $L$ in order to maintain $\UU_{>}
+\UU_{\leqslant} = \mathrm{Id}$. For the moment, we keep these parameters fixed but arbitrary and their precise values will be determined  below in Section \ref{ssec:apr}.}


Including the localizers, we define
\begin{align}\label{eq:44c}
\begin{aligned}
\Phi &:= \llbracket X^3 \rrbracket + 3 (\phi + \psi) \prec \UU_{>}
   \llbracket X^2 \rrbracket + 3 (\phi + \psi)^2 \prec \UU_{>} X,\\
    \Psi &:= \Psi_1 + \Psi_2, \\
     \Psi_1 &:= \phi^3 + 3 \psi \phi^2 + 3 \psi^2 \phi,\\
     \Psi_2 &:= 3 (\phi + \psi) \prec \UU_{\leqslant} \llbracket X^2
   \rrbracket + 3 (\phi + \psi)^2 \prec \UU_{\leqslant} X + 3 (\phi +
   \psi) \succcurlyeq \llbracket X^2 \rrbracket + 3 (\phi + \psi)^2
   \succcurlyeq X.
\end{aligned}
\end{align}   
 
\subsection{A priori estimates}
\label{ssec:apr}
 
 \rmbb{Let us fix a constant $K>0$ to be chosen  after \eqref{eq:44e} based on the $L^{\infty}$-norm of $\phi+\psi$. Given this value of $K$, we  now determine the values of $L$ in the localization of $X$ and $\llbracket X^{2}\rrbracket$ appearing in \eqref{eq:44c}.}
To this end, recall that the stochastic objects can be constructed so that
$$
\|X\|_{\CC^{-\kappa}(\rho^\sigma)},\|\llbracket X^2 \rrbracket\|_{\CC^{-\kappa}(\rho^\sigma)},\|\llbracket X^3 \rrbracket\|_{\CC^{-\kappa}(\rho^\sigma)}\lesssim 1
$$
provided $\rho$ is a polynomial weight of the form $\rho(x)=\langle x\rangle^{-\nu}$ for some $\nu>0$ and $\kappa,\sigma>0$. 
Hence in view of Lemma~\ref{lem:local} we can choose small parameters $\alpha>\kappa>0$ and $\delta=2-\kappa-\alpha>0$, $\beta=\alpha-\kappa>0$ to set up the localization operators so that (in the sequel, the parameter $\sigma$ is always positive but may change from bound to bound)
$$
\| \UU_{>} X \|_{\CC^{\alpha-2} (\rho^{- 1})} \lesssim 2^{- \delta K}\|X\|_{\CC^{-\kappa}(\rho^\sigma)},
   \qquad \| \UU_{\leqslant} X \|_{\CC^{\alpha} (\rho^{1})}
   \lesssim 2^{(\alpha+\kappa) K}\|X\|_{\CC^{-\kappa}(\rho^\sigma)},
$$
and
\[ \| \UU_{>} \llbracket X^2 \rrbracket \|_{\CC^{\alpha-2}} \lesssim
   2^{- \delta K/2}\|\llbracket X^2 \rrbracket\|_{\CC^{-\kappa}(\rho^\sigma)}, \qquad \| \UU_{\leqslant} \llbracket X^2 \rrbracket
   \|_{\CC^{\alpha} (\rho^{2 })} \lesssim 2^{(\alpha+\kappa) K/2}\|\llbracket X^2 \rrbracket\|_{\CC^{-\kappa}(\rho^\sigma)}. \]
\rmbb{Remark that we chose different values of the parameter $L$ in the localization of $X$ and $\llbracket X^{2}\rrbracket$, namely, $L=K$ for $X$ and $L=K/2$ for $\llbracket X^{2}\rrbracket$.}
From this we have
$$
\| \Phi \|_{\CC^{\alpha - 2} (\rho)} \lesssim \|\llbracket X^3 \rrbracket\|_{\CC^{-\kappa}(\rho^\sigma)} + \|\phi+\psi\|_{L^\infty(\rho)}\| \UU_{>} \llbracket X^2 \rrbracket \|_{\CC^{\alpha - 2}}+\|\phi+\psi\|_{L^\infty(\rho)}^2\| \UU_{>} X \|_{\CC^{\alpha - 2} (\rho^{- 1})}
$$
$$
\lesssim 1+ \|\phi+\psi\|_{L^\infty(\rho)} 2^{-(2-\kappa-\alpha)K/2}+\|\phi+\psi\|_{L^\infty(\rho)}^2 2^{-(2-\kappa-\alpha)K}.
$$
Hence it follows from the Schauder estimates that
\begin{equation}\label{eq:44e}
\|\phi\|_{\CC^\alpha(\rho)}\lesssim \| \Phi \|_{\CC^{\alpha - 2} (\rho)} \lesssim 1+  \|\phi+\psi\|_{L^\infty(\rho)} 2^{-(2-\kappa-\alpha)K/2}+\|\phi+\psi\|_{L^\infty(\rho)}^2 2^{-(2-\kappa-\alpha)K}.
\end{equation}


\rmbb{This leads us to the precise choice of the parameter $K$. In particular, we recall that since the equation for $\phi+\psi$ does not contain any localizers, the norm $\|\phi+\psi\|_{L^{\infty} (\rho)}$ does not depend on the particular choice of the localizers and is finite by assumption.
Let $K>0$ be such that $1+\|\phi+\psi\|_{L^\infty(\rho)}= 2^{(2-\kappa-\alpha)K/2}$. Then  in view of the embedding \eqref{eq:emb} we deduce from \eqref{eq:44e}
\begin{equation*}%\label{eq:LL43}
\| \phi \|_{L^{\infty} (\rho)} + \| \phi \|_{\CC^{\alpha} (\rho^{1 +
   \alpha})}+ \| \phi \|_{\CC^{\beta} (\rho^{1 +
   \beta})} \lesssim 1
   \end{equation*}
  where the constant on the right hand side depends on the noise terms but  is independent of $ K$.
  Consequently,
  $$2^{(2-\kappa-\alpha)K/2}\leq 1+\|\phi\|_{L^\infty(\rho)}+\|\psi\|_{L^\infty(\rho)}\lesssim 1+\|\psi\|_{L^\infty(\rho)}$$
 and hence $2^{(\alpha+\kappa)K}\lesssim 1+ \|\psi\|_{L^\infty(\rho)}^{\varepsilon} $ for some $\varepsilon\in (0,1)$ independent of $K$. The implicit constant (here and in the sequel) is also independent of $K$. The parameter $K$ remains fixed for the rest of the analysis of the elliptic $\Phi^{4}_{4}$ model.} 
% 

To proceed with our a priori estimate, we have
\begin{align}\label{eq:psi1}
\begin{aligned}
\| \Psi_1 \|_{\CC^{\beta} (\rho^{3 + \beta})}& \lesssim (1+\|\psi\|_{L^\infty(\rho)})(1+\|\psi\|_{\CC^{\beta}(\rho^{1+\beta})})\\
\| \Psi_1 \|_{L^{\infty} (\rho^3)} &\lesssim  1+
   \| \psi \|^2_{L^{\infty}
   (\rho)}
\end{aligned}
\end{align}
and
   \[ \| \Psi_2 \|_{\CC^{\beta} (\rho^{3 + \beta})} \lesssim (1+\|\psi\|_{L^\infty(\rho)} ) \|\UU_{\leq}\llbracket X^2\rrbracket\|_{\CC^{\beta}(\rho^{2+\beta})}+(1+\|\psi\|^2_{L^\infty(\rho)} ) \|\UU_{\leq} X\|_{\CC^{\beta}(\rho^{1+\beta})}\]
   \[+(1+\|\psi\|_{\CC^\alpha(\rho^{1+\alpha})} ) \|\llbracket X^2\rrbracket\|_{\CC^{-\kappa}(\rho^{2-\kappa})}+(1+\|\psi\|_{\CC^{\alpha}(\rho^{1+\alpha})} )(1+\|\psi\|_{L^\infty(\rho)}) \| X\|_{\CC^{-\kappa}(\rho^{1-\kappa})}\]
   \begin{equation}\label{eq:psi2}
   \lesssim 1+ \|\psi\|_{L^\infty(\rho)}^{2+\varepsilon}+\|\psi\|_{\CC^\alpha(\rho^{1+\alpha})}(1+\|\psi\|_{L^\infty(\rho)}),
   \end{equation}
   which implies also
\begin{equation}\label{eq:psi22}
\| \Psi_2 \|_{L^{\infty} (\rho^3)} \lesssim 1+\|\psi\|_{L^\infty(\rho)}^{2+\varepsilon}+\|\psi\|_{\CC^\alpha(\rho^{1+\alpha})}(1+\|\psi\|_{L^\infty(\rho)}). 
\end{equation}
Thus, according to Lemma \ref{lemma:schauder-ellptic}
\[
\|\psi\|_{\CC^{2+\beta}(\rho^{3+\beta})}\lesssim 1+\|\psi\|_{\CC^\alpha(\rho^{1+\alpha})}(1+\|\psi\|_{L^\infty(\rho)})+\|\psi\|_{L^\infty(\rho)}^{3+\beta} .\]
 Since we have due to Lemma \ref{lemma:interp}
 \[
\|\psi\|_{\CC^{\alpha}(\rho^{1+\alpha})} \lesssim \|\psi\|_{L^\infty(\rho)}^{1-\alpha/(2+\beta)}\|\psi\|^{\alpha/(2+\beta)}_{\CC^{2+\beta}(\rho^{3+\beta})},
\]
the weighted Young inequality yields
\[
\|\psi\|_{\CC^{2+\beta}(\rho^{3+\beta})} \lesssim 1+\|\psi\|_{L^\infty(\rho)}^{3+\beta}.
\]
Now, it holds
 \[ 
 \|\Psi\|_{L^\infty(\rho^3)}\lesssim 1+\|\psi\|_{L^\infty(\rho)}^{2+\varepsilon}
+\|\psi\|_{\CC^\alpha(\rho^{1+\alpha})} (1 + \|\psi\|_{L^\infty(\rho)}) \]
\[
\lesssim 1+\|\psi\|_{L^\infty(\rho)}^{2+\varepsilon}
+\|\psi\|_{L^\infty(\rho)}^{1-\alpha/(2+\beta)}\|\psi\|^{\alpha/(2+\beta)}_{\CC^{2+\beta}(\rho^{3+\beta})}(1 + \|\psi\|_{L^\infty(\rho)})
\]
\begin{equation}\label{eq:psi3}
\lesssim 1+\|\psi\|^{2+\varepsilon}_{L^\infty(\rho)}
\end{equation}
for some $\varepsilon\in (0,1)$.
Consequently, Lemma \ref{lemma:apriori-elliptic} implies 
\[
\|\psi\|_{L^\infty(\rho)}\lesssim 1+ \|\psi\|^{1-\varepsilon}_{L^\infty(\rho)}
\]
for some $\varepsilon\in (0,1)$. Therefore we deduce
\[
   \|\psi\|_{L^\infty(\rho)}\lesssim 1.
   \]








\subsection{Existence}
\label{ssec:ex44}


As the first step towards the  existence of solutions to the elliptic $\Phi^{4}$ model \eqref{eq:phi4} in dimension~$4$, we consider the problem on a large torus $\mathbb{T}_{M}^{4}$ of a fixed size $M\in\N$. As observed in Section \ref{ssec:decomp}, it reduces to solving the system \eqref{eq:44b}, \eqref{eq:44c} with the space white noise $\xi$ as well as the probabilistic objects $X,\llbracket X^{2}\rrbracket$ and $\llbracket X^{3}\rrbracket$ replaced by their periodic approximations $\xi_{M},X_{M},\llbracket X_{M}^{2}\rrbracket$ and $\llbracket X_{M}^{3}\rrbracket$. We refer to Section~\ref{ssec:renorm-el} for details of the probabilistic construction.




The proof of existence will be divided into two steps. First, we construct a suitable fixed point map
$$\mathcal{K}:\CC^{\beta}(\mathbb{T}_{M}^{4})\times \CC^{\beta}(\mathbb{T}_{M}^{4})\to\CC^{\beta}(\mathbb{T}_{M}^{4})\times \CC^{\beta}(\mathbb{T}_{M}^{4}).$$
Second, we apply Schaefer's fixed point theorem \cite[Section 9.2.2, Theorem 4]{E98} to  show that $\mathcal{K}$ has a fixed point.
More precisely, we define the mapping $\mathcal{K}$ as follows: given
$$(\tilde\phi,\tilde\psi)\in \CC^{\beta}(\mathbb{T}_{M}^{4})\times \CC^{\beta}(\mathbb{T}_{M}^{4}),$$
let $\mathcal{K}(\tilde\phi,\tilde\psi)=(\phi,\psi)$ be a solution to 
\begin{equation}\label{eq:44d}
\Q \phi + \Phi(\tilde\phi,\tilde\psi) = 0, \qquad \Q \psi + \psi^3 + \Psi(\tilde\phi,\tilde\psi) = 0,
\end{equation}
where
\begin{align*}
\begin{aligned}
\Phi(\tilde\phi,\tilde\psi)&= \llbracket X^3 \rrbracket + 3 (\tilde\phi + \tilde\psi) \prec \UU_{>}
   \llbracket X^2 \rrbracket + 3 (\tilde\phi + \tilde\psi)^2 \prec \UU_{>} X,\\
\Psi(\tilde\phi,\tilde\psi)&= \Psi_1(\tilde\phi,\tilde\psi) + \Psi_2(\tilde\phi,\tilde\psi), \\
 \Psi_{1}(\tilde\phi,\tilde\psi)&=\tilde\phi^3 + 3 \tilde\psi \tilde\phi^2 + 3 \tilde\psi^2 \tilde\phi,\\
\Psi_{2}(\tilde\phi,\tilde\psi)&= 3 (\tilde\phi + \tilde\psi) \prec \UU_{\leqslant} \llbracket X^2
   \rrbracket + 3 (\tilde\phi + \tilde\psi)^2 \prec \UU_{\leqslant} X + 3 (\tilde\phi +
   \tilde\psi) \succcurlyeq \llbracket X^2 \rrbracket + 3 (\tilde\phi + \tilde\psi)^2
   \succcurlyeq X.
\end{aligned}
\end{align*}   
Note that the first equation in \eqref{eq:44d} always has a (unique) solution $\phi$ which belongs to $\CC^{\alpha}(\mathbb{T}_{M}^{4})$  due to \eqref{eq:44e}. Indeed, in view of the given regularity of $(\tilde\phi,\tilde\psi)$ and the estimates from Section \ref{ssec:apr} imply  (recall that $\alpha=\beta+\kappa$)
$$
\|\Phi(\tilde\phi,\tilde\psi)\|_{\CC^{\alpha-2}(\mathbb{T}_{M}^{4})}\leq c(\|\tilde\phi+\tilde\psi\|_{\CC^{\beta}(\mathbb{T}_{M}^{4})}).
$$
Next,  we observe that due to \eqref{eq:psi1}, \eqref{eq:psi2} (performed on $\mathbb{T}^{4}_{M}$) the term $\Psi(\tilde\phi,\tilde\psi)$ belongs to $\CC^{\gamma}(\mathbb{T}^{4}_{M})$ provided   $(\tilde\phi,\tilde\psi)\in \CC^{\beta}(\mathbb{T}^{4}_{M})\times \CC^{\beta} (\mathbb{T}^{4}_{M})$ and $\gamma=\beta-\kappa$. Hence according to Proposition~\ref{prop:aux-1} there exists $\psi$ which is a unique classical solution of the second equation in \eqref{eq:44d} and belongs to $\CC^{2+\gamma}(\mathbb{T}_{M}^{4})$. This shows that the map $\mathcal{K}$ is well-defined.
As the next step, we will show that the map $\mathcal{K}$ has a fixed point.


\begin{proposition}\label{prop:fixp}
There exists $(\phi,\psi)\in\CC^{\beta}(\mathbb{T}_{M}^{4})\times \CC^{\beta}(\mathbb{T}_{M}^{4})$ such that $(\phi,\psi)=\mathcal{K}(\phi,\psi).$ Moreover, $(\phi,\psi)$ belongs to $\CC^{\alpha} (\mathbb{T}_{M}^{4})\times \CC^{2+\beta}(\mathbb{T}_{M}^{4})$ for $\alpha=\beta+\kappa$.
\end{proposition}


\begin{proof}
We intend to apply the Schaefer's fixed point theorem which can be found in \cite[Section 9.2.2, Theorem 4]{E98}. To this end, it is necessary to verify that the map $\mathcal{K}$ is continuous and compact and the set
\begin{equation}\label{eq:44f}
\{(\phi,\psi)\in\CC^{\beta}(\mathbb{T}_{M}^{4})\times \CC^{\beta}(\mathbb{T}_{M}^{4});\, (\phi,\psi)=\lambda\mathcal{K}(\phi,\psi) \text{ for some } 0\leqslant \lambda\leqslant 1\}
\end{equation}
is bounded.

\emph{Continuity and compactness:} Assume that $(\tilde\phi_{n},\tilde\psi_{n})\to(\tilde\phi,\tilde\psi)$ in $\CC^{\beta}(\mathbb{T}_{M}^{4})\times \CC^{\beta}(\mathbb{T}_{M}^{4})$ and denote $(\phi_{n},\psi_{n})=\mathcal{K}(\tilde\phi_{n},\tilde\psi_{n}).$ First, we observe that a slight modification of \eqref{eq:psi1}, \eqref{eq:psi2} and \eqref{eq:psi22} shows that
\begin{equation}\label{eq:44ha}
\|\Psi(\tilde\phi_{n},\tilde\psi_{n})\|_{\CC^{\gamma}(\mathbb{T}_{M}^{4})}\leqslant c\big(\|\tilde\phi_{n}+\tilde\psi_{n}\|_{\CC^{\beta}(\mathbb{T}_{M}^{4})}\big)\lesssim 1
\end{equation}
uniformly in $n$. Hence due to the Schauder estimates and Lemma \ref{lemma:schauder-ellptic}, it follows
\begin{equation}\label{eq:44hb}
\|\phi_{n}\|_{\CC^{\alpha}(\mathbb{T}_{M}^{4})}+\|\psi_{n}\|_{\CC^{2+\gamma}(\mathbb{T}_{M}^{4})}\lesssim 1
\end{equation}
uniformly in $n$.
According to the compact embedding \eqref{eq:emb} we deduce that there exists a subsequence still denoted by $(\phi_{n},\psi_{n})$ which converges to certain $(\phi,\psi)$ in $\CC^{\beta}(\mathbb{T}_{M}^{4})\times \CC^{\beta}(\mathbb{T}_{M}^{4})$. Moreover, due to the  uniform bound \eqref{eq:44hb}, it holds
\begin{equation*}
\|\phi\|_{\CC^{\alpha}(\mathbb{T}_{M}^{4})}+\|\psi\|_{\CC^{2+\gamma}(\mathbb{T}_{M}^{4})}\lesssim 1.
\end{equation*}
Since $\Phi$ as well as $\Psi$ in \eqref{eq:44d} depends continuously on $(\tilde\phi_{n},\tilde\psi_{n})$, which can be seen by similar estimates as in Section \ref{ssec:apr}, we may pass to the limit and conclude that $(\phi,\psi)=\mathcal{K}(\tilde\phi,\tilde\psi)$. In view of uniqueness, we deduce that every subsequence converges to the same limit which implies that the whole sequence converges and the desired  continuity of $\mathcal{K}$ follows.
Furthermore, compactness of  $\mathcal{K}$ is also a direct consequence of the bound \eqref{eq:44hb}.



\emph{Boundedness of \eqref{eq:44f}:}  If $(\phi,\psi)=\lambda\mathcal{K}(\phi,\psi) $ for some $0< \lambda\leqslant 1$, then $(\lambda^{-1}\phi,\lambda^{-1}\psi)=\mathcal{K}(\phi,\psi) $ hence
\begin{equation}\label{eq:44g}
\Q \phi + \lambda \Phi(\phi,\psi) = 0, \qquad \Q \psi + \frac{1}{\lambda^{2}}\psi^3 + \lambda\Psi(\phi,\psi) = 0.
\end{equation}
We shall modify the a priori estimates from Section \ref{ssec:apr} in order to account for the parameter $\lambda$ and obtain bounds uniform in $\lambda$. First, we observe that the first equation in \eqref{eq:44g} does not cause any difficulties as $\|\lambda\Phi(\phi,\psi)\|_{\CC^{\alpha-2}(\mathbb{T}_{M}^{4})}\leqslant \|\Phi(\phi,\psi)\|_{\CC^{\alpha-2}(\mathbb{T}_{M}^{4})}$. Consequently, as in  \eqref{eq:44e} we deduce that
$$
\|\phi\|_{\CC^{\alpha}(\mathbb{T}_{M}^{4})}\lesssim 1
$$
uniformly in $\lambda$. The same approach can be applied to the bounds \eqref{eq:psi1}, \eqref{eq:psi2} and \eqref{eq:psi22} which remain unchanged and independent of  $ \lambda$. Revisiting the proof of Lemma \ref{lemma:schauder-ellptic} we obtain
\begin{equation}\label{eq:44h}
\|\psi\|_{\CC^{2+\beta}(\mathbb{T}_{M}^{4})}\lesssim \|\Psi\|_{\CC^{\beta}(\mathbb{T}_{M}^{4})}+\frac{1}{\lambda^{2+\beta}}\|\psi\|_{L^{\infty}(\mathbb{T}_{M}^{4})}^{3+\beta}.
\end{equation}
In order to control the right hand side uniformly in $\lambda$ we revisit the proof of Lemma \ref{lemma:apriori-elliptic} and observe that it simplifies since the weight is not needed on the  torus. Then we apply  \eqref{eq:psi3} and we obtain
$$
\|\psi\|_{L^{\infty}(\mathbb{T}_{M}^{4})}\leqslant \lambda\|\Psi\|^{{1/3}}_{L^{\infty}(\mathbb{T}_{M}^{4})}\lesssim \lambda\big(1+\|\psi\|_{L^{\infty}(\mathbb{T}_{M}^{4})}^{1-\varepsilon}\big)
$$
for some $\varepsilon\in (0,1)$.
Hence, by the weighted Young inequality, we deduce
$$
\|\psi\|_{L^{\infty}(\mathbb{T}_{M}^{4})}\lesssim\lambda.
$$
Plugging this into \eqref{eq:44h} and using the bound for $\Psi$ in \eqref{eq:psi1}, \eqref{eq:psi2} leads to
$$
\|\psi\|_{\CC^{2+\beta}(\mathbb{T}_{M}^{4})}\lesssim 1,
$$
uniformly in $\lambda$ and the boundedness of \eqref{eq:44f} follows.

Finally, Schaefer's fixed point theorem \cite[Section 9.2.2, Theorem 4]{E98} gives the existence of a fixed point of $\mathcal{K}$. Moreover, the a priori estimates from Section \ref{ssec:apr} show that $\psi\in \CC^{2+\beta}(\mathbb{T}^{4}_{M}).$
\end{proof}


Therefore, we have proved the following result.

\begin{theorem}\label{thm:ex-44-per}
Let $M\in\mathbb{N}$. Let $\kappa, \alpha\in (0,1)$ be chosen sufficiently small and let $\beta=\alpha-\kappa>0$.  There exists $(\phi,\psi)\in \CC^{\alpha} (\mathbb{T}_{M}^{4})\times \CC^{2+\beta}(\mathbb{T}_{M}^{4})$ which is a solution to \eqref{eq:44b}, \eqref{eq:44c} on $\mathbb{T}_{M}^{4}$.
\end{theorem}

With this in hand, we are able to conclude the proof of existence on $\R^{4}$.

\begin{theorem}\label{thm:ex43aaa}
Let $\kappa, \alpha\in (0,1)$ be chosen sufficiently small and let $\beta=\alpha-\kappa>0$. There exists $ (\phi,\psi)\in \CC^{\alpha}(\rho)\times[\CC^{2+\beta}(\rho^{3+\beta})\cap L^{\infty}(\rho)]$  which is a solution to \eqref{eq:44b}, \eqref{eq:44c} on $\R^{4}$.
\end{theorem}

\begin{proof}
Let  $(\phi_{M},\psi_{M})\in \CC^{\alpha}(\mathbb{T}^{4}_{M})\times\CC^{2+\beta}(\mathbb{T}^{4}_{M})$ denote the solution to \eqref{eq:44b}, \eqref{eq:44c}  constructed in Theorem \ref{thm:ex-44-per}.
Since functions on $\mathbb{T}^{4}_{M}$ can be regarded as periodic functions defined on the full space $\R^{4}$, we may apply the a priori estimates from Section \ref{ssec:apr}. More precisely, in view of Theorem~\ref{thm:renorm}, we conclude that the approximate solutions $(\phi_{M},\psi_{M})$ are bounded uniformly in $M$ in $ \CC^{\alpha}(\rho)\times[\CC^{2+\beta}(\rho^{3+\beta})\cap L^{\infty}(\rho)]$ whenever $\rho$ is a polynomial bound. Due to \eqref{eq:emb}, this space is compactly embedded into $\CC^{\alpha'}(\rho^{1+\alpha'})\times\CC^{2+\beta'}(\rho^{3+\beta''})$ provided $\alpha'<\alpha$ and $\beta'<\beta<\beta''$. Therefore, there exists a subsequence, still denoted $(\phi_{M},\psi_{M})$ which converges in $\CC^{\alpha'}(\rho^{1+\alpha'})\times\CC^{2+\beta'}(\rho^{3+\beta''})$ to certain $(\phi,\psi)\in  \CC^{\alpha}(\rho)\times[\CC^{2+\beta}(\rho^{3+\beta})\cap L^{\infty}(\rho)]$. Passing to the limit in  \eqref{eq:44b}, \eqref{eq:44c} concludes the proof of existence on the full space.
\end{proof}


Finally, we note that a priori we do not know whether the solution constructed in Theorem~\ref{thm:ex43aaa} is a well-defined random variable, that is, if it is measurable with respect to $\omega$ in the underlying probability space. Indeed, the Schaefer's fixed theorem used in Proposition~\ref{prop:fixp} does not guarantee measurability. However, the existence of a measurable selection can be shown by means of Filippov's implicit function theorem \cite[Theorem 18.17]{AB06}.




\section{Elliptic $\Phi^4_5$ model}
\label{sec:45}

In this section we focus on the elliptic $\Phi^{4}$ model \eqref{eq:phi4} in dimension 5.  First we decompose the equation into a system of equations and establish a priori estimates for the involved quantities. Due to the lower regularity of the driving noise, the analysis is more involved than in Section~\ref{sec:42}. In particular, it is necessary to include additional paracontrolled ansatz, which allows to cancel certain irregular term. Consequently, the a priori estimates become rather delicate and are presented in Sections \ref{ssec:phi1}, \ref{ssec:phi2}, \ref{ssec:theta}, \ref{ssec:psi1}, \ref{ssec:psi2} below.
 This will also serve as a basis for the investigation of the parabolic $\Phi^{4}$ model in dimension 3 in Sections \ref{sec:43}, \ref{s:uniq}, \ref{s:d}.


\subsection{Decomposition into simpler equations}
\label{ssec:dec45}

We  study the elliptic equation
\begin{equation}\label{eq:45}
(- \Delta + \mu) \varphi + \varphi^3 + (- 3 a + 3b) \varphi - \xi = 0 
\end{equation}
in $\mathbbm{R}^5$ where $\xi$ is a space white noise and $a,b$ stand for renormalization constants. We let $(- \Delta +
\mu) = \Q$ and introduce the ansatz
\[ \varphi = X - \tthreeone{X} + \phi + \psi\]
with
\[ \Q X = \xi, \quad \llbracket X^3 \rrbracket \assign X^3 - 3 a X, \quad \llbracket X^2
   \rrbracket \assign X^2 - a, \]
\[ \Q \tthreeone{X} = \llbracket X^3 \rrbracket, \quad \Q
   \ttwoone{X} = \llbracket X^2 \rrbracket, \]
\[\tthreetwor{X}=\tthreeone{X}\circ X,\qquad \ttwothreer{X} = \ttwoone{X} \circ \llbracket X^2 \rrbracket -\frac{b}{3}, \qquad
   \tthreethreer{X} = \tthreeone{X} \circ \llbracket X^2 \rrbracket -  bX. \]
   
   
Recall that if $\rho$ is a polynomial weight of the form $\rho(x)=\langle x\rangle^{-\nu}$ for some   $\nu>0$ and $\sigma>0$ then these objects can be constructed in spaces $\CC^{\alpha}(\rho^\sigma)$ where the respective values of $\alpha$ are given in  Table~\ref{t:reg}. The parameter $\kappa>0$ can be chosen arbitrarily small. 

As a consequence, the left hand side of \eqref{eq:45} rewrites as
\begin{align}\label{eq:rhs45}
\begin{aligned}
&\Q \varphi + \varphi^3 + (- 3 a + 3b) \varphi - \xi \\
&= \Q\phi+\Q\psi +3\llbracket X^2 \rrbracket(-\tthreeone{X}+\phi+\psi)+3X(-\tthreeone{X}+\phi+\psi)^2+(-\tthreeone{X}+\phi+\psi)^3+3b\varphi.
\end{aligned}
\end{align}
Our goal is to construct $\psi$ with regularity $2+\alpha$ whereas $\phi$ will be of regularity $\frac{1}{2}+\alpha$ for some $\alpha>0$ small. Consequently, the third  term on the right hand side of \eqref{eq:rhs45} is not expected to be well-defined and difficulties also arise in the fourth term. In order to cancel the most irregular part of the third term, we assume further that $\phi$ is paracontrolled by $\ttwoone{X}$, namely, it holds
\begin{align}\label{eq:th}
\phi=\vartheta -3(-\tthreeone{X} + \phi + \psi)\prec \ttwoone{X}
\end{align}
for some $\vartheta$ which is more regular (we will see below that $\vartheta$ has the regularity $1+\alpha$). Hence, \eqref{eq:45} rewrites as
\begin{align}\label{eq:rhs45a}
\begin{aligned}
0&= \Q\vartheta+\Q\psi +3\llbracket X^2 \rrbracket\preccurlyeq(-\tthreeone{X}+\phi+\psi)-3[\Q,(-\tthreeone{X}+\phi+\psi)\prec]\ttwoone{X}\\
&\quad+3X(-\tthreeone{X}+\phi+\psi)^2+(-\tthreeone{X}+\phi+\psi)^3+3b\varphi.
\end{aligned}
\end{align}

\subsection{Including the localizers}
\label{s:phi}

\rmbb{In this subsection, we introduce a decomposition of \eqref{eq:rhs45} into two equations. To this end, we adopt the following strategy: As the first step,  we decompose the right hand side of \eqref{eq:rhs45} into four parts: in magenta we collect all the contributions of negative regularity containing only various versions of $X$ (belonging at least to $\rmm{\CC^{-1-\kappa}}$), in orange we collect all the terms of negative regularity depending on $\phi+\psi$ (belonging also to $\rmg{\CC^{-1-\kappa}}$), the blue color denotes all the terms belonging locally to $\rmb{L^{\infty}}$ and we keep the term $\psi^3$ separate.
As the next step, we then further decompose each orange term, namely, into a sum of irregular magenta terms (depending also on $\phi+\psi$) and regular blue terms. This leads to the final decomposition \eqref{eq:two} below.}

Within the first step, we write 
\begin{align}\label{eq:2}
&3\llbracket X^2 \rrbracket\succ(-\tthreeone{X}+\phi+\psi)=-\rmm{3\llbracket X^2 \rrbracket\succ\tthreeone{X}}+\rmg{3\llbracket X^2 \rrbracket\succ(\phi+\psi)},
\end{align}
\begin{align*}
&3\llbracket X^2 \rrbracket\preccurlyeq(-\tthreeone{X}+\phi+\psi)=-\rmm{3\llbracket X^2 \rrbracket\prec\tthreeone{X}}+\rmg{3\llbracket X^2 \rrbracket\prec(\phi+\psi)}+3\llbracket X^2 \rrbracket\circ(-\tthreeone{X}+\phi)+\rmb{3\llbracket X^2 \rrbracket\circ\psi}.
\end{align*}
Now we add the last term from the right hand side of \eqref{eq:rhs45a} to obtain
\begin{align*}
&3\llbracket X^2 \rrbracket\circ(-\tthreeone{X}+\phi)+3b\varphi=- 3\tthreethreer{X}+3\llbracket X^2 \rrbracket\circ\phi+3b( - \tthreeone{X} + \phi + \psi)\\
&\quad= - 3\tthreethreer{X}+3\llbracket X^2 \rrbracket\circ\vartheta-9\llbracket X^2 \rrbracket\circ(( - \tthreeone{X} + \phi + \psi)\prec \ttwoone{X})\\
&\qquad\qquad+3b( - \tthreeone{X} + \phi + \psi)\\
&\quad= - 3\tthreethreer{X}+3\llbracket X^2 \rrbracket\circ\vartheta-9( - \tthreeone{X} + \phi + \psi)(\ttwoone{X}\circ\llbracket X^2 \rrbracket)\\
&\qquad\qquad-9\mathrm{com}(-\tthreeone{X} + \phi + \psi,\ttwoone{X},\llbracket X^2 \rrbracket)+3b( - \tthreeone{X} + \phi + \psi)\\
&\quad= - \rmm{3\tthreethreer{X}}+\rmb{3\llbracket X^2 \rrbracket\circ\vartheta}-\rmg{3( - \tthreeone{X} + \phi + \psi)\ttwothreer{X}}\\
&\qquad\qquad-\rmb{9\mathrm{com}(-\tthreeone{X} + \phi + \psi,\ttwoone{X},\llbracket X^2 \rrbracket)}.
\end{align*}
%
Next, we have
\begin{align*}
&3X(-\tthreeone{X}+\phi+\psi)^2=3X(\tthreeone{X})^2-6X\tthreeone{X}(\phi+\psi)+3X(\phi+\psi)^2\\
&=3X\succ(\tthreeone{X})^2+3X\prec(\tthreeone{X})^2+3X\circ(\tthreeone{X})^2\\
&\quad-6X\succ(\tthreeone{X}(\phi+\psi))-6X\prec(\tthreeone{X}(\phi+\psi))-6X\circ(\tthreeone{X}(\phi+\psi))\\
&\quad+3X(\phi+\psi)^2,
\end{align*}
where
$$
3X\circ(\tthreeone{X})^2=6X\circ(\tthreeone{X}\prec \tthreeone{X})+3X\circ (\tthreeone{X}\circ \tthreeone{X})=6\tthreeone{X}\tthreetwor{X}+6\mathrm{com}(\tthreeone{X},\tthreeone{X},X)+3X\circ (\tthreeone{X}\circ \tthreeone{X}).
$$
%Here,
%$$\|R_1(\tthreeone{X})\|_{\CC^{1-2\kappa}(\rho^{2\sigma})}=\|(\tthreeone{X})^2-2\tthreeone{X}\prec \tthreeone{X}\|_{\CC^{1-2\kappa}(\rho^{2\sigma})}\lesssim \|\tthreeone{X}\|^2_{\CC^{\frac{1}{2}-\kappa}(\rho^\sigma)}.$$
Similarly we decompose
$$
6X\circ(\tthreeone{X}\succ(\phi+\psi))=6(\phi+\psi)\tthreetwor{X}+6\mathrm{com}(\phi+\psi,\tthreeone{X},X),
$$
and observe that all the other terms are well-defined.
Thus we obtain
\begin{align*}
3X(-\tthreeone{X}+\phi+\psi)^2&=
\rmm{3X\succ(\tthreeone{X})^2}+\rmm{3X\prec(\tthreeone{X})^2}+\rmm{6\tthreeone{X}\tthreetwor{X}}\\
&\quad+\rmb{6\mathrm{com}(\tthreeone{X},\tthreeone{X},X)}+\rmb{3X\circ (\tthreeone{X}\circ \tthreeone{X})}\\
&\quad-\rmg{6X\succ(\tthreeone{X}(\phi+\psi))}-\rmg{6X\prec(\tthreeone{X}(\phi+\psi))}-\rmb{6X\circ(\tthreeone{X}\preccurlyeq(\phi+\psi))}\\
&\quad-\rmg{6(\phi+\psi)\tthreetwor{X}}-\rmb{6\mathrm{com}(\phi+\psi,\tthreeone{X},X)}+\rmg{3X(\phi+\psi)^2}.
\end{align*}
For the remaining term in \eqref{eq:rhs45a} we write
$$
(-\tthreeone{X}+\phi+\psi)^3=\rmb{(-\tthreeone{X}+\phi)^3}+\rmb{3(-\tthreeone{X}+\phi)^2\psi}+\rmb{3(-\tthreeone{X}+\phi)\psi^2}+\psi^3
$$


As the next step, we refine the above decomposition even further. To be more precise, we employ the localization  operators $\UU_{>}
$ and $\UU_{\leqslant}$ such that $\UU_{>}
+\UU_{\leqslant} = \mathrm{Id}$  (see Section \ref{ssec:local} for their construction) and carefully separate certain contributions of the orange terms above. \rmbb{We point out that the localizers depend on a parameter $L>0$ whose precise value will be determined below in Section \ref{ssec:phi1}. Moreover, we will choose different values of $L$ for different stochastic objects while keeping in mind that $\UU_{>}$ and $\UU_{\leq}$ of one object are given by the same parameter $L$ in order to maintain $\UU_{>}
+\UU_{\leqslant} = \mathrm{Id}$.}

\rmbb{Following the  regularity rules outlined at the beginning of Section \ref{s:phi}}, all the orange terms will be written as a sum of  magenta and blue terms, which will lead to our final decomposition.  Namely,
\begin{align*}
&\rmg{3\llbracket X^2 \rrbracket\succ(\phi+\psi)}=\rmm{3\UU_>\llbracket X^2 \rrbracket\succ(\phi+\psi)}+\rmb{3\UU_\leq\llbracket X^2 \rrbracket\succ(\phi+\psi)}
\end{align*}
\begin{align*}
\rmg{3\llbracket X^2 \rrbracket\prec(\phi+\psi)}&=\rmm{3\UU_>\llbracket X^2 \rrbracket\prec(\phi+\psi)}+\rmb{3\UU_\leqslant\llbracket X^2 \rrbracket\prec(\phi+\psi)},
\end{align*}
\begin{align*}
-\rmg{3( - \tthreeone{X} + \phi + \psi)\ttwothreer{X}}&=\rmm{3\tthreeone{X}\ttwothreer{X}}-\rmm{3( \phi + \psi)\prec\UU_>\ttwothreer{X}}\\
&\quad-\rmb{3( \phi + \psi)\prec\UU_\leqslant\ttwothreer{X}}-\rmb{3( \phi + \psi)\succcurlyeq\ttwothreer{X}},
\end{align*}
\begin{align}\label{eq:92}
\rmg{6X\succ(\tthreeone{X}(\phi+\psi))}=\rmm{6\UU_>X\succ(\tthreeone{X}(\phi+\psi))}+\rmb{6\UU_\leqslant X\succ(\tthreeone{X}(\phi+\psi))},
\end{align}
\begin{align}\label{eq:93}
\rmg{6X\prec(\tthreeone{X}(\phi+\psi))}=\rmm{6\UU_> X\prec(\tthreeone{X}(\phi+\psi))}+\rmb{6\UU_\leqslant X\prec(\tthreeone{X}(\phi+\psi))},
\end{align}
\begin{align*}
\rmg{6(\phi+\psi)\tthreetwor{X}}=\rmm{6(\phi+\psi)\prec\UU_{>}\tthreetwor{X}}+\rmb{6(\phi+\psi)\prec\UU_{\leqslant}\tthreetwor{X}}+\rmb{6(\phi+\psi)\succcurlyeq\tthreetwor{X}},
\end{align*}
\begin{align}\label{eq:91}
\rmg{3X(\phi+\psi)^2}=\rmm{3\UU_> X\succ(\phi+\psi)^2}+\rmb{3\UU_\leqslant X\succ (\phi+\psi)^2}+\rmb{3 X\preccurlyeq(\phi+\psi)^2}.
\end{align}

As mentioned above,  the concrete choice of the localizers $\UU_\leq$, $\UU_{>}$ in the above changes from line to line.  In particular, it will be seen below that the localization of $X$ in \eqref{eq:92}, \eqref{eq:93} is different from \eqref{eq:91}. The precise choice of these parameters will be made in Section \ref{ssec:phi1} below.

Now, let $\rmm{\Phi}$ be the sum of all the  magenta terms above and $\rmb{\Psi}$ the sum of all the blue terms. \rmbb{More precisely,} 
\begin{align*}
\rmm{\Phi}&:=-\rmm{3\llbracket X^2 \rrbracket\succ\tthreeone{X}}-\rmm{3\llbracket X^2 \rrbracket\prec\tthreeone{X}}- \rmm{3\tthreethreer{X}}+\rmm{3X\succ(\tthreeone{X})^2}+\rmm{3X\prec(\tthreeone{X})^2}+\rmm{6\tthreeone{X}\tthreetwor{X}}\\
&\qquad+\rmm{3\UU_>\llbracket X^2 \rrbracket\succ(\phi+\psi)}+\rmm{3\UU_>\llbracket X^2 \rrbracket\prec(\phi+\psi)}+\rmm{3\tthreeone{X}\ttwothreer{X}}-\rmm{3( \phi + \psi)\prec\UU_>\ttwothreer{X}}\\
&\qquad+\rmm{6\UU_> X\prec(\tthreeone{X}(\phi+\psi))}+\rmm{6\UU_> X\prec(\tthreeone{X}(\phi+\psi))}+\rmm{6(\phi+\psi)\prec\UU_{>}\tthreetwor{X}}\\
&\qquad+\rmm{3\UU_> X\succ(\phi+\psi)^2},
\end{align*}
\begin{align*}
\rmb{\Psi}&:=\rmb{3\llbracket X^2 \rrbracket\circ\psi}+\rmb{3\llbracket X^2 \rrbracket\circ\vartheta}-\rmb{9\mathrm{com}(-\tthreeone{X} + \phi + \psi,\ttwoone{X},\llbracket X^2 \rrbracket)}+\rmb{6\mathrm{com}(\tthreeone{X},\tthreeone{X},X)}\\
&\qquad+\rmb{3X\circ (\tthreeone{X}\circ \tthreeone{X})}-\rmb{6X\circ(\tthreeone{X}\preccurlyeq(\phi+\psi))}-\rmb{6\mathrm{com}(\phi+\psi,\tthreeone{X},X)}\\
&\qquad+\rmb{(-\tthreeone{X}+\phi)^3}+\rmb{3(-\tthreeone{X}+\phi)^2\psi}+\rmb{3(-\tthreeone{X}+\phi)\psi^2}\\
&\qquad +\rmb{3\UU_\leq\llbracket X^2 \rrbracket\succ(\phi+\psi)}
+\rmb{3\UU_\leqslant\llbracket X^2 \rrbracket\prec(\phi+\psi)}-\rmb{3( \phi + \psi)\prec\UU_\leqslant\ttwothreer{X}}-\rmb{3( \phi + \psi)\succcurlyeq\ttwothreer{X}}\\
&\qquad +\rmb{6\UU_\leqslant X\prec(\tthreeone{X}(\phi+\psi))}+\rmb{6\UU_\leqslant X\prec(\tthreeone{X}(\phi+\psi))}\\
&\qquad+\rmb{6(\phi+\psi)\prec\UU_{\leqslant}\tthreetwor{X}}+\rmb{6(\phi+\psi)\succcurlyeq\tthreetwor{X}}+\rmb{3\UU_\leqslant X\succ (\phi+\psi)^2}+\rmb{3 X\preccurlyeq(\phi+\psi)^2}.
\end{align*}
We require that, separately,
\begin{align}\label{eq:two}
 \Q \phi + \rmm{\Phi} = 0, \qquad \Q \psi + \psi^3 + \rmb{\Psi} = 0. 
 \end{align}
Note that in order to have the term $\llbracket X^2\rrbracket\circ\vartheta$ well-defined, it is necessary that $\vartheta$ is at least of regularity $1+\alpha$ for some $\alpha>0.$  This will be shown below. 




\subsection{Bound for $\phi$ in $\CC^\alpha(\rho)$}
\label{ssec:phi1}

At this point we only consider the equation for $\phi$ and intend to show that it belongs to $\CC^{\alpha}(\rho)$ for some $\alpha>0$.
Therefore we aim to estimate $\rmm{\Phi}$ in $\CC^{-2+\alpha}(\rho)$. Recall that before we included the localization operators $\UU_>,$ $\UU_\leqslant$ above, all the magenta and all the orange terms were actually better, namely, of regularity at least $-1-\kappa$. Thanks to the operator $\UU_>$ we are able to profit from this difference of actual and wanted regularity. More precisely, we gain a small factor in all the terms in $\rmm{\Phi}$ containing $\phi+\psi$. As a consequence, a suitable choice of the parameter $L$ in the construction of $\UU_>,$ $\UU_\leqslant,$ yields  a bound for $\rmm{\Phi}$ that only depends on the data of the problem but not on the solution. 

\rmbb{Fix a parameter $K>0$ which will be chosen at the end of this subsection depending on the $L^{\infty}$-norm of $\phi+\psi$. As the next step, given $K$, we shall determine the precise value of $L$ for each application of the localization operators $\UU_{\leq},\UU_{>}$.} First of all, we observe that all the magenta terms that do not contain $\phi+\psi$ can be bounded in $\CC^{-1-\kappa}(\rho^\sigma)$. For the remaining terms, it holds
\begin{align*}
&\|3\UU_>\llbracket X^2 \rrbracket\succ(\phi+\psi)\|_{\CC^{-2+\alpha}(\rho)}+\|3\UU_>\llbracket X^2 \rrbracket\prec(\phi+\psi)\|_{\CC^{-2+\alpha}(\rho)}\\&\qquad\lesssim \|\phi+\psi\|_{L^\infty(\rho)}\|\UU_>\llbracket X^2 \rrbracket\|_{\CC^{-2+\alpha}}\lesssim 2^{-(1-\alpha-\kappa)K} \|\phi+\psi\|_{L^\infty(\rho)},
\end{align*}
provided
\begin{equation}\label{eq:l1}
\|\UU_> \llbracket X^2 \rrbracket\|_{\CC^{-2+\alpha}}\lesssim 2^{-(1-\alpha-\kappa)K} \|\llbracket X^2 \rrbracket\|_{\CC^{-1-\kappa}(\rho^\sigma)}.
\end{equation}
Similarly,
\begin{align*}
\|3( \phi + \psi)\prec\UU_>\ttwothreer{X}\|_{\CC^{-2+\alpha}(\rho)}
\lesssim  \|\phi+\psi\|_{L^\infty(\rho)}\|\UU_>\ttwothreer{X}\|_{\CC^{-2+\alpha}}\lesssim 2^{-(2-\alpha-\kappa)K/2}\|\phi+\psi\|_{L^\infty(\rho)},
\end{align*}
provided
\begin{equation}\label{eq:l2}
\|\UU_> \ttwothreer{X}\|_{\CC^{-2+\alpha}}\lesssim 2^{-(2-\alpha-\kappa)K/2} \|\ttwothreer{X} \|_{\CC^{-\kappa}(\rho^\sigma)},
\end{equation}
\begin{align*}
&\|6\UU_>X\succ(\tthreeone{X}(\phi+\psi))\|_{\CC^{-2+\alpha}(\rho)}+\|6\UU_> X\prec(\tthreeone{X}(\phi+\psi))\|_{\CC^{-2+\alpha}(\rho)}\\
&\quad\lesssim\|\phi+\psi\|_{L^\infty(\rho)}\|\UU_>X\|_{\CC^{-2+\alpha}(\rho^{-\alpha})}\\
&\quad\lesssim 2^{-(\frac{3}{2}-\alpha-\kappa)\frac{2}{3}K}\|\phi+\psi\|_{L^\infty(\rho)},
\end{align*}
provided
\begin{equation}\label{eq:l3}
\|\UU_> X\|_{\CC^{-2+\alpha}(\rho^{-\alpha})}\lesssim 2^{-(\frac{3}{2}-\alpha-\kappa)\frac{2}{3}K} \|X\|_{\CC^{-\frac{1}{2}-\kappa}(\rho^\sigma)},
\end{equation}
\begin{align*}
\|6(\phi+\psi)\prec\UU_{>}\tthreetwor{X}\|_{\CC^{-2+\alpha}(\rho)}
&\lesssim \|\phi+\psi\|_{L^\infty(\rho)}\|\UU_>\tthreetwor{X}\|_{\CC^{-2+\alpha}}\\
&\lesssim 2^{-(2-\alpha-\kappa)K/2}\|\phi+\psi\|_{L^\infty(\rho)},
\end{align*}
provided
\begin{equation}\label{eq:l4}
\|\UU_> \tthreetwor{X}\|_{\CC^{-2+\alpha}}\lesssim 2^{-(2-\alpha-\kappa)K/2} \|\tthreetwor{X} \|_{\CC^{-\kappa}(\rho^\sigma)}.
\end{equation}
and finally
\begin{align*}
\|3\UU_> X\succ(\phi+\psi)^2\|_{\CC^{-2+\alpha}(\rho)}
&\lesssim \|\phi+\psi\|^2_{L^\infty(\rho)}\|\UU_>X\|_{\CC^{-2+\alpha}(\rho^{-1})}\\
&\lesssim 2^{-(\frac{3}{2}-\alpha-\kappa)\frac{4}{3}K}\|\phi+\psi\|_{L^\infty(\rho)}^2,
\end{align*}
provided
\begin{equation}\label{eq:l5}
\|\UU_> X\|_{\CC^{-2+\alpha}(\rho^{-1})}\lesssim 2^{-(\frac{3}{2}-\alpha-\kappa)\frac{4}{3}K} \|X\|_{\CC^{-\frac{1}{2}-\kappa}(\rho^\sigma)}.
\end{equation}

In view of Lemma \ref{lem:local}, once the weight $\rho$ is fixed, the value of $L$ completely determines how the associated localizers $\UU_>$ and $\UU_{\leq}$ are defined. The above considerations  and in particular \eqref{eq:l1}, \eqref{eq:l2}, \eqref{eq:l3}, \eqref{eq:l4}, \eqref{eq:l5} lead us  the  values of $L$ for various objects in our expansion summarized  (in chronological order) in Table \ref{t:loc}.

  \begin{table}
    \begin{center}
  \begin{tabular}{| r | c | c | c | c|c|}
  \hline
  Object &  $ \llbracket X^2 \rrbracket$ & $ \ttwothreer{X}$  & $X$ &  $\tthreetwor{X}$ & $ X$\\ \hline 
  $L$ & $K$ & $\frac{K}{2}$ & $\frac{2K}{3} $ & $\frac{K}{2}$ & $\frac{4K}{3}$  \\
  \hline
  \end{tabular}
  \caption{The value of parameter $L$ used to construct $\UU_>$, $\UU_\leq$.}
  \label{t:loc}
    \end{center}
    \end{table}



Collecting all the above estimates and using the Schauder estimates we  deduce that 
\begin{align*}
&\|\phi\|_{L^\infty(\rho)}\lesssim  \|\phi\|_{\CC^\alpha(\rho)} \\
&\lesssim\|\Phi\|_{\CC^{-2+\alpha}(\rho)}\lesssim 1+2^{-(1-\alpha-\kappa)K}\|\phi+\psi\|_{L^\infty(\rho)}+2^{-(\frac{3}{2}-\alpha-\kappa)\frac{4}{3}K}\|\phi+\psi\|^2_{L^\infty(\rho)}
\end{align*}
\rmbb{\rmbb{and this leads us to a precise choice of the parameter $K>0$. In particular, we proceed as in Section \ref{ssec:apr} and let $K>0$ be such that $1+\|\phi+\psi\|_{L^\infty(\rho)}= 2^{(1-\kappa-\alpha)K}$. Then  in view of the embedding \eqref{eq:emb}}
\begin{equation}\label{eq:10}
\|\phi\|_{L^\infty(\rho)}+ \|\phi\|_{\CC^\alpha(\rho)}\lesssim 1,
\end{equation}
  \rmbb{where the right hand side is independent of $ K$.} So we have that
  \begin{equation}\label{eq:5.15}
  2^{(1-\kappa-\alpha)K}\lesssim 1+\|\psi\|_{L^\infty(\rho)}.
  \end{equation}
 \rmbb{The parameter $K$ as well as the corresponding localizers given according to Table \ref{t:loc} remain fixed for the rest of the analysis of the elliptic $\Phi^{4}_{5}$ model}. }



%Consequently, in view of the embedding \eqref{eq:emb}, whenever $K>0$ is such that $\|\phi+\psi\|_{L^\infty(\rho)}\lesssim 2^{(1-\alpha-\kappa)K}$ it follows
%\begin{equation}\label{eq:10}
%\|\phi\|_{L^\infty(\rho)}+ \|\phi\|_{\CC^\alpha(\rho)}\lesssim 1.
%\end{equation}
%So we can choose $L>0$ such that $\|\phi+\psi\|_{L^\infty(\rho)}\lesssim 1+\|\psi\|_{L^\infty(\rho)}\simeq 2^{(1-\alpha-\kappa)K}$. This value of $L$ as well as the corresponding localizers given according to Table \ref{t:loc} are fixed for the rest of the proof.
% 
 

\subsection{Bound for $\phi$ in $\CC^{\frac{1}{2}+\alpha}(\rho^{\frac{3}{2}+\alpha})$}
\label{ssec:phi2}

As the next step, we estimate $\Phi$ in $\CC^{-\frac{3}{2}+\alpha}(\rho^{\frac{3}{2}+\alpha})$. It will be seen below that all the terms except for the one which is quadratic in $\phi+\psi$ can be even estimated in $\CC^{-\frac{3}{2}+\alpha}(\rho)$. 
 \rmbb{Recall that the parameter $K>0$  fixed in the previous section  determined the value of $L$ for each application of the localization operators, see Table \ref{t:loc}. }
In view of Lemma \ref{lem:local}, the sequence $(L_k)_{k\in\N_0}$ is therefore also fixed (and possibly different for each application of the  localizing operators). \rmbb{Below, we use the bound \eqref{eq:5.15} in order to control various powers of $2^{ K}$ by  $\|\psi\|^{\varepsilon}_{L^{\infty}(\rho)}$, where  $\varepsilon\in(0,1)$ is a generic constant \rmbb{independent of $K$} whose value changes from line to line. Accordingly, Lemma \ref{lem:local} yields}
$$
\|\UU_> \llbracket X^2 \rrbracket\|_{\CC^{-\frac{3}{2}+\alpha}}\lesssim 2^{-(\frac{1}{2}-\alpha-\kappa)K} \|\llbracket X^2 \rrbracket\|_{\CC^{-1-\kappa}(\rho^\sigma)},
$$
$$
\|\UU_> \ttwothreer{X}\|_{\CC^{-\frac{3}{2}+\alpha}}\lesssim 2^{-(\frac{3}{2}-\alpha-\kappa)K/2} \|\ttwothreer{X} \|_{\CC^{-\kappa}(\rho^\sigma)},
$$
$$
\|\UU_> X\|_{\CC^{-\frac{3}{2}+\alpha}(\rho^{-\alpha})}\lesssim 2^{-(1-\alpha-\kappa)\frac{2}{3}K} \|X\|_{\CC^{-\frac{1}{2}-\kappa}(\rho^\sigma)},
$$
$$
\|\UU_> \tthreetwor{X}\|_{\CC^{-\frac{3}{2}+\alpha}}\lesssim 2^{-(\frac{3}{2}-\alpha-\kappa)K/2} \|\tthreetwor{X} \|_{\CC^{-\kappa}(\rho^\sigma)},
$$
$$
\|\UU_> X\|_{\CC^{-\frac{3}{2}+\alpha}(\rho^{-\frac{1}{2}+\alpha})}\lesssim 2^{-(1-\alpha-\kappa)\frac{4}{3}K} \|X\|_{\CC^{-\frac{1}{2}-\kappa}(\rho^\sigma)},
$$
which implies
\begin{align*}
&\|3\UU_>\llbracket X^2 \rrbracket\succ(\phi+\psi)\|_{\CC^{-\frac{3}{2}+\alpha}(\rho)}+\|3\UU_>\llbracket X^2 \rrbracket\prec(\phi+\psi)\|_{\CC^{-\frac{3}{2}+\alpha}(\rho)}\\
&\quad\lesssim \|\phi+\psi\|_{L^\infty(\rho)}\|\UU_>\llbracket X^2 \rrbracket\|_{\CC^{-\frac{3}{2}+\alpha}}\lesssim 2^{-(\frac{1}{2}-\alpha-\kappa)K} \|\phi+\psi\|_{L^\infty(\rho)}\\
&\quad\lesssim 2^{-(\frac{1}{2}-\alpha-\kappa)K}2^{(1-\alpha-\kappa)K}\lesssim 1+\|\psi\|_{L^\infty(\rho)}^{\varepsilon},
\end{align*}
\begin{align*}
\|3( \phi + \psi)\prec\UU_>\ttwothreer{X}\|_{\CC^{-\frac{3}{2}+\alpha}(\rho)}
&\lesssim  \|\phi+\psi\|_{L^\infty(\rho)}\|\UU_>\ttwothreer{X}\|_{\CC^{-\frac{3}{2}+\alpha}}\lesssim 2^{-(\frac{3}{2}-\alpha-\kappa)K/2}\|\phi+\psi\|_{L^\infty(\rho)}\\
&\lesssim  1+\|\psi\|_{L^\infty(\rho)}^{\varepsilon},
\end{align*}
\begin{align*}
&\|6\UU_>X\succ(\tthreeone{X}(\phi+\psi))\|_{\CC^{-\frac{3}{2}+\alpha}(\rho)}+\|6\UU_> X\prec(\tthreeone{X}(\phi+\psi))\|_{\CC^{-\frac{3}{2}+\alpha}(\rho)}\\
&\quad\lesssim  \|\phi+\psi\|_{L^\infty(\rho)}\|\UU_>X\|_{\CC^{-\frac{3}{2}+\alpha}(\rho^{-\alpha})}\\
&\quad\lesssim 2^{-(1-\alpha-\kappa)\frac{2}{3}K}\|\phi+\psi\|_{L^\infty(\rho)}\lesssim 1+\|\psi\|_{L^\infty(\rho)}^{\varepsilon},
\end{align*}
\begin{align*}
\|6(\phi+\psi)\prec\UU_{>}\tthreetwor{X}\|_{\CC^{-\frac{3}{2}+\alpha}(\rho)}
&\lesssim \|\phi+\psi\|_{L^\infty(\rho)}\|\UU_>\tthreetwor{X}\|_{\CC^{-\frac{3}{2}+\alpha}}\\
&\lesssim 2^{-(\frac{3}{2}-\alpha-\kappa)K/2}\|\phi+\psi\|_{L^\infty(\rho)}\lesssim 1+\|\psi\|_{L^\infty(\rho)}^{\varepsilon}.
\end{align*}
For the last term, we note that a higher power of $\rho$ is necessary and estimate
\begin{align*}
\|3\UU_> X\succ(\phi+\psi)^2\|_{\CC^{-\frac{3}{2}+\alpha}(\rho^{\frac{3}{2}+\alpha})}
&\lesssim \|\phi+\psi\|^2_{L^\infty(\rho)}\|\UU_>X\|_{\CC^{-\frac{3}{2}+\alpha}(\rho^{-\frac{1}{2}+\alpha})}\\
&\lesssim 2^{-(1-\alpha-\kappa)\frac{4}{3}K}\|\phi+\psi\|_{L^\infty(\rho)}^2\lesssim 1+\|\psi\|_{L^\infty(\rho)}^{\varepsilon},
\end{align*}

To summarize, collecting all the above estimates and using the Schauder estimates we  deduce that 
\begin{align}\label{eq:11}
& \|\phi\|_{\CC^{\frac{1}{2}+\alpha}(\rho^{\frac{3}{2}+\alpha})}\lesssim\|\Phi\|_{\CC^{-\frac{3}{2}+\alpha}(\rho^{\frac{3}{2}+\alpha})}\lesssim 1+\|\psi\|_{L^\infty(\rho)}^{\varepsilon}.
\end{align}




\subsection{Bound for $\vartheta$ in $\CC^{1+\alpha}(\rho^{2+\alpha})$}
\label{ssec:theta}

As the next step, we derive a bound for $\vartheta$ in $\CC^{1+\alpha}(\rho^{2+\alpha})$ which will be needed in the sequel in order to control $\psi$. In view of the paracontrolled ansatz \eqref{eq:th}, equation \eqref{eq:rhs45a} as well as the decomposition \eqref{eq:two}, we observe that the most irregular part of $\Phi$, namely the two magenta terms coming from \eqref{eq:2}, cancel out, and additionally the blue term coming from \eqref{eq:2} and a commutator appear. More precisely,  $\vartheta$ solves
$$
\Q\vartheta + \Theta =0,
$$
with
$$
\Theta=\Phi +\rmm{3\llbracket X^2 \rrbracket\succ\tthreeone{X}}-\rmm{3\UU_>\llbracket X^2 \rrbracket\succ(\phi+\psi)}+\rmb{3\UU_\leq\llbracket X^2 \rrbracket\succ(\phi+\psi)}-3[\Q,(-\tthreeone{X}+\phi+\psi)\prec]\ttwoone{X}.
$$
Next, we observe that all the remaining terms from $\Phi$ can be estimated in $\CC^{-1+\alpha}(\rho^{2+\alpha})$. Indeed, all the terms that do not contain $\phi+\psi$ are bounded in this space and for the terms containing $\phi+\psi$, we observe that
$$
\|\UU_> \llbracket X^2 \rrbracket\|_{\CC^{-\frac{3}{2}}(\rho^{\frac{1}{2}})}\lesssim 2^{-(\frac{1}{2}-\kappa)K} \|\llbracket X^2 \rrbracket\|_{\CC^{-1-\kappa}(\rho^\sigma)},
$$
$$
\|\UU_\leq\llbracket X^2 \rrbracket\|_{\CC^{-1+\alpha}(\rho^{1+\alpha})}\lesssim 2^{(\alpha+\kappa)K} \|\llbracket X^2 \rrbracket\|_{\CC^{-1-\kappa}(\rho^{\sigma})},
$$
$$
\|\UU_> \ttwothreer{X}\|_{\CC^{-1+\alpha}}\lesssim 2^{-(1-\alpha-\kappa)K/2} \|\ttwothreer{X} \|_{\CC^{-\kappa}(\rho^\sigma)},
$$
$$
\|\UU_> X\|_{\CC^{-1+\alpha}(\rho^{-\alpha})}\lesssim 2^{-(\frac{1}{2}-\alpha-\kappa)\frac{2}{3}K} \|X\|_{\CC^{-\frac{1}{2}-\kappa}(\rho^\sigma)},
$$
$$
\|\UU_> \tthreetwor{X}\|_{\CC^{-1+\alpha}}\lesssim 2^{-(1-\alpha-\kappa)K/2} \|\tthreetwor{X} \|_{\CC^{-\kappa}(\rho^\sigma)},
$$
$$
\|\UU_> X\|_{\CC^{-1+\alpha}(\rho^\alpha)}\lesssim 2^{-(\frac{1}{2}-\alpha-\kappa)\frac{4}{3}K} \|X\|_{\CC^{-\frac{1}{2}-\kappa}(\rho^\sigma)},
$$
and consequently
\begin{align*}
\|3\UU_>\llbracket X^2 \rrbracket\prec(\phi+\psi)\|_{\CC^{-1+\alpha}(\rho^{2+\alpha})}
&\lesssim \|\phi+\psi\|_{\CC^{\frac{1}{2}+\alpha}(\rho^{\frac{3}{2}+\alpha})}\|\UU_>\llbracket X^2 \rrbracket\|_{\CC^{-\frac{3}{2}}(\rho^{\frac{1}{2}})}\\
&\lesssim 2^{-(\frac{1}{2}-\kappa)K}  \|\phi+\psi\|_{\CC^{\frac{1}{2}+\alpha}(\rho^{\frac{3}{2}+\alpha})}\lesssim 1+\|\psi\|_{L^\infty(\rho)}^{\varepsilon}+ \|\psi\|_{\CC^{\frac{1}{2}+\alpha}(\rho^{\frac{3}{2}+\alpha})}.
\end{align*}
\begin{align*}
\|3( \phi + \psi)\prec\UU_>\ttwothreer{X}\|_{\CC^{-1+\alpha}(\rho)}
&\lesssim  \|\phi+\psi\|_{L^\infty(\rho)}\|\UU_>\ttwothreer{X}\|_{\CC^{-1+\alpha}}\lesssim 2^{-(1-\alpha-\kappa)K/2}\|\phi+\psi\|_{L^\infty(\rho)}\\
&\lesssim 1+\|\psi\|_{L^\infty(\rho)}^{\varepsilon},
\end{align*}
\begin{align*}
&\|6\UU_>X\succ(\tthreeone{X}(\phi+\psi))\|_{\CC^{-1+\alpha}(\rho)}+\|6\UU_> X\prec(\tthreeone{X}(\phi+\psi))\|_{\CC^{-1+\alpha}(\rho)}\\
&\quad\lesssim \|\phi+\psi\|_{L^\infty(\rho)}\|\UU_>X\|_{\CC^{-1+\alpha}(\rho^{-\alpha})}\\
&\quad\lesssim 2^{-(\frac{1}{2}-\alpha-\kappa)\frac{2}{3}K}\|\phi+\psi\|_{L^\infty(\rho)}\lesssim 1+\|\psi\|_{L^\infty(\rho)}^{\varepsilon},
\end{align*}
\begin{align*}
\|6(\phi+\psi)\prec\UU_{>}\tthreetwor{X}\|_{\CC^{-1+\alpha}(\rho)}
&\lesssim \|\phi+\psi\|_{L^\infty(\rho)}\|\UU_>\tthreetwor{X}\|_{\CC^{-1+\alpha}}\\
&\lesssim 2^{-(1-\alpha-\kappa)K/2}\|\phi+\psi\|_{L^\infty(\rho)}\\
&\lesssim 1+\|\psi\|_{L^\infty(\rho)}^\varepsilon,
\end{align*}
and finally
\begin{align*}
\|3\UU_> X\succ(\phi+\psi)^2\|_{\CC^{-1+\alpha}(\rho^{2+\alpha})}
&\lesssim \|\phi+\psi\|^2_{L^\infty(\rho)}\|\UU_>X\|_{\CC^{-1+\alpha}(\rho^\alpha)}\\
&\lesssim 2^{-(\frac{1}{2}-\alpha-\kappa)\frac{4}{3}K}\|\phi+\psi\|_{L^\infty(\rho)}^2\\
&\lesssim 1+\|\psi\|_{L^\infty(\rho)}^{1+\varepsilon}.
\end{align*}
Hence, we have shown that
\begin{align*}
\|\Phi +3\llbracket X^2 \rrbracket\succ\tthreeone{X}-3\UU_>\llbracket X^2 \rrbracket\succ(\phi+\psi)\|_{\CC^{-1+\alpha}(\rho^{2+\alpha})}\lesssim 1+\|\psi\|_{L^\infty(\rho)}^{1+\varepsilon}+ \|\psi\|_{\CC^{\frac{1}{2}+\alpha}(\rho^{\frac{3}{2}+\alpha})} .
\end{align*}
Similarly,
\begin{align*}
\|3\UU_\leq\llbracket X^2 \rrbracket\succ(\phi+\psi)\|_{\CC^{-1+\alpha}(\rho^{2+\alpha})}&\lesssim \|\phi+\psi\|_{L^\infty(\rho)}\|\UU_\leq\llbracket X^2 \rrbracket\|_{\CC^{-1+\alpha}(\rho^{1+\alpha})}\\
&\lesssim 2^{(\alpha+\kappa)K}(1+\|\psi\|_{L^\infty(\rho)} )\\
&\lesssim 1+\|\psi\|_{L^\infty(\rho)}^{1+\varepsilon},
\end{align*}
and for the commutator, we obtain
\begin{align*}
&\|3[\Q,(-\tthreeone{X}+\phi+\psi)\prec]\ttwoone{X}||_{\CC^{-1+\alpha}(\rho^{2+\alpha})}\\
&\quad=\|3\Delta(-\tthreeone{X}+\phi+\psi)\prec \ttwoone{X}+6(\nabla(-\tthreeone{X}+\phi+\psi))\prec(\nabla \ttwoone{X})\|_{\CC^{-1+\alpha}(\rho^{2+\alpha})}\\
&\quad\lesssim \|-\tthreeone{X}+\phi+\psi\|_{\CC^{\frac{1}{2}-\kappa}(\rho^{\frac{3}{2}+\alpha})}\| \ttwoone{X}\|_{\CC^{1-\kappa}(\rho^\sigma)}\lesssim 1+\|\phi+\psi\|_{\CC^{\frac{1}{2}+\alpha}(\rho^{\frac{3}{2}+\alpha})}\\
&\quad \lesssim 1+\|\psi\|^\varepsilon_{L^\infty(\rho)}+\|\psi\|_{\CC^{\frac{1}{2}+\alpha}(\rho^{\frac{3}{2}+\alpha})}.
\end{align*}

To summarize, we have proved that
\begin{equation}\label{eq:12}
\|\vartheta\|_{\CC^{1+\alpha}(\rho^{2+\alpha})}\lesssim 1 +  \|\psi\|^{1+\varepsilon}_{L^\infty(\rho)}+\|\psi\|_{\CC^{\frac{1}{2}+\alpha}(\rho^{\frac{3}{2}+\alpha})}.
\end{equation}







\subsection{Bound for $\psi$ in $\CC^{2+\gamma}(\rho^{3+\gamma})$}
\label{ssec:psi1}



In this section we make use of the estimates \eqref{eq:10}, \eqref{eq:11}, \eqref{eq:12} in order to estimate  $\psi$ in $\CC^{2+\gamma}(\rho^{3+\gamma})$ for  $\gamma=\alpha-\kappa>0$ such that $\gamma\leq\frac{1}{2}-3\kappa$ (which can be achieved by a suitable choice of $\alpha, \kappa>0$. In view of \eqref{eq:two} and Lemma \ref{lemma:schauder-ellptic}, it is therefore necessary to estimate $\Psi$ in $\CC^{\gamma}(\rho^{3+\gamma})$.
We estimate as follows
\begin{align*}
\|3\llbracket X^2\rrbracket\circ\psi\|_{\CC^\gamma(\rho^{3+\gamma})}+\|3\llbracket X^2\rrbracket\circ\vartheta\|_{\CC^\gamma(\rho^{3+\gamma})}&\lesssim \|\psi\|_{\CC^{1+\alpha}(\rho^{2+\alpha})}+\|\vartheta\|_{\CC^{1+\alpha}(\rho^{2+\alpha})}\\
&\lesssim\|\psi\|_{\CC^{1+\alpha}(\rho^{2+\alpha})}+1+\|\psi\|^{1+\varepsilon}_{L^\infty(\rho)}+\|\psi\|_{\CC^{\frac{1}{2}+\alpha}(\rho^{\frac{3}{2}+\alpha})} ,
\end{align*}
and according to Lemma \ref{lem:com}
\begin{align*}
\|9\mathrm{com}(-\tthreeone{X} + \phi + \psi,\ttwoone{X},\llbracket X^2 \rrbracket)\|_{\CC^\gamma(\rho^{3+\gamma})}&\lesssim 1+\|\phi+\psi\|_{\CC^{\frac{1}{2}}(\rho^{2+\alpha})}\lesssim 1+\|\psi\|_{L^\infty(\rho)}^\varepsilon+\|\psi\|_{\CC^{{1+\alpha}}(\rho^{2+\alpha})},
\end{align*}
\begin{align*}
\|6\mathrm{com}(\tthreeone{X},\tthreeone{X},X)\|_{\CC^\gamma(\rho^{3+\gamma})}\lesssim 1,
\end{align*}
\begin{align*}
\|6X\circ(\tthreeone{X}\preccurlyeq(\phi+\psi))\|_{\CC^\gamma(\rho^{3+\gamma})}\lesssim\|\phi+\psi\|_{\CC^{\frac{1}{2}+\alpha}(\rho^{2+\alpha})}\lesssim  \|\psi\|_{L^\infty(\rho)}^\varepsilon+\|\psi\|_{\CC^{1+\alpha}(\rho^{2+\alpha})},
\end{align*}
\begin{align*}
\|6\mathrm{com}(\phi+\psi,\tthreeone{X},X)\|_{\CC^\gamma(\rho^{3+\gamma})}\lesssim\|\phi+\psi\|_{\CC^{\frac{1}{2}}(\rho^{2+\alpha})}\lesssim \|\psi\|_{L^\infty(\rho)}^\varepsilon+ \|\psi\|_{\CC^{1+\alpha}(\rho^{2+\alpha})},
\end{align*}
\begin{align*}
&\|(-\tthreeone{X}+\phi)^3+3(-\tthreeone{X}+\phi)^2\psi+3(-\tthreeone{X}+\phi)\psi^2\|_{\CC^\gamma(\rho^{3+\gamma})}\\
&\quad\lesssim (1+\|\psi\|_{L^\infty(\rho)})(1+\|\psi\|_{\CC^\gamma(\rho^{1+\gamma})}).
\end{align*}
Next, we observe that due to our choice of $K$ at the end of Section \ref{s:phi}, it holds that $$2^{(1+\gamma+\kappa)K}\simeq 1+\|\psi\|_{L^\infty(\rho)}^{1+\varepsilon},\quad 2^{(\gamma+\kappa)K/2}\simeq 1+\|\psi\|_{L^\infty(\rho)}^\varepsilon,$$
$$
2^{(\frac{1}{2}+\gamma+\kappa)\frac{2}{3}K}\simeq 1+\|\psi\|_{L^\infty(\rho)}^{\varepsilon}, \quad 2^{(\frac{1}{2}+\gamma+\kappa)\frac{4}{3}K}\simeq 1+\|\psi\|_{L^\infty(\rho)}^{\varepsilon}
$$
 for some $\varepsilon\in(0,1)$ (whose value possibly changes from bound to bound); and in view of Table~\ref{t:loc} we have
 $$
\|\UU_\leqslant\llbracket X^2 \rrbracket\|_{\CC^\gamma(\rho^{2})}\lesssim 2^{(1+\gamma+\kappa)K}\|\llbracket X^2 \rrbracket\|_{\CC^{-1-\kappa}(\rho^{\sigma})},
$$
\begin{align*}
\|\UU_\leqslant \ttwothreer{X}\|_{\CC^{\gamma}(\rho^{2+\gamma})}\lesssim 2^{(\gamma+\kappa)K/2}\|\ttwothreer{X}\|_{\CC^{-\kappa}(\rho^\sigma)},
\end{align*}
\begin{align*}
\|\UU_\leqslant X\|_{\CC^\gamma(\rho^{2})}\lesssim 2^{(\frac{1}{2}+\gamma+\kappa)\frac{2}{3}K} \|X\|_{\CC^{-\frac{1}{2}-\kappa}(\rho^\sigma)},
\end{align*}
$$
\|\UU_{\leq}X\|_{\CC^{-\frac{1}{2}-\kappa}(\rho^{\sigma})}\lesssim \|X\|_{\CC^{-\frac{1}{2}-\kappa}(\rho^{\sigma})},
$$
\begin{align*}
\|\UU_\leqslant \tthreetwor{X}\|_{\CC^\gamma(\rho^{2+\gamma})}\lesssim 2^{(\gamma+\kappa)K/2} \| \tthreetwor{X}\|_{\CC^{-\kappa}(\rho^{\sigma})}
\end{align*}
\begin{align*}
\|3\UU_\leqslant X\|_{\CC^\gamma(\rho^{1+\gamma})}\lesssim 2^{(\frac{1}{2}+\gamma+\kappa)\frac{4}{3}K}\|X\|_{\CC^{-\frac{1}{2}-\kappa}(\rho^{\sigma})}
\end{align*}
This leads to
\begin{align*}
\begin{aligned}
\|3\UU_\leqslant\llbracket X^2 \rrbracket\succ(\phi+\psi)\|_{\CC^\gamma(\rho^{3+\gamma})}&\lesssim \|\phi+\psi\|_{L^\infty(\rho)}\|\UU_\leqslant\llbracket X^2 \rrbracket\|_{\CC^\gamma(\rho^{2+\gamma})}\\
&\lesssim 2^{(1+\gamma+\kappa)K}\|\phi+\psi\|_{L^\infty(\rho)}\lesssim 1+\|\psi\|^{2+\varepsilon}_{L^\infty(\rho)},
\end{aligned}
\end{align*}
\begin{align*}
\begin{aligned}
\|3\UU_\leqslant\llbracket X^2 \rrbracket\prec(\phi+\psi)\|_{\CC^\gamma(\rho^{3+\gamma})}&\lesssim\|\phi+\psi\|_{\CC^{\gamma}(\rho^{1+\gamma})} \|\UU_\leqslant\llbracket X^2 \rrbracket\|_{\CC^\gamma(\rho^{2})}\\
&\lesssim 2^{(1+\gamma+\kappa)K}\|\phi+\psi\|_{\CC^{\gamma}(\rho^{1+\gamma})} \lesssim (1+\|\psi\|^{1+\varepsilon}_{L^\infty(\rho)})(1+\|\psi\|_{\CC^{\gamma}(\rho^{1+\gamma})})\\
&\lesssim 1+ \|\psi\|^{1+\varepsilon}_{L^\infty(\rho)} +\|\psi\|^{1+\varepsilon}_{L^\infty(\rho)}\|\psi\|_{\CC^{\gamma}(\rho^{1+\gamma})},
\end{aligned}
\end{align*}
\begin{align*}
\|3( \phi + \psi)\prec\UU_\leqslant\ttwothreer{X}\|_{\CC^\gamma(\rho^{3+\gamma})}&\lesssim \|\phi+\psi\|_{L^\infty(\rho)}\|\UU_\leqslant \ttwothreer{X}\|_{\CC^{\gamma}(\rho^{2+\gamma})}\lesssim 2^{(\gamma+\kappa)K/2}\|\phi+\psi\|_{L^\infty(\rho)}\\
&\lesssim 1+\|\psi\|_{L^\infty(\rho)}^{1+\varepsilon},
\end{align*}
\begin{align*}
\|3( \phi + \psi)\succcurlyeq\ttwothreer{X}\|_{\CC^\gamma(\rho^{3+\gamma})}\lesssim \|\phi+\psi\|_{\CC^{\gamma+\kappa}(\rho^{2+\gamma+\kappa})}\lesssim 1+\|\psi\|_{\CC^{1+\alpha}(\rho^{2+\alpha})},
\end{align*}
\begin{align*}
\|6\UU_\leqslant X\succ(\tthreeone{X}(\phi+\psi))\|_{\CC^\gamma(\rho^{3+\gamma})}&\lesssim \|\phi+\psi\|_{L^\infty(\rho)}\|\UU_\leqslant X\|_{\CC^\gamma(\rho^{2})}\lesssim 2^{(\frac{1}{2}+\gamma+\kappa)\frac{2}{3}K}\|\phi+\psi\|_{L^\infty(\rho)}\\
&\lesssim 1+\|\psi\|_{L^\infty(\rho)}^{1+\varepsilon},
\end{align*}
\begin{align*}
\|6\UU_\leqslant X\prec(\tthreeone{X}(\phi+\psi))\|_{\CC^\gamma(\rho^{3+\gamma})}&\lesssim \|\UU_{\leq}X\|_{\CC^{-\frac{1}{2}-\kappa}(\rho^{\sigma})}\|\phi+\psi\|_{\CC^{\frac{1}{2}+\gamma+\kappa}(\rho^{2+\gamma+\kappa})}\lesssim \|\phi+\psi\|_{\CC^{\frac{1}{2}+\alpha}(\rho^{2+\alpha})}\\
&\lesssim 1+\|\psi\|_{L^\infty(\rho)}^\varepsilon+ \|\psi\|_{\CC^{1+\alpha}(\rho^{2+\alpha})},
\end{align*}
\begin{align*}
\|6(\phi+\psi)\prec\UU_{\leqslant}\tthreetwor{X}\|_{\CC^\gamma(\rho^{3+\gamma})}&\lesssim \|\phi+\psi\|_{L^\infty(\rho)}\|\UU_\leqslant \tthreetwor{X}\|_{\CC^\gamma(\rho^{2+\gamma})}\lesssim 2^{(\gamma+\kappa)K/2} \|\phi+\psi\|_{L^\infty(\rho)}\\
&\lesssim  1+\|\psi\|_{L^\infty(\rho)}^{1+\varepsilon},
\end{align*}
\begin{align*}
\|6(\phi+\psi)\succcurlyeq\tthreetwor{X}\|_{\CC^\gamma(\rho^{3+\gamma})}\lesssim\|\phi+\psi\|_{\CC^{\gamma+\kappa}(\rho^{2+\gamma+\kappa})}\lesssim 1+\|\psi\|_{\CC^{1+\alpha}(\rho^{2+\alpha})},
\end{align*}
\begin{align*}
\begin{aligned}
\|3\UU_\leqslant X\succ (\phi+\psi)^2\|_{\CC^\gamma(\rho^{3+\gamma})}
&\lesssim\|3\UU_\leqslant X\|_{\CC^\gamma(\rho^{1+\gamma})}\|\phi+\psi\|^2_{L^\infty(\rho)}\lesssim 2^{(\frac{1}{2}+\gamma+\kappa)\frac{4}{3}K}\|\phi+\psi\|^2_{L^\infty(\rho)}\\
&\lesssim 1+\|\psi\|_{L^\infty(\rho)}^{2+\varepsilon}
\end{aligned}
\end{align*}
and finally
\begin{align*}
\|3 X\preccurlyeq(\phi+\psi)^2\|_{\CC^\gamma(\rho^{3+\gamma})}
&\lesssim \|(\phi+\psi)^2\|_{\CC^{\frac{1}{2}+\gamma+\kappa}(\rho^{3})}\lesssim \|\phi+\psi\|_{L^\infty(\rho)}\|\phi+\psi\|_{\CC^{\frac{1}{2}+\gamma+\kappa}(\rho^{2})}\\
&\lesssim (1+\|\psi\|_{L^\infty(\rho)})(1+\|\psi\|_{L^\infty(\rho)}^\varepsilon+\|\psi\|_{\CC^{\frac{1}{2}+\alpha}(\rho^{\frac{3}{2}+\alpha})})\\
&\lesssim 1 +\|\psi\|_{L^\infty(\rho)}^{1+\varepsilon}+\|\psi\|_{\CC^{\frac{1}{2}+\alpha}(\rho^{\frac{3}{2}+\alpha})}+\|\psi\|_{L^\infty(\rho)}\|\psi\|_{\CC^{\frac{1}{2}+\alpha}(\rho^{\frac{3}{2}+\alpha})}.
\end{align*}

To summarize, we have shown that
\begin{align*}
\|\Psi\|_{\CC^{\gamma}(\rho^{3+\gamma})}
&\lesssim 1+ \|\psi\|_{\CC^{1+\alpha}(\rho^{2+\alpha})} +\|\psi\|_{\CC^{\frac{1}{2}+\alpha}(\rho^{\frac{3}{2}+\alpha})}
\\
&\quad +\|\psi\|^{2+\varepsilon}_{L^\infty(\rho)} 
+\|\psi\|^{1+\varepsilon}_{L^\infty(\rho)}\|\psi\|_{\CC^{\gamma}(\rho^{1+\gamma})}\\
&\quad +\|\psi\|_{L^\infty(\rho)}\|\psi\|_{\CC^{\frac{1}{2}+\alpha}(\rho^{\frac{3}{2}+\alpha})}.
\end{align*}
Due to interpolation from Lemma \ref{lemma:interp} we estimate
\begin{align*}
\|\psi\|_{\CC^{1+\alpha}(\rho^{2+\alpha})}\lesssim \|\psi\|^{1-\frac{1+\alpha}{2+\gamma}}_{L^\infty(\rho)} \|\psi\|^\frac{1+\alpha}{2+\gamma}_{\CC^{2+\gamma}(\rho^{3+\gamma})},
\end{align*}
\begin{align*}
\|\psi\|_{\CC^{\gamma}(\rho^{1+\gamma})}\lesssim \|\psi\|^{1-\frac{\gamma}{2+\gamma}}_{L^\infty(\rho)} \|\psi\|^\frac{\gamma}{2+\gamma}_{\CC^{2+\gamma}(\rho^{3+\gamma})},
\end{align*}
\begin{align*}
\|\psi\|_{\CC^{\frac{1}{2}+\alpha}(\rho^{\frac{3}{2}+\alpha})}\lesssim \|\psi\|^{1-\frac{\frac{1}{2}+\alpha}{2+\gamma}}_{L^\infty(\rho)} \|\psi\|^{\frac{\frac{1}{2}+\alpha}{2+\gamma}}_{\CC^{2+\gamma}(\rho^{3+\gamma})},
\end{align*}
therefore,  Lemma \ref{lemma:schauder-ellptic} together with the weighted Young inequality implies
\begin{align}\label{eq:97}
\|\psi\|_{\CC^{2+\gamma}(\rho^{3+\gamma})}\lesssim \|\Psi\|_{\CC^{\gamma}(\rho^{3+\gamma})}+\|\psi\|^{3+\gamma}_{L^\infty(\rho)} \lesssim 1 +\|\psi\|^{3+\gamma}_{L^\infty(\rho)} .
\end{align}

\subsection{Bound for $\psi$ in $L^\infty(\rho)$}
\label{ssec:psi2}

As the next step, towards the application of Lemma \ref{lemma:apriori-elliptic}, it is necessary to estimate $\Psi$ in $L^\infty(\rho^3)$. We observe that for most of the terms we may use the estimates above, only the cubic term is estimated as follows
\begin{align*}
&\|(-\tthreeone{X}+\phi)^3+3(-\tthreeone{X}+\phi)^2\psi+3(-\tthreeone{X}+\phi)\psi^2\|_{L^\infty(\rho^{3})}\lesssim 1+\|\psi\|^2_{L^\infty(\rho)},
\end{align*}
and we may also improve the bound
\begin{align*}
\begin{aligned}
\|3\UU_\leqslant\llbracket X^2 \rrbracket\prec(\phi+\psi)\|_{L^\infty(\rho^{3})}&\lesssim\|\phi+\psi\|_{L^\infty(\rho)} \|\UU_\leqslant\llbracket X^2 \rrbracket\|_{\CC^\gamma(\rho^{2})}\\
&\lesssim 2^{(1+\gamma+\kappa)K}\|\phi+\psi\|_{L^\infty(\rho)} \lesssim 1+ \|\psi\|^{2+\varepsilon}_{L^\infty(\rho)}.
\end{aligned}
\end{align*}
Therefore, we deduce
\begin{align*}
\|\Psi\|_{L^\infty(\rho^{3})}
&\lesssim 1+ \|\psi\|_{\CC^{1+\alpha}(\rho^{2+\alpha})} +\|\psi\|_{\CC^{\frac{1}{2}+\alpha}(\rho^{\frac{3}{2}+\alpha})}
\\
&\quad +\|\psi\|^{2+\varepsilon}_{L^\infty(\rho)} 
 +\|\psi\|_{L^\infty(\rho)}\|\psi\|_{\CC^{\frac{1}{2}+\alpha}(\rho^{\frac{3}{2}+\alpha})}
\end{align*}
and applying again the interpolation from Lemma \ref{lemma:interp} together with \eqref{eq:97} leads to
\begin{align*}
\|\Psi\|_{L^\infty(\rho^{3})}
&\lesssim 1+ \|\psi\|_{L^\infty(\rho)}^{2+\varepsilon} .
\end{align*}
Finally, according to Lemma \ref{lemma:apriori-elliptic} and weighted Young inequality we conclude
$$
\|\psi\|_{L^\infty(\rho)}\lesssim 1+ \|\Psi\|_{L^\infty(\rho)}^{1/3}\lesssim 1 + \|\psi\|_{L^\infty(\rho)}^\varepsilon\lesssim 1,
$$
and the proof is complete.

 





\subsection{Existence}
\label{ssec:ex45}

The construction of a solution proceeds similarly to Section~\ref{ssec:ex44}. More precisely,  we first consider the problem on a large torus of size $M$ and establish existence based on Schaefer's fixed point theorem \cite[Section 9.2.2, Theorem 4]{E98}. Then we make use of the a priori estimates from Sections~\ref{ssec:phi1}, \ref{ssec:phi2}, \ref{ssec:theta}, \ref{ssec:psi2}, \ref{ssec:psi2} together with Theorem \ref{thm:renorm2} and a compactness argument to pass to the limit as $M\to\infty$.

Recall that in view of the computations in Sections \ref{ssec:dec45}, \ref{s:phi}, system \eqref{eq:phi4} in dimension~5 reduces to equations \eqref{eq:two}, \eqref{eq:th}.


\begin{theorem}\label{thm:ex43}
Let $\kappa, \alpha\in (0,1)$ be chosen sufficiently small and let $\gamma=\alpha-\kappa>0$. There exists
$$(\phi,\psi)\in  [ \CC^{\frac 12+\alpha}(\rho^{\frac 32+\alpha})\cap \CC^{\alpha}(\rho)]\times[\CC^{2+\gamma}(\rho^{3+\gamma})\cap L^{\infty}(\rho)]$$
which is a solution to \eqref{eq:two}, \eqref{eq:th} on $\R^{5}$.
\end{theorem}

\begin{proof} \emph{Step 1 -- existence on a large torus:}
Similarly to the proof of Theorem \ref{thm:ex-44-per}, we define a fixed point map
$$\mathcal{K}:\CC^{\frac 12+\beta}(\mathbb{T}_{M}^{5})\times \CC^{1+\beta}(\mathbb{T}_{M}^{5})\to\CC^{\frac 12+\beta}(\mathbb{T}_{M}^{5})\times \CC^{1+\beta}(\mathbb{T}_{M}^{5})$$ for a small parameter $\beta\in(0,1)$ as follows: given
$$(\tilde\phi,\tilde\psi)\in \CC^{\frac 12+\beta}(\mathbb{T}_{M}^{5})\times \CC^{ 1+\beta}(\mathbb{T}_{M}^{5}),$$
let $\mathcal{K}(\tilde\phi,\tilde\psi)=(\phi,\psi)$ be a solution to 
\begin{equation}\label{eq:45d}
\Q \phi + \Phi(\tilde\phi,\tilde\psi) = 0, \qquad \Q \psi + \psi^3 + \Psi(\tilde\phi,\tilde\psi) = 0,
\end{equation}
where $\Phi(\tilde\phi,\tilde\psi)$ and $\Psi(\tilde\phi,\tilde\psi)$ contain all the magenta and blue terms from Section \ref{s:phi}, respectively, with $\phi,\psi$ replaced by $\tilde\phi,\tilde
\psi$. 


The first equation in \eqref{eq:45d} always has a (unique) solution $\phi$ which belongs to $\CC^{\alpha}(\mathbb{T}_{M}^{5})$  due to the bounds in Section \ref{ssec:phi1}. Moreover, Section \ref{ssec:phi2} shows that $\phi\in\CC^{\frac 12+\alpha}(\mathbb{T}_{M}^{5})$ and we may choose $\alpha>\beta$. Furthermore, similarly to \eqref{eq:th} we denote
$$
\vartheta:=\phi +3(-\tthreeone{X} + \tilde\phi + \tilde\psi)\prec \ttwoone{X}
$$
and observe that due to Section \ref{ssec:theta} and Section \ref{ssec:psi1} (performed on $\mathbb{T}^{5}_{M}$) the right hand side $\Psi(\tilde\phi,\tilde\psi)$ belongs to $\CC^{\gamma}(\mathbb{T}^{5}_{M})$ provided   $(\tilde\phi,\tilde\psi)\in \CC^{\frac 12+\beta}(\mathbb{T}^{5}_{M})\times \CC^{1+\beta} (\mathbb{T}^{5}_{M})$ and $\gamma=\beta-\kappa$, $\gamma\leq\frac 12-3\kappa$. Hence Proposition \ref{prop:aux-1} implies  existence of a unique classical solution to the second equation in \eqref{eq:45d}. Hence the map $\mathcal{K}$ is well-defined.


Next, we deduce that  the map $\mathcal{K}$ has a fixed point $(\phi,\psi)\in \CC^{\frac 12+\alpha} (\mathbb{T}_{M}^{5})\times \CC^{2+\gamma}(\mathbb{T}_{M}^{5})$ for $\alpha=\beta+\kappa$ and $\gamma=\beta-\kappa$. More precisely, the proof follows the lines of Proposition \ref{prop:fixp} and employs the estimates from Sections \ref{ssec:phi1}, \ref{ssec:phi2},  \ref{ssec:theta}, \ref{ssec:psi1}, \ref{ssec:psi2}. The proof of existence on $\mathbb{T}^{5}_{M}$ is therefore complete.




 \emph{Step 2 -- existence on the full space:} For $M\in\N$ let  $(\phi_{M},\psi_{M})$ denote the solution to \eqref{eq:two}, \eqref{eq:th} on $\mathbb{T}^{5}_{M}$  constructed above.
Then  the a priori estimates from Sections \ref{ssec:phi1}, \ref{ssec:phi2},  \ref{ssec:theta}, \ref{ssec:psi1}, \ref{ssec:psi2} apply and, in view of Theorem~\ref{thm:renorm2}, we conclude that the approximate solutions $(\phi_{M},\psi_{M})$ are bounded uniformly in $M$ in
$$[ \CC^{\frac 12+\alpha}(\rho^{\frac 32+\alpha})\cap \CC^{\alpha}(\rho)]\times[\CC^{2+\gamma}(\rho^{3+\gamma})\cap L^{\infty}(\rho)]$$
whenever $\rho$ is a polynomial bound. Due to \eqref{eq:emb}, this space is compactly embedded into $\CC^{\frac 12+\alpha'}(\rho^{\frac 32+\alpha''})\times\CC^{2+\gamma'}(\rho^{3+\gamma''})$ provided $\alpha'<\alpha<\alpha''$ and $\gamma'<\gamma<\gamma''$. Therefore, there exists a subsequence, still denoted $(\phi_{M},\psi_{M})$ which converges in $\CC^{\frac 12+\alpha'}(\rho^{\frac 32+\alpha''})\times\CC^{2+\gamma'}(\rho^{3+\gamma''})$ to certain
$$(\phi,\psi)\in  [ \CC^{\frac 12+\alpha}(\rho^{\frac 32+\alpha})\cap \CC^{\alpha}(\rho)]\times[\CC^{2+\gamma}(\rho^{3+\gamma})\cap L^{\infty}(\rho)].$$
Passing to the limit in  \eqref{eq:two}, \eqref{eq:th} concludes the proof of existence on the full space.
\end{proof}





\section{Parabolic $\Phi^4_2$ model}
\label{sec:42}

The analysis of the parabolic $\Phi^{4}$ model on $\R^{2}$, that is,
\begin{equation}\label{eq:42}
\partial_{t}\varphi+ (- \Delta + \mu) \varphi + \varphi^3 - 3 a \varphi - \xi = 0,\qquad\varphi(0)=\varphi_{0},
\end{equation}
where $\xi$ is a space-time white noise, is very similar to the elliptic $\Phi^{4}_{4}$ model on $\R^{4}$. Indeed, the regularity of the space white noise in dimension $4$ is the same as the regularity of the space-time white noise in dimension $2$. Without loss of generality we assume that the mass $\mu$ is strictly positive (otherwise we add a linear term with positive mass to both sides of \eqref{eq:42} and consider the original massive term as a right hand side, see Remark \ref{rem:mu}). Then we proceed as in Section~\ref{sec:44}, let $(- \Delta +
\mu) = \Q$ and $\LL=\partial_{t}+\Q\hspace{-.2em}$ and introduce the ansatz
\[ \varphi = X + \phi + \psi, \]
with
\begin{equation*}%\label{eq:42a}
\LL X = \xi, \qquad \llbracket X^3 \rrbracket \assign X^3 - 3 a X, \quad
   \llbracket X^2 \rrbracket \assign X^2 - a.
\end{equation*}
Recall that $X$ is chosen stationary.
This leads us to the system of equations
\begin{equation}\label{eq:42b}
\LL \phi + \Phi = 0, \quad \phi(0)=\phi_{0}=\varphi_{0},\qquad \LL \psi + \psi^3 + \Psi = 0,\quad \psi(0)=0,
\end{equation}
where $\varphi_{0}\in \CC^{\alpha}(\rho_{0})$; and  $\Phi$ and $\Psi$ are given as in \eqref{eq:44c} but employing the parabolic localizers $\VV_{>},\VV_{\leq}$ instead of $\UU_{>},\UU_{\leq}$.



The existence of a solution can now be proved by choosing a smooth and space periodic approximation $\xi_{\varepsilon}$ of the driving space-time white noise $\xi$, defined on the torus of size $\frac 1\varepsilon$ and solving \eqref{eq:42} on the approximate level with the associated renormalization constant $a_{\varepsilon}$. Subsequently, we may pass to the limit using  the above uniform estimates together with  compactness. To be more precise, let $\xi_{\varepsilon}$ be a periodic version of a {space-time mollification} of $\xi$ defined on $[0,\8)\times\mathbb{T}^{2}_{1/\varepsilon}$ and define $X_{\varepsilon}$ as the stationary solution to
\begin{equation*}
\LL X_{\varepsilon} = \xi_{\varepsilon}.
\end{equation*}
The other stochastic objects were defined in Theorem \ref{thm:renorm42}. Throughout this section, $\rho$ denotes   a polynomial space-time weight. 

 


Let $\varphi_{\varepsilon,0}$ be a mollification of the initial condition $\varphi_{0}$. Then according to Proposition \ref{prop:42d}, for every $\varepsilon\in (0,1)$ there exists $\varphi_{\varepsilon}\in C^{\infty}([0,\8)\times \mathbb{T}^{2}_{1/\varepsilon})$ which is the unique classical solution to
\begin{equation*}%\label{eq:42d}
\LL \varphi_{\varepsilon} + \varphi_{\varepsilon}^3 - 3 a_{\varepsilon} \varphi_{\varepsilon} - \xi_{\varepsilon} = 0,\qquad\varphi(0)=\varphi_{\varepsilon,0}.
\end{equation*}
As the next step, we proceed with the same decomposition $\varphi_{\varepsilon}=X_{\varepsilon}+\phi_{\varepsilon}+\psi_{\varepsilon}$ as above, only starting from the mollified version $\xi_{\varepsilon}$ instead of from $\xi$. 
According to Corollary \ref{cor:42d} it holds for every $\varepsilon\in (0,1)$ that $\varphi_{\varepsilon}\in  C\CC^{2 + \kappa} (\rho^{3 + \kappa})\cap C^{1}L^{\infty}(\rho^{3+\kappa}) \cap L^{\infty}L^{\infty} (\rho)$ and the  same regularity holds for $\phi_{\varepsilon},\psi_{\varepsilon}$. Hence we follow the lines of  Section~\ref{ssec:apr} and employ  Lemmas \ref{lemma:interp2}, \ref{lem:local2}, \ref{lemma:sch}, \ref{lemma:schauder-par}, \ref{lemma:apriori-parabolic}, in order to deduce that the following bound holds true uniformly in $\varepsilon\in (0,1)$
\begin{equation}\label{eq:un42}
\|\phi_{\varepsilon}\|_{C\CC^{\alpha}(\rho)}+\|\phi_{\varepsilon}\|_{C^{\alpha/2}L^{\infty}(\rho)}+\|\psi_{\varepsilon}\|_{C\CC^{2+\beta}(\rho^{3+\beta})}+\|\psi_{\varepsilon}\|_{C^{1}L^{\infty}(\rho^{3+\beta})}+\|\psi_{\varepsilon}\|_{L^{\infty}L^{\infty}(\rho)}\lesssim 1.
\end{equation}
Based on this uniform bound we are able to pass to the limit.

\begin{theorem}\label{thm:ex42}
Let $\kappa, \alpha\in (0,1)$ be chosen sufficiently small and let $\beta=\alpha-\kappa>0$. If $\varphi_{0}\in \CC^{\alpha}(\rho_{0})$ then there exists
$$(\phi,\psi)\in [C\CC^{\alpha}(\rho)\cap C^{\alpha/2}L^{\infty}(\rho)] \times[C\CC^{2+\beta}(\rho^{3+\beta})\cap C^{1}L^{\infty}(\rho^{3+\beta})\cap L^{\infty}L^{\infty}(\rho)]$$
which is a solution to  \eqref{eq:42b}.
\end{theorem}

\begin{proof}
Due to \eqref{eq:un42}, the approximate solutions $(\phi_{\varepsilon},\psi_{\varepsilon})$ are bounded uniformly in $\varepsilon
\in (0,1)$ in
$$[C\CC^{\alpha}(\rho)\cap C^{\alpha/2}L^{\infty}(\rho)] \times[C\CC^{2+\beta}(\rho^{3+\beta})\cap C^{1}L^{\infty}(\rho^{3+\beta})\cap L^{\infty}L^{\infty}(\rho)]$$
whenever $\rho$ is a polynomial bound. Due to \eqref{eq:emb}, Arzel\`a-Ascoli and Aubin-Lions-type argument (see \cite[Lemma 1, Theorem 5]{S87}) this space is compactly embedded into
$$[C_{\rm{loc}}\CC^{\alpha-\iota}(\rho^{1+\delta})\cap C^{(\alpha-\iota)/2}_{\rm{loc}}\CC^{-\gamma}(\rho^{1+\delta}) ]\times [C_{\rm{loc}}\CC^{2+\beta-\iota}(\rho^{3+\beta+\delta})\cap C^{1-\iota}_{\rm{loc}}\CC^{-\gamma}(\rho^{3+\beta+\delta}) ]$$ provided  $\iota\in (0,\alpha\wedge\beta\wedge1)$ and $\gamma,\delta\in (0,1)$ are chosen small. Therefore, there exists a subsequence, still denoted $(\phi_{\varepsilon},\psi_{\varepsilon})$ which converges in this space to certain $(\phi,\psi)$ and we intend to pass to the limit in  \eqref{eq:42b}.

To this end, we fix $T>0$. Note that due to Theorem~\ref{thm:renorm42}, the linearity of the localizers $\VV_{>},\VV_{\leq}$ and Lemma \ref{lem:local2} it follows that
\begin{align*}
\VV_{>} X_{\varepsilon}\to \VV_{>} X\quad\text{in}\quad C_{T}\CC^{\alpha-2} (\rho^{- 1}),
   &\qquad \VV_{\leqslant} X_{\varepsilon} \to \VV_{\leqslant} X\quad\text{in}\quad C_{T}\CC^{\alpha} (\rho^{1}),\\
 \VV_{>} \llbracket X_{\varepsilon}^2 \rrbracket \to  \VV_{>} \llbracket X^2 \rrbracket \quad\text{in}\quad C_{T}\CC^{\alpha-2},
 &\qquad  \VV_{\leqslant} \llbracket X_{\varepsilon}^2 \rrbracket\to  \VV_{\leqslant} \llbracket X^2 \rrbracket
  \quad\text{in}\quad C_{T}\CC^{\alpha} (\rho^{2}).
\end{align*}
Note that  we employed the same spaces as for the a priori estimates in Section \ref{ssec:apr}. As a consequence and in view of the estimates from Section \ref{ssec:apr}, we observe that there exists $K>0$ such that
$$\Phi_{\varepsilon}\to\Phi\quad\text{in}\quad C_{T}\CC^{\alpha-2}(\rho^{K}),\qquad \Psi_{\varepsilon}\to\Psi\quad\text{in}\quad C_{T}\CC^{\beta-\iota}(\rho^{K}),$$
where $\Phi,\Psi$ are defined as in \eqref{eq:44c}. The constant $K>0$ needs to be  chosen sufficiently large in order to compensate for the lack of convergence of $\phi_{\varepsilon}$ and $\psi_{\varepsilon}$ in $C_{T}L^{\infty}(\rho)$, which has to be replaced by $C_{T}L^{\infty}(\rho^{1+\delta})$ and $C_{T}L^{\infty}(\rho^{3+\beta+\delta})$, respectively. Passing to the limit in the remaining terms in \eqref{eq:42b} is straightforward, and therefore, the couple $(\phi,\psi)$ solves \eqref{eq:42b}, which is understood in distributional sense.

It remains to show that
$$(\phi,\psi)\in [C\CC^{\alpha}(\rho)\cap C^{\alpha/2}L^{\infty}(\rho)] \times[C\CC^{2+\beta}(\rho^{3+\beta})\cap C^{1}L^{\infty}(\rho^{3+\beta})\cap L^{\infty}L^{\infty}(\rho)].$$
To this end, we observe that according to the above convergence $(\phi_{\varepsilon},\psi_{\varepsilon})\to (\phi,\psi)$ it follows that $\Delta_{i}\phi_{\varepsilon}(t,x)\to \Delta_{i}\phi(t,x)$ and $\Delta_{i}\psi_{\varepsilon}(t,x)\to \Delta_{i}\psi(t,x)$ for every $i\geq -1$ and almost every $(t,x)\in [0,\8)\times\R^{2}$. In addition, the Littlewood--Paley blocks $\Delta_{i}\phi_{\varepsilon}, \Delta_{i}\psi_{\varepsilon}$ satisfy the uniform  bounds (even uniform in $\varepsilon,i$ and $t$)
$$
\|\rho_{t}\Delta_{i}\phi_{\varepsilon}(t)\|_{L^{\infty}}\lesssim 1,\qquad \|\rho_{t}\Delta_{i}\psi_{\varepsilon}(t)\|_{L^{\infty}}\lesssim 1.
$$
Consequently,  $\rho_{t}\Delta_{i}\phi_{\varepsilon}(t)\to \rho_{t}\Delta_{i}\phi(t)$ and $\rho_{t}\Delta_{i}\psi_{\varepsilon}(t)\to \rho_{t}\Delta_{i}\psi(t)$  weak-star in $L^{\infty}(\R^{2})$ for every $i\geq -1$ and almost every $t\in [0,\8)$. Since the $L^{\infty}$-norm is weak-star lower semicontinuous, we obtain
$$
\|\rho_{t}\Delta_{i}\phi(t)\|_{L^{\infty}}\leqslant \liminf_{\varepsilon\to 0}\|\rho_{t}\Delta_{i}\phi_{\varepsilon}(t)\|_{L^{\infty}}\leqslant \liminf_{\varepsilon\to 0}\|\phi_{\varepsilon}\|_{C_{T}\CC^{\alpha}(\rho)}2^{-i\alpha}\lesssim 2^{-i\alpha},
$$
$$
\|\rho_{t}\Delta_{i}\psi(t)\|_{L^{\infty}}\leqslant \liminf_{\varepsilon\to 0}\|\rho_{t}\Delta_{i}\psi_{\varepsilon}(t)\|_{L^{\infty}}\leqslant \liminf_{\varepsilon\to 0}\|\psi_{\varepsilon}\|_{C_{T}L^{\infty}(\rho)}\lesssim 1,
$$
and by the same argument
$$
\|\rho_{t}^{3+\beta}\Delta_{i}\psi(t)\|_{L^{\infty}}\leqslant \liminf_{\varepsilon\to 0}\|\rho_{t}^{3+\beta}\Delta_{i}\psi_{\varepsilon}(t)\|_{L^{\infty}}\leqslant \liminf_{\varepsilon\to 0}\|\psi_{\varepsilon}(t)\|_{C_{T}\CC^{2+\beta}(\rho^{3+\beta})}\lesssim 2^{-i(2+\beta)}.
$$
This implies that
$$(\phi,\psi)\in L^{\infty}\CC^{\alpha}(\rho) \times[L^{\infty}\CC^{2+\beta}(\rho^{3+\beta})\cap  L^{\infty}L^{\infty}(\rho)].$$
Now, using the convergence $(\phi_{\varepsilon},\psi_{\varepsilon})\to (\phi,\psi)$ in $C_{ \rm{loc}}^{\alpha/2}L^{\infty}(\rho)\times C^{1}_{ \rm{loc}}L^{\infty}(\rho)$ we obtain
$$
\|\phi(t)-\phi(s)\|_{L^{\infty}(\rho)}=\lim_{\varepsilon\to 0}\|\phi_{\varepsilon}(t)-\phi_{\varepsilon}(s)\|_{L^{\infty}(\rho)}\leqslant \lim_{\varepsilon\to 0}\|\phi_{\varepsilon}\|_{C^{\alpha/2}L^{\infty}(\rho)}|t-s|^{\alpha/2}\lesssim |t-s|^{\alpha/2}
$$
and similarly for the norm of $\psi$ in $C^{1}_{T}L^{\infty}(\rho^{3+\beta}).$ Hence
$$(\phi,\psi)\in C^{\alpha/2}_{T}L^{\infty}(\rho) \times C^{1}_{T}L^{\infty}(\rho^{{3+\beta}}).$$
Now,  we apply the Schauder estimates for both $\phi$ and $\psi$ (i.e. Lemma \ref{lemma:sch} and Lemma \ref{lemma:schauder-par}) to obtain continuity in time, namely,
$$
(\phi,\psi)\in C_{T}\CC^{\alpha}(\rho)\times C_{T}\CC^{2+\beta}(\rho^{3+\beta}).
$$
The proof is complete.
\end{proof}



 


\section{Parabolic $\Phi^4_3$ model}
\label{sec:43}

We proceed by similar arguments as  in the elliptic $\Phi^4_5$ model discussed in Section \ref{sec:45}. More precisely, we intend to study the parabolic equation
\begin{equation}\label{eq:43}
\partial_{t} \varphi+(- \Delta + \mu) \varphi + \varphi^3 + (- 3 a + 3b) \varphi - \xi = 0,\qquad\varphi(0)=\varphi_{0},
\end{equation}
in $\mathbbm{R}^3$ where $\xi$ is a space-time white noise, $a,b$ stand for renormalization constants. We recall that according to Theorem \ref{thm:renorm43} the renormalization constant $b$ depends on time and is bounded.
Without loss of generality we assume $\mu>0$ (see Remark~\ref{rem:mu}).

While the existence of solutions to the parabolic $\Phi^{4}$ model in dimension 2 was a more or less straightforward consequence of the elliptic a priori estimates for the $\Phi^{4}_{4}$ model from Section \ref{ssec:apr}, the situation is more involved in dimension $3$, which will be seen in the sequel. To be more precise, we let $(- \Delta +
\mu) = \Q$ and $\LL=\partial_{t}+\Q\hspace{-.2em}$
 and we introduce the ansatz
\[ \varphi = X - \tthreeone{X} + \phi + \psi\]
with $X$ being stationary and
\[\LL X = \xi, \quad \llbracket X^2
   \rrbracket \assign X^2 - a,\quad \llbracket X^3 \rrbracket \assign X^3 - 3 a X, \]
\[  \LL \tthreeone{X} = \llbracket X^3 \rrbracket, \quad \tthreeone{X}(0)=0,\qquad\LL
   \ttwoone{X} = \llbracket X^2 \rrbracket,\quad\ttwoone{X}(0)=0, \]
\[\tthreetwor{X}=\tthreeone{X}\circ X,\qquad \ttwothreer{X} = \ttwoone{X} \circ \llbracket X^2 \rrbracket -\frac{b}{3}, \qquad
   \tthreethreer{X} = \tthreeone{X} \circ \llbracket X^2 \rrbracket -  bX. \]
Thus, the left hand side of \eqref{eq:43} rewrites as
\begin{align*}%\label{eq:rhs43}
\begin{aligned}
&\LL \varphi + \varphi^3 + (- 3 a + 3b) \varphi - \xi \\
&\quad= \LL\phi+\LL\psi +3\llbracket X^2 \rrbracket(-\tthreeone{X}+\phi+\psi)+3X(-\tthreeone{X}+\phi+\psi)^2+(-\tthreeone{X}+\phi+\psi)^3+3b\varphi.
\end{aligned}
\end{align*}


In Section \ref{sec:42} we were able to apply the elliptic a priori estimates pointwise in time.
In view of the decomposition of the elliptic $\Phi^{4}_{5}$ model, this is no longer possible here. The difficulty arises in the paracontrolled ansatz similar to \eqref{eq:th} and the associated commutator as in \eqref{eq:rhs45a}. Indeed, since the differential operator in \eqref{eq:43} includes also time derivative, we require sufficient time regularity in order to control the  commutator. To this end, we make use of the time-mollified paraproduct, introduced in Section \ref{ssec:para}, and  apply Lemma \ref{lem:5.1}. More precisely, we assume that $\phi$ is paracontrolled by $\ttwoone{X}$, namely, it holds
\begin{align}\label{eq:th43}
\phi=\vartheta -3(-\tthreeone{X} + \phi + \psi)\Prec \ttwoone{X}
\end{align}
for some $\vartheta$ which is more regular (we will see below that $\vartheta$ has the regularity $1+\alpha$).

Furthermore, it can be seen in Lemma \ref{lemma:schauder-par} that the expected time regularity for $\psi$ is not optimal. Indeed,  we cannot go beyond $C^{1}$ with respect to time, which would be natural for taking the full advantage of the interpolation (in time) from Lemma \ref{lemma:interp2} and mimicking the strategy of Section \ref{sec:45}. Therefore, the terms requiring time regularity of $\phi+\psi$ have to be treated differently. To this end, it is necessary to consider a higher power of the weight $\rho$ in the bounds for $\vartheta$, which will help us compensate for sub-optimal time interpolation. Therefore, we aim at estimating $\vartheta$ in $C\CC^{1+\alpha}(\rho^{3+\gamma'})$ where $0<\gamma'<\gamma$. This issue  will become even more challenging in the coming down from infinity in Section \ref{s:d}, where no time interpolation is available. 


With \eqref{eq:th43} at hand, \eqref{eq:43} rewrites as
\begin{align}\label{eq:rhs43a}
\begin{aligned}
0&= \LL\vartheta+\LL\psi -3\Big( (-\tthreeone{X}+\phi+\psi)\Prec\llbracket X^{2}\rrbracket- (-\tthreeone{X}+\phi+\psi)\prec\llbracket X^{2}\rrbracket\Big)\\
&\quad+3\llbracket X^2 \rrbracket\preccurlyeq(-\tthreeone{X}+\phi+\psi)-3[\LL,(-\tthreeone{X}+\phi+\psi)\Prec]\ttwoone{X}\\
&\quad+3X(-\tthreeone{X}+\phi+\psi)^2+(-\tthreeone{X}+\phi+\psi)^3+3b\varphi.
\end{aligned}
\end{align}
Next, we proceed with the same decomposition into regular (blue) and irregular (magenta) part as in Section \ref{s:phi}. It leads  to
\begin{align}\label{eq:two43}
\begin{aligned}
& \LL \phi + \rmm{\Phi} = 0,\quad\phi(0)=\phi_{0}=\varphi_{0},\\
&  \LL \psi + \psi^3 + \rmb{\tilde\Psi} = 0,\quad\psi(0)=0,
  \end{aligned}
 \end{align}
 where
 \begin{equation}\label{eq:two43aa}
 \rmb{\tilde\Psi}:=  \rmb{\Psi} -9\llbracket X^{2}\rrbracket \circ[(-\tthreeone{X}+\phi+\psi)\precprec \ttwoone{X}-(-\tthreeone{X}+\phi+\psi)\prec \ttwoone{X}],
 \end{equation}
 and $\rmm{\Phi}$, $\rmb{\Psi}$ are given exactly as in Section \ref{sec:45} (using the parabolic localizers $\VV_{\leq},\VV_{>}$ instead of the elliptic ones $\UU_{\leq},\UU_{>}$) with the same bounds applied pointwise in time, and $\varphi_{0}\in \CC^{1+\alpha}(\rho^{3+\gamma'}_{0})$ for some $\gamma'\in(0,\gamma)$ to be chosen below (the role of the parameter $\gamma$ is the same as in Section \ref{sec:45}, in particular, $\gamma=\alpha-\kappa$). The additional commutator in \eqref{eq:two43aa} is bounded as
 \begin{align*}
&\| 9\llbracket X^{2}\rrbracket \circ[(-\tthreeone{X}+\phi+\psi)\precprec \ttwoone{X}-(-\tthreeone{X}+\phi+\psi)\prec \ttwoone{X}]\|_{C\CC^{\gamma}(\rho^{3+\gamma})}\\
&\quad\lesssim \|\llbracket X^{2} \rrbracket\|_{C\CC^{-1-\kappa}(\rho^{\sigma})}\|-\tthreeone{X}+\phi+\psi\|_{C^{(\alpha+\kappa)/2}L^{\infty}(\rho^{3+\gamma'})} \|\ttwoone{X}\|_{C\CC^{1-\kappa}(\rho^{\sigma})}\\
&\quad\lesssim 1+\|\phi+\psi\|_{C^{(\alpha+\kappa)/2}L^{\infty}(\rho^{3+\gamma'})}.
 \end{align*}

 The equation for $\vartheta$ now reads as
 \begin{align*}%\label{eq:th43}
 \LL \vartheta + \Theta = 0,\quad\vartheta(0)=\phi_{0},
 \end{align*}
with
\begin{align*}
\Theta&=\rmm{\Phi}+\rmm{3\llbracket X^2 \rrbracket\succ\tthreeone{X}}-\rmm{3\VV_>\llbracket X^2 \rrbracket\succ(\phi+\psi)}+\rmb{3\VV_\leq\llbracket X^2 \rrbracket\succ(\phi+\psi)}\\
&\quad -3\Big( (-\tthreeone{X}+\phi+\psi)\Prec\llbracket X^{2}\rrbracket- (-\tthreeone{X}+\phi+\psi)\prec\llbracket X^{2}\rrbracket\Big)-3[\LL,(-\tthreeone{X}+\phi+\psi)\Prec]\ttwoone{X}.
\end{align*}
Here the two new terms are estimated using Lemma \ref{lem:5.1} as follows
\begin{align*}
&3\big\| (-\tthreeone{X}+\phi+\psi)\Prec\llbracket X^{2}\rrbracket- (-\tthreeone{X}+\phi+\psi)\prec\llbracket X^{2}\rrbracket\big\|_{C\CC^{-1+\alpha}(\rho^{3+\gamma'})}\\
& \quad\lesssim \|-\tthreeone{X}+\phi+\psi\|_{C^{(\alpha+\kappa)/2}L^{\infty}(\rho^{ 3+\gamma''})} \|\llbracket X^{2} \rrbracket\|_{C\CC^{-1-\kappa}(\rho^{\sigma})}\lesssim 1+\|\phi+\psi\|_{C^{(\alpha+\kappa)/2}L^{\infty}(\rho^{ 3+\gamma''})},
\end{align*}
\begin{align*}
&3\|[\LL,(-\tthreeone{X}+\phi+\psi)\Prec]\ttwoone{X}\|_{C\CC^{-1+\alpha}(\rho^{3+\gamma'})}\\
&\quad\lesssim \Big( \|-\tthreeone{X}+\phi+\psi\|_{C^{(\alpha+\kappa)/2}L^{\infty}(\rho^{3+\gamma''})}+\|-\tthreeone{X}+\phi+\psi\|_{C_{T}\CC^{\alpha+\kappa}(\rho^{3+\gamma''})}  \Big) \|\ttwoone{X}\|_{C\CC^{1-\kappa}(\rho^{\sigma})}\\
&\quad \lesssim 1 +  \|\phi+\psi\|_{C^{( \alpha+\kappa)/2}L^{\infty}(\rho^{ 3+\gamma''})}+\|\phi+\psi\|_{C\CC^{\frac 12 +\alpha}(\rho^{\frac 32+\alpha})},
\end{align*}
where $0<\gamma''<\gamma'$ and we chose $\alpha$ and $\kappa$ sufficiently small such that  $\alpha+\kappa<\delta$ for $\delta$ from Theorem \ref{thm:renorm43}.
All the other terms in $\Theta$ can be estimated pointwise in time by the same approach as in Section \ref{ssec:theta}.
Therefore, it only remains to bound the time regularity of $\phi+\psi$. Choosing $\gamma''$ sufficiently large and using Lemma \ref{lemma:interp2}, we obtain for $\lambda\in (0,1)$
\begin{align*}
\|\phi+\psi\|_{C^{( \alpha+\kappa)/2}L^{\infty}(\rho^{ 3+\gamma''})}&\lesssim \|\phi+\psi\|_{L^{\8}L^{\infty}(\rho)}+\lambda\|\phi+\psi\|_{C^{(\frac 12+ \alpha)/2}L^{\infty}(\rho^{ 3+\gamma})}\\
&\lesssim \|\phi+\psi\|_{L^{\8}L^{\infty}(\rho)}+\|\phi\|_{C^{(\frac 12+ \alpha)/2}L^{\infty}(\rho^{ \frac 32+\alpha})}+\lambda\|\psi\|_{C^{1}L^{\infty}(\rho^{  3+\gamma})}.
\end{align*}
Observe that the small parameter $\lambda$ is only needed in order to absorb the $C^{1}$-norm of $\psi$ into the left hand side.
And the same bound holds true for $\phi+\psi\in C^{( \alpha+\kappa)/2}L^{\infty}(\rho^{ 3+\gamma'})$ needed to estimate the additional term in $\rmb{\tilde\Psi}$.
This is sufficient in order to obtain the desired uniform bounds for $\phi,\vartheta,\psi$.

To be more precise, as in the case of the parabolic $\Phi^{4}$ model in dimension 2, we show existence via a smooth approximation and compactness. To this end, let $\xi_{\varepsilon}$ be a smooth and periodic approximation  of the driving space-time white noise $\xi$, defined on the torus of size $\frac 1\varepsilon$.
If  $\varphi_{\varepsilon,0}$ is a smooth approximation  of the initial condition $\varphi_{0}$, then according to Proposition \ref{prop:42d}, for every $\varepsilon\in (0,1)$ there exists $\varphi_{\varepsilon}\in C^{\infty}([0,T]\times \mathbb{T}^{3}_{1/\varepsilon})$ which is the unique classical solution to
\begin{equation*}%\label{eq:43d}
\LL \varphi_{\varepsilon} + \varphi_{\varepsilon}^3+ ( - 3 a_{\varepsilon} +3 b_{\varepsilon})\varphi_{\varepsilon} - \xi_{\varepsilon} = 0,\qquad\varphi(0)=\varphi_{\varepsilon,0}.
\end{equation*}
Now, we proceed with the same decomposition $\varphi_{\varepsilon}=X_{\varepsilon} - \tthreeone{X_{\varepsilon}} +\phi_{\varepsilon}+\psi_{\varepsilon}$ as above, only starting from the mollified noise $\xi_{\varepsilon}$ instead of  $\xi$. 
%
Next, in view of the bounds from Theorem \ref{thm:renorm43}, the a priori estimates from Sections~\ref{ssec:phi1}, \ref{ssec:phi2}, \ref{ssec:theta}, \ref{ssec:psi1}, \ref{ssec:psi2} apply mutatis mutandis  (using  Lemmas~\ref{lemma:interp2}, \ref{lem:local2}, \ref{lemma:sch}, \ref{lemma:schauder-par}, \ref{lemma:apriori-parabolic}), with only slight modification due to the necessary time regularity needed for the additional commutators in $\rmb{\tilde\Psi}$ and $\Theta$.


To summarize, we deduce that the following bounds hold true uniformly in $\varepsilon\in (0,1)$
\begin{align}\label{eq:un43}
\begin{aligned}
\|\phi_{\varepsilon}\|_{C\CC^{\alpha}(\rho)}+\|\phi_{\varepsilon}\|_{C^{\alpha/2}L^{\infty}(\rho)}+\|\phi_{\varepsilon}\|_{C\CC^{\frac 12+\alpha}(\rho^{\frac 32+\alpha})}+\|\phi_{\varepsilon}\|_{C^{(\frac 12+\alpha)/2}L^{\infty}(\rho^{\frac 32+\alpha})}&\lesssim 1,\\
\|\vartheta_{\varepsilon}\|_{C\CC^{1+\alpha}(\rho^{3+\gamma'})}+\|\vartheta_{\varepsilon}\|_{C^{(1+\alpha)/2}L^{\infty}(\rho^{3+\gamma'})}&\lesssim 1\\
\|\psi_{\varepsilon}\|_{C\CC^{2+\gamma}(\rho^{3+\gamma})}+\|\psi_{\varepsilon}\|_{C^{1}L^{\infty}(\rho^{3+\gamma})}+\|\psi_{\varepsilon}\|_{L^{\infty}L^{\infty}(\rho)}&\lesssim 1.
\end{aligned}
\end{align}




Consequently, we are able  to pass to the limit.

\begin{theorem}\label{thm:ex3d}
Let $\kappa, \alpha\in (0,1)$ be chosen sufficiently small and let $\gamma=\alpha-\kappa>0$ and $\gamma'\in(0,\gamma)$ sufficiently large. If $\varphi_{0}\in \CC^{1+\alpha}(\rho^{3+\gamma'}_{0})$ then there exist
\begin{align*}
\phi&\in C\CC^{\alpha}(\rho)\cap C^{\alpha/2}L^{\infty}(\rho)\cap C\CC^{\frac 12+\alpha}(\rho^{\frac 32+\alpha})\cap C^{(\frac 12+\alpha)/2}L^{\infty}(\rho^{\frac 32+\alpha}),\\
\vartheta&\in C\CC^{1+\alpha}(\rho^{3+\gamma'}) \cap C^{(1+\alpha)/2}L^{\infty}(\rho^{3+\gamma'}),\\
\psi&\in C\CC^{2+\gamma}(\rho^{3+\gamma})\cap C^{1}L^{\infty}(\rho^{3+\gamma})\cap L^{\infty}L^{\infty}(\rho),
\end{align*}
which is a solution to  \eqref{eq:two43}, \eqref{eq:th43}.
\end{theorem}

\begin{proof}
The proof follows the lines of Theorem \ref{thm:ex42}. Based on the uniform bounds from \eqref{eq:un43} we obtain compactness of the sequence of approximate solutions $(\phi_{\varepsilon},\vartheta_{\varepsilon},\psi_{\varepsilon})$ in a slightly worse space. In view of Theorem \ref{thm:renorm43}, this allows us to pass to the limit in the approximate version of \eqref{eq:two43}, \eqref{eq:th43}. Finally, we obtain that the limit solutions belong to the spaces where the uniform bounds hold.
\end{proof}














\section{Uniqueness for the parabolic models}
\label{s:uniq}


This  section is concerned with uniqueness to the parabolic
$\Phi^4_d$ model for $d = 2, 3$.

\begin{theorem}
  \label{thm:uniq3d}
The parabolic $\Phi^{4}$ model \eqref{eq:phi4p} in dimension 2 and 3 has a unique solution: Let $(\phi,\psi)$, $(\tilde\phi,\tilde\psi)$ be two solutions in the sense of Theorem \ref{thm:ex42} and Theorem \ref{thm:ex3d}, respectively, starting from an initial condition $\phi_{0}+\psi_{0}=\tilde\phi_{0}+\tilde\psi_{0}=\varphi_{0}$. Then $\phi+\psi=\tilde\phi+\tilde\psi$.
\end{theorem}


In what follows we present the proof of this result in the more involved setting of $d=3$. The two dimensional case follows the same pattern but is significantly easier and  the details are left  to the reader. Since the cubic term does not seem to be helpful for uniqueness, namely, when we study the equation for a difference of two solutions, it is necessary to find another mechanism which could handle the loss of weight in the terms of lower order. This issue can  be easily seen on the model equation
$$
\LL v = v f, \quad v(0)=0,
$$
which is the form the equation for the difference takes. Intuitively, if $f$ can only be bounded in a weighted space and accordingly also $v$ is bounded in a weighted space, then the product $vf$ can only be bounded when multiplied by the product of the two weights. Hence the right hand side requires higher weight than the left hand side which causes difficulties in closing the estimates.

We overcome this problem by introducing an exponential weight of the form $\rho(x)\pi(t,x):=\langle x \rangle^{- a} e^{- t \langle x
\rangle^b}$ for $a \in \mathbbm{R}$ and $b \in (0, 1)$. As usual $t\in [0,\8)$ denotes the time variable. However, this is not an admissible weight in the sense of Section \ref{ssec:besov} and consequently the definition of the associated weighted Besov spaces requires a different approach, either employing ultra-distributions (see \cite{MR891189}) or Gevrey classes (see \cite{MR1249275}). In Section \ref{ss:b} we recall the basic ideas  based on Gevrey classes following the detailed presentation of \cite{MW17}, where we also refer the reader for further
details. Note that exponential weights have already been employed in \cite{MR3358965, MR3779690}.

With suitable weighted Besov spaces at hand, we employ the classical $L^{2}$-energy technique. First, and similarly to the previous sections, we decompose the equation for the difference of two solutions into its regular and irregular components. Then we test both equations by a suitable test function, which corresponds to  the chain rule for certain Besov norm in the $L^{2}$-scale. This way we obtain a control of the $B_{2,2}^{\beta}(\pi_{t})$-norm of the regular component and the $B_{2,2}^{-\beta}(\pi_{t})$-norm of the irregular one, for some $\beta\in (0,1)$. The advantage of the exponential weight $\pi$ (which depends on time) originates in the form of its time derivative. More precisely, this gives a good term on the left hand side with weight of the form $\pi\rho^{-2b}$, that is, explosive at infinity in the space variable. This term is essential in order to control all the terms on the right hand side.




\subsection{Besov spaces with exponential weights}
\label{ss:b}

For the proof of Theorem \ref{thm:uniq3d}, we will employ weighted Besov
spaces with weights of the form $\rho(x)\pi(t,x):=\langle x \rangle^{- a} e^{- t \langle x
\rangle^b}$ for $a \in \mathbbm{R}$ and $b \in (0, 1)$ and $t\in [0,\8)$, which stands for the time variable. In order to compensate for the exponential growth, the definition of the
corresponding Besov spaces relies on the so-called Gevrey classes rather than
on Schwartz functions. Since multiplication by the polynomial weight $\langle
x \rangle^{- a}$ only introduces a logarithmic correction, namely, $\langle x
\rangle^{- a} e^{- t \langle x \rangle^b} = e^{- t \langle x \rangle^b - a
\log \langle x \rangle}$, we may work with the same Gevrey class
$\mathcal{G}^{\theta}$ of index $\theta \in (1, 1 / b)$ as for the case of only exponential weight $e^{- t \langle x
\rangle^b}$.
Consequently, the  results of \cite[Section~2]{MW17} remain valid and given $T>0$ the corresponding bounds are uniform over all $t\in [0,T]$.

Next,  we define the weighted Besov spaces (based on a partition of
unity from $\mathcal{G}^{\theta}$), as the completion of $C^{\infty}_c$ with
respect to the norm
\[ \| f \|_{B^{\alpha}_{p, q} (\pi_{t}\rho)} \assign \left( \sum_{k \geqslant - 1}
   (2^{\alpha k} \| \Delta_k f \|_{L^p (\pi_{t}\rho)})^q \right)^{\frac{1}{q}} =
   \left( \sum_{k \geqslant - 1} (2^{\alpha k} \| \pi_{t}\rho \Delta_k f \|_{L^p})^q
   \right)^{\frac{1}{q}}, \]
where $\pi_{t}(\cdot)=\pi(t,\cdot)$.
Note that unlike \cite{MW17} we pull the weight inside the $L^p$-norm, which is consistent with our definition of weighted Besov spaces in Section \ref{ssec:besov}. The corresponding results of \cite[Section 3.3]{MW17} (with straightforward modifications due to the weights) remain valid.
More precisely, the following paraproduct estimates will be used in the sequel.

\begin{lemma}
  \label{lem:par1}Let $\kappa \in [0, 1]$, $\beta \in \mathbbm{R}$ and $\delta > 0$. Then it
  holds uniformly in $t\geq 0$
  \[ \| f \prec g \|_{B_{2, 2}^{\beta} (\pi_{t} \rho)} \lesssim \| f \|_{L^2
     (\pi_{t})} \| g \|_{\CC^{\beta + \delta} (\rho)}\wedge \| f
     \|_{L^{\infty} (\rho)} \| g \|_{B^{\beta}_{2, 2} (\pi_{t})}, \]
  and if $\alpha < 0$ then uniformly over $t\geq 0$
    \[ \| f \prec g \|_{B_{2, 2}^{\alpha + \beta} (\pi_{t} \rho)} \lesssim \| f
     \|_{B^{\alpha}_{2, 2} (\pi_{t})} \| g \|_{\CC^{\beta+\delta} (\rho)} \wedge\| f
     \|_{\CC^{\alpha} (\rho)} \| g \|_{B^{\beta}_{2, 2} (\pi_{t})} . \]
  If $\alpha, \beta \in \mathbbm{R}$ such that $\alpha + \beta > 0$ then it
  holds uniformly in $t\geq 0$
  \[ \| f \circ g \|_{B_{2, 2}^{\alpha + \beta} (\pi_{t} \rho^{})} \lesssim \| f
     \|_{\CC^{\alpha} (\rho)} \| g \|_{B^{\beta}_{2, 2} (\pi_{t})} . \]
\end{lemma}

\begin{proof}
Let  $0<\gamma<\delta$. As a consequence of  \cite[Theorem 3.17]{MW17} and embeddings of Besov spaces, we have
  \[ \| f \prec g \|_{B_{2, 2}^{\beta} (\pi_{t} \rho)} \lesssim \| f \|_{B^{-
     \gamma}_{2, \infty} (\pi_{t})} \| g \|_{B^{\beta+\gamma}_{\infty, 2} (\rho)}
     \lesssim \| f \|_{L^2 (\pi_{t})} \| g \|_{B^{\beta + \delta}_{\infty, \infty}
     (\rho)} . \]
     So the first bound follows and similarly we obtain the third bound.
     The remaining bounds follow directly from  \cite[Theorem 3.17]{MW17}.
\end{proof}

Similarly to Lemma \ref{lem:com} we obtain the following result, whose  proof is a straightforward modification of \cite[Lemma 2.4]{GIP}.
\begin{lemma}
  \label{lem:com1}
  Let $\alpha \in (0, 1)$ and $\beta,
  \gamma \in \mathbbm{R}$ such that $\alpha + \beta + \gamma > 0$ and $\beta +
  \gamma < 0$. Then for every $\delta > 0$ it holds uniformly in $t\ge 0$
  \[ \| \tmop{com} (f, g, h) \|_{B^{\alpha + \beta + \gamma}_{2, 2} (\pi_{t}
     \rho_{1}\rho_{2})} \lesssim \| f \|_{B^{\alpha}_{2, 2} (\pi_{t})} \| g \|_{\CC^{\beta}
     (\rho_{1})} \| h \|_{\CC^{\gamma + \delta} (\rho_{2})} . \]
\end{lemma}


\subsection{Proof of Theorem \ref{thm:uniq3d}}

\begin{proof}
We prove the result for the case $d=3$.  Recall that if $\varphi$ is a solution to \eqref{eq:phi4p}  in the sense of  Theorem \ref{thm:ex3d} then $\varphi = X - \tthreeone{X} + \phi + \psi$, where $\phi$ is paracontrolled by $\ttwoone{X}$
 and equations \eqref{eq:two43}, \eqref{eq:th43} are satisfied.
For the purposes of the proof of uniqueness, it is not necessary to consider the modified paraproduct and therefore we may work with a similar decomposition as in the elliptic setting in Section \ref{sec:45}.
We define
  \[ \phi = \theta - 3 (- \tthreeone{X} + \phi + \psi) \prec
     \ttwoone{X} \]
  where (formally)
  \begin{align*}
      0 & = \LL \theta + \LL \psi  + 3 \llbracket X^2 \rrbracket \preccurlyeq (- \tthreeone{X} +
      \phi + \psi) - 3 [\LL, (- \tthreeone{X} + \phi + \psi) \prec]
      \ttwoone{X}\\
      & \quad + 3 X (- \tthreeone{X} + \phi + \psi)^2 + (- \tthreeone{X} +
      \phi + \psi)^3 + 3 b \varphi .
  \end{align*}
For notational simplicity, we chose  to write the above equation in this not rigorous  form -- with the infinite constant $b$ appearing -- instead of introducing the full decomposition with all the trees. Indeed,  we  are actually interested in a difference of the corresponding equations for two solutions $\varphi$ and $\tilde\varphi$ starting from the same initial condition. Thus the decomposition will simplify as the terms that do not depend on the solutions  cancel out.

  So if we denote by $\tilde{\varphi} = X - \tthreeone{X} + \tilde{\phi} +
  \tilde{\psi}$ another solution starting from the same initial condition and set $\zeta = \varphi - \tilde{\varphi}$
  and $\eta = \theta - \tilde{\theta} + \psi - \tilde{\psi}$, then we
  have
  \[ \zeta = \phi + \psi - \tilde{\phi} - \tilde{\psi} = - 3 \zeta \prec
     \ttwoone{X} + \eta . \]
  In addition, it holds
  \[ 0 = \LL \zeta + 3 \llbracket X^2 \rrbracket \zeta + 3 b \zeta + \Upsilon
     \zeta, \qquad \zeta(0)=0,\]
  \[ 0 = \LL \eta + 3 \llbracket X^2 \rrbracket \preccurlyeq
     \zeta + 3 b \zeta - 3 [\LL, \zeta \prec] \ttwoone{X} + \Upsilon
     \zeta,  \qquad \eta(0)=0,\]
  with
  \begin{align*}
      \Upsilon = & 3 X (- 2 \tthreeone{X} + \phi + \psi + \tilde{\phi} +
      \tilde{\psi}) \\
      & + [(- \tthreeone{X} + \phi + \psi)^2 + (- \tthreeone{X} + \phi +
      \psi) (- \tthreeone{X} + \tilde{\phi} + \tilde{\psi}) + (- \tthreeone{X}
      + \tilde{\phi} + \tilde{\psi})^2] .
  \end{align*}
  From the estimates of the solutions from Theorem \ref{thm:ex3d} we obtain $\Upsilon \in \CC^{-1/2- \kappa}
  (\rho^{\sigma})$. However, we stress that the term $3 \llbracket X^2
  \rrbracket \circ \zeta + 3 b \zeta$ is to be understood in the following sense
  \[ 3 \llbracket X^2 \rrbracket \circ \zeta + 3 b \zeta = - 9 \llbracket X^2
     \rrbracket \circ (\zeta \prec \ttwoone{X}) + 3 b \zeta + 3 \llbracket
     X^2 \rrbracket \circ \eta \]
  \[ =  - 9 \zeta \left( \llbracket X^2
     \rrbracket \circ \ttwoone{X} - \frac{b}{3} \right) - 9 \tmop{com} (\zeta,
     \ttwoone{X}, \llbracket X^2 \rrbracket) + 3 \llbracket X^2 \rrbracket
     \circ \eta, \]
  where $\ttwothreer{X}=\llbracket X^2 \rrbracket \circ \ttwoone{X} - \frac{b}{3}$ is the
  renormalized resonant product in $\CC^{- \kappa} (\rho^{\sigma})$.
  
  Now, we introduce a time dependent weight $\pi (t, x) = \exp (- t \rho^{- 2
  b} (x))$ where $\rho (x) = \langle x \rangle^{- 1}$ is a polynomial weight
  and $b \in (0, 1 / 2)$. Let $\beta, \gamma \in (0, 1)$ to be chosen below.
  First, we use the fact that $\Delta_k \LL = \LL \Delta_k$ in order to derive
  the equation for the Littlewood--Paley block $\Delta_k \eta$. Now, we test
  this equation by $\pi^2 \Delta_k \eta$ and use $\partial_t \pi = - \pi
  \rho^{- 2 b}$ to get
  \begin{equation}
    \frac{1}{2} \partial_t \| \Delta_k \eta \|_{L^2 (\pi)}^2 + \| \nabla
    \Delta_k \eta \|_{L^2 (\pi)}^2 + \mu \| \Delta_k \eta \|_{L^2 (\pi)}^2 +
    \| \rho^{- b} \Delta_k \eta \|_{L^2 (\pi)}^2 \label{eq:eta}
  \end{equation}
  \[ = \left\langle \pi^{} \Delta_k \eta, \pi \Delta_k \left( \LL \eta \right)
     \right\rangle - 2 \langle \pi \Delta_k \eta, \frac{\nabla \pi}{\pi} \pi
     \nabla \Delta_k \eta \rangle . \]
  Since $\nabla \pi = \pi t 2 b \rho^{- 2 b - 1} \nabla \rho$ and $| \nabla
  \rho / \rho^2 | \lesssim 1$, we obtain $\left| \frac{\nabla \pi}{\pi}
  \right| \lesssim t \rho^{1 - 2 b} \lesssim_T 1$ as a consequence of $b \in
  (0, 1 / 2)$ and $t \in [0, T]$. Hence by Young's inequality
  \[ | \langle \pi \Delta_k \eta, \frac{\nabla \pi}{\pi} \pi \nabla \Delta_k
     \eta \rangle | \leqslant C_{T, \delta} \|  \Delta_k \eta
     \|_{L^2 (\pi)}^2 + \delta \| \nabla \Delta_k \eta \|_{L^2 (\pi)}^2 . \]
  Now we estimate by duality, for some $\gamma \in (0, 1)$ and some $0 < a < b
  / 2$,
  \[ \left| \left\langle \pi^{} \Delta_k \eta, \pi \Delta_k \left( \LL \eta
     \right) \right\rangle \right| \lesssim 2^{- 2 \beta k} 2^{(2 \beta +
     \gamma) k} \| \Delta_k \eta \|_{L^2 (\pi \rho^{- a})} 2^{- \gamma k}
     \left\| \Delta_k \left( \LL \eta \right) \right\|_{L^2 (\pi \rho^a)} \]
  \[ \leqslant 2^{- 2 \beta k} \left( C_{\delta} 2^{2 (2 \beta + \gamma) k} \|
     \Delta_k \eta \|^2_{L^2 (\pi \rho^{- a})} + \delta 2^{- 2 \gamma k}
     \left\| \Delta_k \left( \LL \eta \right) \right\|^2_{L^2 (\pi \rho^a)}
     \right) . \]
 \rmbb{ Moreover, by a suitable choice of $a \in (0, b / 2)$ and $\beta, \gamma \in
  (0, 1)$ such that $\beta < 2 \beta + \gamma < \beta + 1$ (hence $\beta +
  \gamma < 1$), we may interpolate as follows. Let $\lambda\in (0,1)$ such that $a=(1-\lambda) b/2$ and $2\beta+\gamma=(1-\lambda)\beta+\lambda(\beta+1)$. Then H\"older's, Bernstein's and Young's  inequalities imply
  \[ 2^{2 (2 \beta + \gamma) k} \| \Delta_k \eta \|^2_{L^2 (\pi \rho^{-
     a})} \leqslant 2^{2(1-\lambda)\beta k} 2^{2\lambda(\beta+1)k} \|\rho^{-(1-\lambda)b/2} |\Delta_k \eta|^{1-\lambda} |\Delta_k \eta|^{\lambda} \|^2_{L^2 (\pi )}
     \]
\[
\leqslant 2^{2 ( 1-\lambda ) \beta k}  \| \rho^{- b/2} \Delta_k \eta \|^{2(1-\lambda)}_{ L^2 (\pi ) }
	2^{ 2 \lambda ( \beta+1 ) k}\|  \Delta_k \eta \|^{2\lambda}_{ L^2 (\pi ) }
\] 
\[
\leqslant 2^{2 ( 1-\lambda ) \beta k}  \| \rho^{- b/2} \Delta_k \eta \|^{2(1-\lambda)}_{ L^2 (\pi ) }
	2^{ 2 \lambda \beta k}\| \nabla \Delta_k \eta \|^{2\lambda}_{ L^2 (\pi ) }
\]     
\[
     \leqslant C_{\delta} 2^{2 \beta k} \| \rho^{- b / 2} \Delta_k
     \eta \|^2_{L^2 (\pi)} + \delta 2^{2 \beta k} \| \nabla \Delta_k \eta
     \|^2_{L^2 (\pi)} \]
  \[ \leqslant C_{\delta} 2^{2 \beta k} \| \Delta_k \eta \|^2_{L^2 (\pi)} +
     \delta 2^{2 \beta k} \| \rho^{- b} \Delta_k \eta \|^2_{L^2 (\pi)} +
     \delta 2^{2 \beta k} \| \nabla \Delta_k \eta \|_{L^2 (\pi)}^2 . \]}
  Similarly, we may test the equation for $\Delta_k \zeta$ by $\pi^2 \Delta_k
  \zeta$ to obtain
  \begin{equation}
    \frac{1}{2} \partial_t \| \Delta_k \zeta \|_{L^2 (\pi)}^2 + \| \nabla
    \Delta_k \zeta \|_{L^2 (\pi)}^2 + \mu \| \Delta_k \zeta \|_{L^2 (\pi)}^2 +
    \| \rho^{- b} \Delta_k \zeta \|_{L^2 (\pi)}^2 \label{eq:zeta}
  \end{equation}
  \[ = \left\langle \pi^{} \Delta_k \zeta, \pi \Delta_k \left( \LL \zeta
     \right) \right\rangle - 2 \langle \pi \Delta_k \zeta, \frac{\nabla
     \pi}{\pi} \pi \nabla \Delta_k \zeta \rangle . \]
  Next, we estimate
  \[ | \langle \pi \Delta_k \zeta, \frac{\nabla \pi}{\pi} \pi \nabla \Delta_k
     \zeta \rangle | \leqslant C_{T, \delta} \| \rho^{1 - 2 b} \Delta_k \zeta
     \|_{L^2 (\pi)}^2 + \delta \| \nabla \Delta_k \zeta \|_{L^2 (\pi)}^2 \]
  \[ \leqslant C_{T, \delta} \| \Delta_k \zeta \|_{L^2 (\pi)}^2 + \delta \|
    \nabla \Delta_k \zeta \|_{L^2 (\pi)}^2, \]
  and by duality
  \[ \left| \left\langle \pi^2 \Delta_k \zeta, \Delta_k \LL \zeta
     \right\rangle \right| \lesssim 2^{2 \beta k} 2^{(1 + \kappa - 2 \beta) k}
     \| \Delta_k \zeta \|_{L^2 (\pi \rho^{- c})} 2^{- (1 + \kappa) k} \left\|
     \Delta_k \left( \LL \zeta \right) \right\|_{L^2 (\pi \rho^c)} \]
  \[ \leqslant 2^{2 \beta k} \left( C_{\delta} 2^{2 (1 + \kappa - 2 \beta) k}
     \| \Delta_k \zeta \|^2_{L^2 (\pi \rho^{- c})} + \delta 2^{- 2 (1 +
     \kappa) k} \left\| \Delta_k \left( \LL \zeta \right) \right\|^2_{L^2 (\pi
     \rho^c)} \right), \]
  where by a suitable choice of $c \in (0, b / 2)$ we again interpolate to obtain
  \[ 2^{2 (1 + \kappa - 2 \beta) k} \| \Delta_k \zeta \|^2_{L^2 (\pi \rho^{-
     c})} \leqslant C_{\delta} 2^{- 2 \beta k} \| \rho^{- b / 2} \Delta_k
     \zeta \|^2_{L^2 (\pi)} + \delta 2^{- 2 \beta k} \| \nabla \Delta_k \zeta
     \|^2_{L^2 (\pi)} \]
  \[ \leqslant C_{\delta} 2^{- 2 \beta k} \| \Delta_k \zeta \|^2_{L^2
     (\pi)} + \delta 2^{- 2 \beta k} \| \rho^{- b} \Delta_k \zeta \|^2_{L^2
     (\pi)} + \delta 2^{- 2 \beta k} \| \nabla \Delta_k \zeta \|_{L^2
     (\pi)}^2, \]
  provided $- \beta < 1 + \kappa - 2 \beta < 1 - \beta$ hence $\kappa < \beta
  < 1 + \kappa$.
  
  As the next step, we multiply (\ref{eq:eta}) by $2^{2 \beta k}$,
  $(\ref{eq:zeta})$ by $2^{- 2 \beta k}$, integrate both inequalities over
  $(0, t)$ for some $t \in (0, T]$. In addition, we choose $\delta$
  sufficiently small  in order to absorb some of the
  terms into the left hand side. Finally, we sum the two inequalities, use the
  fact that $\zeta (0) = 0, \eta (0) = 0$, and sum over $k$ to obtain
  \[ \frac{1}{2} \| \eta (t) \|_{B^{\beta}_{2, 2} (\pi)}^2 + \frac{1}{2}\int_0^t \|
     \nabla \eta \|_{B^{\beta}_{2, 2} (\pi)}^2 \mathd s + \frac{1}{2} \| \zeta
     (t) \|_{B^{- \beta}_{2, 2} (\pi)}^2 + \frac{1}{2}\int_0^t \| \nabla \zeta \|_{B^{-
     \beta}_{2, 2} (\pi)}^2 \mathd s \]
  \begin{equation}\label{eq:s1}
    \le C_{\delta} \int_{0}^{t}\| \eta \|^2_{B^{\beta}_{2,2}
     (\pi)}\mathd s +C_{\delta} \int_{0}^{t}\| \zeta \|^2_{B^{-\beta}_{2,2}
     (\pi)}\mathd s 
     \end{equation}
     \begin{equation}+ \delta \int_0^t \left\| \LL \eta \right\|^2_{B^{- \gamma}_{2, 2}
    (\pi \rho^a)} \mathd s + \delta \int_0^t \left\| \LL \zeta \right\|^2_{B^{-
    (1 + \kappa)}_{2, 2} (\pi \rho^c)} \mathd s. \label{eq:s}
  \end{equation}
  We remark that in the above inequality we omitted the two terms containing $\rho^{-b}$ as well as the two terms containing $\mu$ on the left hand side since they are no longer needed.
  So it remains to control $\left\| \LL \eta \right\|^2_{B^{- \gamma}_{2, 2}
  (\pi \rho^a)}$ and $\left\| \LL \zeta \right\|^2_{B^{- (1 + \kappa)}_{2, 2}
  (\pi \rho^c)}$. Note that both terms are multiplied by a small constant
  $\delta > 0$, which will be needed in order to absorb them into the left hand side.
  
 We set $\beta = 2 \kappa$, $\gamma = \frac 12 +\kappa$ for some $\kappa > 0$
  sufficiently small, which is also the parameter to be used in order to
  estimate the stochastic terms according to the Table~\ref{t:reg}. Then $a=(1/2-3\kappa)b/2$ whereas $c=\kappa b/2$
hence  $c<a$, which will be used below in order to control the time derivative of $\zeta$.
  
 In view of Lemma \ref{lem:par1} and Lemma
  \ref{lem:com1}, we may estimate all the terms in $\LL \eta$ as follows
  \[ \| 3 \llbracket X^2 \rrbracket \prec \zeta (t) \|_{B^{- \gamma}_{2, 2}
     (\pi_{t} \rho^a)} \lesssim \| \zeta (t) \|_{B^{1 - \beta}_{2, 2} (\pi_{t})}, \]
  \[ \left\| 9 \zeta \left( \llbracket X^2 \rrbracket \circ \ttwoone{X} -
     \frac{b}{3} \right) (t) \right\|_{B^{- \gamma}_{2, 2} (\pi _{t}\rho^a)}
     \lesssim \| \zeta (t) \|_{B^{1 - \beta}_{2, 2} (\pi_{t})}, \]
  \[ \left\| 9 \tmop{com} (\zeta, \ttwoone{X}, \llbracket X^2 \rrbracket) (t)
     \right\|_{B^{- \gamma}_{2, 2} (\pi_{t} \rho^a)} \lesssim \| \zeta (t)
     \|_{B^{1 - \beta}_{2, 2} (\pi_{t})}, \]
  \[ \| 3 \llbracket X^2 \rrbracket \circ \eta (t) \|_{B^{- \gamma}_{2, 2}
     (\pi _{t}\rho^a)} \lesssim \| \eta (t) \|_{B^{1 + \beta}_{2, 2} (\pi_{t})}, \]
  \[ \left\| 3 [\LL, \zeta \prec] \ttwoone{X} (t)\right\|_{B^{-
     \gamma}_{2, 2} (\pi _{t}\rho^a)} \lesssim \left\| (\partial_{t}\zeta )\prec \ttwoone{X} (t)\right\|_{ B^{-
     \gamma}_{2, 2} (\pi _{t}\rho^a)}+\left\| (\Delta\zeta )\prec \ttwoone{X}(t) \right\|_{B^{-
     \gamma}_{2, 2} (\pi _{t}\rho^a)}\]
     \[
     +\left\| (\nabla\zeta )\prec \nabla\ttwoone{X} (t)\right\|_{ B^{-
     \gamma}_{2, 2} (\pi_{t} \rho^a)}\lesssim \| \partial_{t}\zeta (t)\|_{ B^{-1-\beta}_{2,2} (\pi_{t} \rho^{c})} + \| \zeta (t)\|_{ B^{1 -
     \beta}_{2, 2} (\pi_{t})}, \]
  \[ \| \Upsilon \zeta (t) \|_{B^{- \gamma}_{2, 2} (\pi_{t} \rho^a)} \lesssim \|
     \zeta (t) \|_{B^{1 - \beta}_{2, 2} (\pi_{t})} . \]
     We used the fact that $c<a$ to estimate the commutator above. 
  Besides, note that $\LL \zeta$ contains only one term which does not appear
  in $\LL \eta$ and it does not contain any term which requires time regularity. The additional term is controlled by
  \[ \| 3 \llbracket X^2 \rrbracket \succ \zeta (t) \|_{B^{- 1 - \kappa}_{2,
     2} (\pi_{t} \rho^c)} \lesssim \| \zeta (t) \|_{B^{1 - \beta}_{2, 2} (\pi_{t})} .
  \]
Using the equation for $\zeta$ we get
  \[ \| \partial_t \zeta(t) \|_{ B^{- 1 - \beta}_{2, 2} (\pi_{t} \rho^c)}
     \lesssim \left\| (\mu - \Delta) \zeta (t) - \LL \zeta (t) \right\|_{
     B^{- 1 - \beta}_{2, 2} (\rho^c \pi_{t})} \]
  \[ \lesssim \| \zeta (t)\|_{B^{1 - \beta}_{2, 2} (\pi_{t})} + \left\| \LL
     \zeta (t) \right\|_{B^{- 1 - \kappa}_{2, 2} (\pi_{t} \rho^c)}\lesssim \| \zeta(t) \|_{ B^{1 - \beta}_{2, 2} (\pi_{t})}  + \| \eta(t)
     \|_{ B^{1 + \beta}_{2, 2} (\pi_{t})}, \]
   where we applied the above estimates for $\LL \zeta$ again. Therefore, it
  follows
  \[ \left\| \LL \eta(t) \right\|_{ B^{- \gamma}_{2, 2} (\pi _{t}\rho^a)} +
     \left\| \LL \zeta (t)\right\|_{ B^{- 1 - \kappa}_{2, 2} (\pi_{t} \rho^c)}
     \lesssim \|\zeta(t) \|_{B^{1 - \beta}_{2, 2} (\pi_{t})} + \| \eta(t)
     \|_{ B^{1 + \beta}_{2, 2} (\pi_{t})} . \]
  Finally we have all in hand to conclude. Choosing $\delta$ sufficiently
  small allows to absorb the two terms in
  (\ref{eq:s}) into the left hand side and by Gronwall's lemma we absorb the terms in \eqref{eq:s1}. Accordingly $\zeta \equiv \eta \equiv
  0$ on $[0, T]$ and the proof is complete.
\end{proof}




\section{Coming down from infinity}
\label{s:d}

The goal of this section is to establish refined a priori estimates for solutions to the parabolic $\Phi^{4}$ model in dimension 2 and 3, which are valid independently of the initial condition. In particular this shows  that the solutions come down from infinity.

For the purposes of this section, we introduce a further time weight of the form $\tau (t) = 1 - e^{- t}$ so
that $\partial_t \tau (t) =  e^{- t} = 1 - \tau (t)$.
First of all, we prepare the initial data as follows. Let $\varphi_{0}\in \CC^{ - 1 + \varepsilon}
     (\rho^{1+\varepsilon}_0)$ for some $\varepsilon>0$.  Let $L>0$ be such that 
  \[ 2^{\varepsilon L} \sim \| \varphi_0 \|_{\CC^{ - 1 + \varepsilon}
     (\rho^{1+\varepsilon}_0)} . \]
Define 
$$\phi_{0}:=\UU_{>}
 \varphi_{0}-X(0),\qquad \psi_{0}:=\UU_{\leq}\varphi_{0},$$
 where $\UU_{>},\UU_{\leq}$ are the localizers corresponding to $L$. We recall that $X$ was chosen stationary and $X(0)\in \CC^{-1/2-\kappa}(\rho^{\sigma}_{0})$ for any $\sigma>0$  (see Theorem \ref{thm:renorm43}). Then it follows from Lemma \ref{lem:local} that
  \begin{equation}\label{eq:phi0}
  \|  \phi_0 \|_{\CC^{ - 1} (\rho_0)} \lesssim 1 
  \end{equation}
  uniformly over $\varphi_0 \in \CC^{ - 1 + \varepsilon} (\rho^{1+\varepsilon}_0)$ and $\varepsilon>0$.
Now we have all in hand to formulate the main result of this section.








\begin{theorem}
  \label{th:come-down-3}
  Let $\kappa, \alpha\in (0,1)$ be chosen sufficiently small and let $\gamma=\alpha-\kappa>0$. Let $\varphi_{0}\in \CC^{ - 1 + \varepsilon}
     (\rho^{1+\varepsilon}_0)$ for some $\varepsilon>0$. Let $(\phi, \psi)$ be a solution to the parabolic problem \eqref{eq:two43}, \eqref{eq:th43}
  in $d = 3$ with initial condition $(\phi_{0},\psi_{0})$ defined above.  
  Then, uniformly in $\varphi_{0}$,
  \[ \phi \in C \CC^{\alpha} (\tau^{\frac{1}{2}} \rho) \cap C \CC^{\frac{1}{2}
     + \alpha} ((\tau^{\frac{1}{2}} \rho)^{\frac{3}{2} + \alpha}), \]
  and
  \[ \psi \in C \CC^{2 + \gamma} ((\tau^{\frac{1}{2}} \rho)^{3 + \gamma}) \cap
     L^{\infty} L^{\infty} (\tau^{\frac{1}{2}} \rho) . \]
\end{theorem}


An analogous result holds true also in dimension 2. In this case, we construct the initial condition in the same way and obtain the following result.


\begin{theorem}
  \label{th:come-down-2}
  Let $\kappa, \alpha\in (0,1)$ be chosen sufficiently small and let $\beta=\alpha-\kappa>0$. Let $(\phi, \psi)$ a solution of the parabolic problem \eqref{eq:42b}
  in $d = 2$ with initial condition $(\phi_{0},\psi_{0})$ defined above. Then
  \[ (\phi, \psi) \in C \CC^{\alpha} (\tau^{\frac 12} \rho) \times [C \CC^{2
     + \beta} ((\tau^{ \frac 12} \rho)^{3 + \beta}) \cap L^{\infty}
     L^{\infty} (\tau^{\frac 12} \rho)] \]
  uniformly in the initial condition.
\end{theorem}


In the following, we discuss the necessary preliminary results and finally prove Theorem~\ref{th:come-down-3} in Section \ref{ss:thm91}. The proof of
Theorem~\ref{th:come-down-2} will not be given since it is substantially simpler
(as one does not need the paracontrolled ansatz) and follows the same pattern.
As a corollary,  we obtain that
\[ \| \phi (t) + \psi (t) \|_{L^{\infty}(\rho)} \lesssim 1 + t^{- 1 / 2} \]
independently of the initial condition: the solution comes down from infinity in a
finite time. This has been first observed by~\cite{MWcomedown} in the periodic setting.




\begin{remark}
We point out that the assumption on the regularity of initial condition in Theorems \ref{th:come-down-3}, \ref{th:come-down-2} is very weak and the existence for such singular initial conditions is not guaranteed  by the respective existence results, Theorems \ref{thm:ex3d}, \ref{thm:ex42}. However, for instance in case of the $\Phi^{4}$ model  \eqref{eq:phi4p} on $\mathbb{T}^{3}$, if $\varphi_{0}$ belongs to the natural space $\CC^{-1/2-\kappa}$, one may use the short time existence of a unique solution from \cite{CC} together with Theorem \ref{th:come-down-3} to deduce global existence and the coming down from infinity property. Furthermore, revisiting the proof of our a~priori estimates we see that the proportionality constants depend polynomially on the noise, which implies integrability of all the moments. This way, we recover the result of~\cite{MWcomedown}.
\end{remark}



\subsection{Interpolation and localization}
\label{ss:d0}


First, we notice that an interpolation similar to Lemma \ref{lemma:interp2} remains valid and the proof follows the same lines.

\begin{lemma}
  \label{lemma:tau-space-interpolation}For $\alpha \in [0, 2 + \kappa]$ we
  have
  \[ \| \tau^{ (1 + \alpha) / 2} \psi \|_{C \CC^{\alpha} (\rho^{1 + \alpha})}
     \lesssim \| \tau^{1 / 2} \psi \|_{C L^{\infty} (\rho)}^{1 - \alpha / (2 + 
     \kappa)} \| \tau^{(3 + \kappa) / 2} \psi \|_{C \CC^{2 + \kappa} (\rho^{3
     + \kappa})}^{\alpha / (2 + \kappa)}, \]
  or more generally
  \[ \| \tau^{ (1 + \alpha) / 2} \psi \|_{C \CC^{\alpha} (\rho^{1 + \alpha})}
     \lesssim \| \tau^{(1 + \delta) / 2} \psi \|_{C \CC^{\delta} (\rho^{1 +
     \delta})}^{1 - \theta} \| \tau^{(3 + \kappa) / 2} \psi \|_{C \CC^{2 +
     \kappa} (\rho^{3 + \kappa})}^{\theta}, \]
  whenever $\delta \in [0, \alpha]$ and $\alpha = (1 - \theta) \delta + \theta
  (2 + \kappa)$ for some $\theta \in [0, 1]$.
\end{lemma}

We stress that unlike Lemma \ref{lemma:interp2}, we do not include any interpolation in terms of time regularity into Lemma \ref{lemma:tau-space-interpolation}. Indeed, since the weight $\tau$ vanishes at zero, the equivalence  \eqref{eq:teq} is not valid anymore for the  corresponding weighted H\"older  spaces (in time). Therefore we proceed differently than in  Section \ref{sec:43}: below, we introduce a new modified paracontrolled ansatz which eventually leads us to the requirement $\tau^{\beta}(\phi+\psi)\in C^{{\delta}}L^{\infty}(\rho^{\sigma})$, for certain $\beta,\delta,\sigma>0$. In other words, instead of time regularity of $\phi+\psi$ in a space weighted by $\tau^{\beta}\rho^{\sigma}$, we require time regularity of $\tau^{\beta}(\phi+\psi)$ in a space weighted by an admissible space-time weight $\rho^{\sigma}$, which falls in the framework of Section \ref{ssec:besov}.


We will also need the parabolic localization \eqref{eq:locd} together with Lemma \ref{lem:local2}. However, since this is only applied to the stochastic objects that do not require any $\tau$ weight, no $\tau$-adapted version of Lemma \ref{lem:local2} is needed.




\subsection{Weighted Schauder estimates}
\label{ss:d1}


In this section we formulate new Schauder estimates  adapted to the particular weight $\tau$ which is not bounded away from zero. 
In particular, with the interpolation in hand, we may  employ Lemma~\ref{lemma:basic-tau-schauder} to deduce the following.


\begin{lemma}
  \label{lemma:tau-weighted-schauder}
   Let $\alpha
  > - 2$, $\gamma = (3 + \alpha) / 2$ and $\beta_i \in [0, 2)$. Assume that
  $\LL v = \sum_i V_i$. Then the following a priori estimate holds true
  \[ \| v \|_{C \CC^{2 + \alpha} (\tau^{\gamma} \rho)} \lesssim \| v \|_{C
     \CC^{\delta} (\tau^{(1 + \delta) / 2} \rho)} + \sum_i \| \tau^{\beta_i /
     2} V_i \|_{C \CC^{\alpha + \beta_i} (\tau^{\gamma} \rho)}, \]
  whenever $\delta \in [0, 1 + \alpha]$ is given by $1 + \alpha = (1 - \theta)
  \delta + \theta (2 + \alpha)$ for some $\theta \in [0, 1]$. Consequently, it
  also holds
  \[ \| v \|_{C \CC^{2 + \alpha} (\tau^{\gamma} \rho)} \lesssim \| v
     \|_{L^{\infty} L^{\infty} (\tau^{1 / 2} \rho)} + \sum_i \| \tau^{\beta_i
     / 2} V_i \|_{C \CC^{\alpha + \beta_i} (\tau^{\gamma} \rho)} . \]
\end{lemma}

\begin{proof}
First observe that
  \[ (\partial_t + \mu - \Delta) (\tau^{\gamma} \rho v) = \tau^{\gamma} \rho \sum_i V_i
     - \left( \frac{\Delta \rho}{\rho} \right) \tau^{\gamma} \rho v - 2
     \frac{\nabla \rho}{\rho} \nabla (\tau^{\gamma} \rho v) + 2 \left(
     \frac{\nabla \rho}{\rho} \right)^2 \tau^{\gamma} \rho v  \]
  \[ + \left(
     \frac{\partial_t \rho}{\rho} \right) \tau^{\gamma} \rho v+ \gamma \tau^{- 1} (1 - \tau) \tau^{\gamma} \rho v. \]
\rmbb{  By a slight modification of Lemma~\ref{lemma:basic-tau-schauder} where we choose different $\beta$ for different terms, we deduce
  $$
  \| \tau^{\gamma} \rho v \|_{C \CC^{2 + \alpha}}\lesssim \sum_i \| \tau^{\beta_i /
     2} V_i \|_{C \CC^{\alpha + \beta_i} (\tau^{\gamma} \rho)} + \| \gamma \tau^{- 1} (1 - \tau) \tau^{\gamma} \rho v
     \|_{C \CC^{2 +
     \alpha'} (\tau^{(2 - \alpha + \alpha') / 2})}
 $$
 $$  
 + \left\| -\left( \frac{\Delta \rho}{\rho} \right) \tau^{\gamma} \rho v - 2
     \frac{\nabla \rho}{\rho} \nabla (\tau^{\gamma} \rho v) + 2 \left(
     \frac{\nabla \rho}{\rho} \right)^2 \tau^{\gamma} \rho v +\left(
     \frac{\partial_t \rho}{\rho} \right) \tau^{\gamma} \rho v \right\|_{C \CC^{ \alpha}} 
  $$}
  \[  \lesssim  \sum_i \| \tau^{\beta_i /
     2} V_i \|_{C \CC^{\alpha + \beta_i} (\tau^{\gamma} \rho)} + \| \tau^{\gamma - 1} \rho v \|_{C \CC^{2 +
     \alpha'} (\tau^{(2 - \alpha + \alpha') / 2})} + \| \tau^{\gamma} \rho v
     \|_{C \CC^{1 + \alpha}}, \]
  where $\alpha < 2 + \alpha' < 2 + \alpha$. Now using  $\tau^{(3+\alpha)/2}\leq \tau^{(2+\alpha)/2}$ for the last term on the right hand side, we obtain by interpolation from
  Lemma~\ref{lemma:tau-space-interpolation} that 
  \[ \| v \|_{C \CC^{1 + \alpha} (\tau^{(3 + \alpha) / 2} \rho)} + \| v \|_{C
     \CC^{2 + \alpha'} (\tau^{(3 + \alpha') / 2} \rho)} \leqslant C \| v \|_{C
     \CC^{\delta} (\tau^{(1 + \delta) / 2} \rho)} + \frac{1}{2} \| v \|_{C
     \CC^{2 + \alpha} (\tau^{(3 + \alpha) / 2} \rho)} . \]
  This allows us to absorb the residual terms in the left hand side giving the
  final statement.
\end{proof}

Below we need also some specific Schauder estimate for time regularity of
solutions to the heat equation with a precise control of the $\tau$-weights in
the source term. We derive it here.

\begin{lemma}
  \label{lemma:time-tau-schauder}For any $\alpha \in (0, 2)$ and $\beta_i \in
  [0, 2)$ such that $\alpha + \beta_i - 2 < 0$ we have
  \[ \| v \|_{C^{\alpha / 2} L^{\infty} (\rho)} \lesssim \| v \|_{C
     \CC^{\alpha} (\rho)} + \sum_i \| \tau^{\beta_i / 2} V_i \|_{C \CC^{\alpha
     + \beta_i - 2} (\rho)} \]
  where \ $\LL v = \sum_i V_i$.
\end{lemma}

\begin{proof}
  Let $f = (\partial_t + \mu - \Delta) v$ and recall that we denoted by $P_{t}=e^{t(\Delta-\mu)}$ the  semigroup of operators generated by $\Delta-\mu$ with $\mu>0$. Fix $t>s \geqslant t-1$ and let $k \in \mathbbm{N}_0$ be such that $2^{- 2 k} \sim | t
  - s |$. Then using the fact that the weight $\rho$ is nonincreasing in time we obtain
  \[ \|v (t) - v (s) \|_{L^{\infty} (\rho_{t})} \lesssim \| \Delta_{\leqslant k}
     (P_{t - s} - \tmop{Id}) v (s) \|_{L^{\infty} (\rho_{s})} \]
  \[ + \int_s^t \| \Delta_{\leqslant k} P_{t - u} f (u) \|_{L^{\infty} (\rho_{u})}
     \mathd u +  \|\Delta_{>k} (v(t)-v(s)) \|_{ L^{\infty} (\rho_{t})} \backassign I_1
     + I_2 + I_3 \]
  Now
  $$
  I_{3}\lesssim \|\Delta_{>k}v(t) \|_{L^{\infty} (\rho_{t})} +\|\Delta_{>k}v(s) \|_{L^{\infty} (\rho_{s})} \lesssim 2^{- \alpha k} \| v \|_{C \CC^{\alpha} (\rho)},
  $$
  \[ I_1 \lesssim | t - s |^{\alpha / 2} \| v \|_{C \CC^{\alpha} (\rho)}, \]
  and if $t\le 2$ then
  \[ I_2 \lesssim \sum_i 2^{(2 - \alpha - \beta_i) k} \int_s^t \tau (u)^{-
     \beta_i / 2} \mathd u \| \tau^{\beta_i / 2} V_i \|_{C \CC^{\alpha +
     \beta_i - 2} (\rho)} \]
  \[ \lesssim \sum_i 2^{(2 - \alpha - \beta_i) k} 
     \int_s^t u^{- \beta_i / 2} \mathd u \| \tau^{\beta_i / 2} V_i \|_{C
     \CC^{\alpha + \beta_i - 2} (\rho)} \]
  \[ \lesssim | t - s |^{\alpha / 2} \left[ \sum_i \frac{(t^{1 - \beta_i / 2}
     - s^{1 - \beta_i / 2})}{| t - s |^{(1 - \beta_i / 2)}} \| \tau^{\beta_i /
     2} V_i \|_{C \CC^{\alpha + \beta_i - 2} (\rho)} \right], \]
Since the function $t\mapsto t^{\delta}$ is $\delta$-H\"older continuous, it holds true
  \[ \frac{(t^{1 - \beta_i / 2} - s^{1 - \beta_i / 2})}{| t - s |^{(1 -
     \beta_i / 2)}} \lesssim 1, \]
  hence
  \[ I_2 \lesssim | t - s |^{\alpha / 2} \sum_i \| \tau^{\beta_i / 2} V_i
     \|_{C \CC^{\alpha + \beta_i - 2} (\rho)} . \]
  If $t>2$ then $s\geq 1$ and consequently
  \[ I_2 \lesssim \sum_i 2^{(2 - \alpha - \beta_i) k} (t - s) \| \tau^{\beta_i
     / 2} V_i \|_{C \CC^{\alpha + \beta_i - 2} (\rho)} \]
  \[ \lesssim | t - s |^{\alpha / 2} \left[ \sum_i \frac{(t^{} - s^{})}{| t -
     s |^{(1 - \beta_i / 2)}} \| \tau^{\beta_i / 2} V_i \|_{C \CC^{\alpha +
     \beta_i - 2} (\rho)} \right] \lesssim | t - s |^{\alpha / 2} \sum_i \|
     \tau^{\beta_i / 2} V_i \|_{C \CC^{\alpha + \beta_i - 2} (\rho)} . \]
Hence we can conclude that
  \[ \sup_{\substack{s,t\in [0,\8)\\t-1\leq s < t}} \frac{\|v (t) - v (s) \|_{L^{\infty} (\rho_{t})}}{| t - s
     |^{\alpha / 2}} \lesssim \sum_i \| \tau^{\beta_i / 2} V_i \|_{C
     \CC^{\alpha + \beta_i - 2} (\rho)} + \| v \|_{C \CC^{\alpha} (\rho)} \]
and the claim follows.
\end{proof}

Similarly to Lemma \ref{lemma:schauder-par} we finally  derive a Schauder estimate for equations including a cubic nonlinearity.

\begin{lemma}
  \label{lemma:schauder-par1}
Fix $\kappa > 0$ and let $\psi \in C\CC^{2 +
  \kappa} (\rho^{3 + \kappa}) \cap C^{1}L^{\infty} (\rho^{3+\kappa})\cap L^{\infty}L^{\infty}(\rho)$ be a classical solution
  to
  \[ \partial_{t}\psi+(- \Delta + \mu) \psi + \psi^3 = \Psi ,\qquad \psi(0)=\psi_{0}.\]
Then
  \[ \| \psi \|_{C\CC^{2 + \kappa} ((\tau^{1/2}\rho)^{3 + \kappa})} \lesssim  1+ \|
     \Psi \|_{C\CC^{\kappa} ((\tau^{1/2}\rho)^{3 + \kappa})} + \| \psi \|_{L^{\infty}L^{\infty}
     (\tau^{1/2}\rho)}^{3 + \kappa}. \]
\end{lemma}


\begin{proof}
The proof follows from Lemma \ref{lemma:tau-space-interpolation} and Lemma \ref{lemma:tau-weighted-schauder} using the same approach as in the proof of Lemma \ref{lemma:schauder-par}.
\end{proof}





\subsection{Weighted  coercive estimate}
\label{ss:d2}

Next, we show that also the coercive estimates remain valid for the weight $\tau$, the proof uses the same ideas as Lemma \ref{lemma:apriori-parabolic}.

\begin{lemma}
  \label{lemma:weighted-coercive}
  Fix $\kappa > 0$ and let $\psi \in C\CC^{2 +
  \kappa} (\rho^{3 + \kappa}) \cap C^{1}L^{\infty} (\rho^{3+\kappa})\cap L^{\infty}L^{\infty}(\rho)$ be a classical solution
  to
  \[ \partial_t \psi + (\mu - \Delta) \psi + \psi^3 = \Psi, \qquad \psi (0) =
     \psi_0 . \]
  Then the following a priori estimate holds
  \[ \| \psi \|_{L^{\infty} L^{\infty} (\tau^{1 / 2} \rho)} \lesssim 1+ \|
     \Psi \|^{1 / 3}_{L^{\infty} L^{\infty} (\tau^{3 / 2} \rho^3)}  \]
     independently of the initial condition.
\end{lemma}

\begin{proof}
  Let $\bar{\psi} = \tau^{1 / 2} \psi \rho$ and assume for the moment that
  $\bar{\psi}$ attains its (global) maximum $M = \bar{\psi} (t^{\ast},
  x^{\ast})$ at the point $(t^{\ast}, x^{\ast})$. If $M \leq 0$, then it is
  necessary to investigate the minimum point (or alternatively the maximum of
  $- \bar{\psi}$), which we discuss below. Let us therefore assume that $M >
  0$. Then necessarily $t^{\ast} > 0$ since $\bar{\psi} (0) = 0$ and
  \[ \tau \rho^2 \partial_t  \bar{\psi} + \tau \rho^2 (- \Delta + \mu)
     \bar{\psi} + \bar{\psi}^3 = \tau^{3 / 2} \rho^3 \Psi + 
     (\tau \rho \partial_t \rho + \rho^2 \tau^{1 / 2} \partial_t \tau^{1 / 2})
     \bar{\psi} - \tau^{3 / 2} \rho^2  (\Delta \rho) \psi - 2 \tau^{3 / 2}
     \rho^2 \nabla \rho \nabla \psi . \]
  and
  \[ \partial_t  \bar{\psi} (t^{\ast}, x^{\ast}) = 0, \qquad \nabla \bar{\psi}
     (t^{\ast}, x^{\ast}) = 0, \qquad \Delta \bar{\psi} (t^{\ast}, x^{\ast})
     \leq 0 \]
  hence $\rho \nabla \psi = - \psi \nabla \rho$. Consequently $- \rho^2 \Delta
  \bar{\psi} (t^{\ast}, x^{\ast}) \geq 0$ and also $\rho \partial_t \rho
  \bar{\psi} (t^{\ast}, x^{\ast}) \leqslant 0$ since $\partial_t \rho
  \leqslant 0$. Hence
  \[ M^3 \leqslant [\tau^{3 / 2} \rho^3 \Psi ]_{|_{(t^{\ast}, x^{\ast})}} + [-\mu \tau\rho^2+\rho^2-\tau\rho\Delta\rho+2\tau |\nabla\rho|^{2}]_{|_{(t^{\ast}, x^{\ast})}} M \]
  \[ \leqslant \| \Psi \|_{L^{\infty} L^{\infty} (\tau^{3 / 2} \rho^3)} +
     c_{\rho, \mu} \| \bar{\psi} \|_{L^{\infty} L^{\infty}} . \]
  Therefore we deduce that
  \[ \bar{\psi} (t^{*},x^{*})= M \lesssim \| \Psi \|^{1 / 3}_{L^{\infty} L^{\infty}
     (\tau^{3 / 2} \rho^3)} + c_{\rho, \mu} \| \psi \|^{1 / 3}_{L^{\infty}
     L^{\infty} (\tau^{1 / 2} \rho)} . \]
  If $M < 0$ we can apply the same argument to $- \bar{\psi}$ to get
  \[ - \bar{\psi} \lesssim \| \Psi \|^{1 / 3}_{L^{\infty} L^{\infty} (\tau^{3
     / 2} \rho^3)} + c_{\rho, \mu} \| \psi \|^{1 / 3}_{L^{\infty} L^{\infty}
     (\tau^{1 / 2} \rho)} . \]
  hence
  \[ \| \psi \|_{L^{\infty} L^{\infty} (\tau^{1 / 2} \rho)} \lesssim \| \Psi
     \|^{1 / 3}_{L^{\infty} L^{\infty} (\tau^{3 / 2} \rho^3)} + c_{\rho, \mu}
     \| \psi \|^{1 / 3}_{L^{\infty} L^{\infty} (\tau^{1 / 2} \rho)} . \]
  and applying the weighted Young inequality yields the claim.


The conclusion in the case when $\tau^{1 / 2} \rho \psi$ does not
  attain its global maximum follows the same argument as in Lemma \ref{lemma:apriori-parabolic}.
\end{proof}


\begin{remark}\label{rem:down}
We note that the choice of $1/2$ as the power of the weight $\tau$ is dictated by the cubic nonlinearity through Lemma \ref{lemma:weighted-coercive}. More precisely, taking $\tau^{\alpha}$ instead of $\tau^{1/2}$ for some $\alpha>0$ and repeating the proof of the above maximum principle, it is necessary to control $\tau^{\alpha}\partial_{t}\tau^{\alpha}$ which leads to the condition $\alpha\geq 1/2$. Hence the choice $\alpha=1/2$ gives the best result.
\end{remark}


\subsection{Proof of Theorem \ref{th:come-down-3}}
\label{ss:thm91}

\begin{proof}
  Using the approach of Section \ref{sec:45} while choosing localization
  operators \eqref{eq:locd} according to the weighted norm $\| \phi + \psi \|_{L^{\infty}
  L^{\infty} (\tau^{1 / 2} \rho)}$ and the constant $L$ of various objects as in Table \ref{t:loc}, we obtain
     \begin{equation}\label{eq:Ph}
 \| \Phi \|_{C \CC^{- 2 + \alpha} (\tau^{(1 + \alpha) / 2} \rho)} \lesssim
     1.
     \end{equation}
The choice of the weight above is due to the following application of the weighted Schauder estimates from Lemma~\ref{lemma:tau-weighted-schauder}.   Namely, Lemma~\ref{lemma:tau-weighted-schauder} implies
  \begin{equation}
    \| \phi \|_{C \CC^{\alpha} (\tau^{(1 + \alpha) / 2} \rho^{})} \lesssim \|
    \phi \|_{L^{\infty} L^{\infty} (\tau^{1 / 2} \rho)} + \| \Phi \|_{C \CC^{-
    2 + \alpha} (\tau^{(1 + \alpha) / 2} \rho)} . \label{eq:down-phi}
  \end{equation}
  We note that this is not yet sufficient to bound $\phi$ in $L^{\infty}
  L^{\infty} (\tau^{1 / 2} \rho)$, since the power of $\tau$ on the left hand side of
  (\ref{eq:down-phi}) is bigger than $1 / 2$, which matters for $t$ small. In order to fill this gap, let $t\in (0,1)$ and
  $k \in \mathbbm{N}$ be such that $2^{- 2 k} \sim \lambda^{2 / \alpha} \tau (t)$
  for some small $\lambda > 0$ to be chosen below (the proportionality constants do not depend on time). Write
  \[ \| \tau^{(1 + \alpha) / 2} \phi (t) \|_{L^{\infty} (\rho_t)} \leqslant \|
     \tau^{(1 + \alpha) / 2} \Delta_{\leqslant k} \phi (t) \|_{L^{\infty}
     (\rho_t)} + \| \tau^{(1 + \alpha) / 2} \Delta_{> k} \phi (t)
     \|_{L^{\infty} (\rho_t)}, \]
  where $\rho_t (\cdot)$ denotes the weight $\rho (t, \cdot)$. For the first
  term on the right hand side we use the equation satisfied by $\phi$ together
  with the fact that the weight $\rho$ is nonincreasing in time. This gives
  \[ \| \tau^{(1 + \alpha) / 2} \phi (t) \|_{L^{\infty} (\rho_t)} \lesssim
     \tau^{(1 + \alpha) / 2} (t) \| \Delta_{\leqslant k} P_t \phi (0)
     \|_{L^{\infty} (\rho_0)} \]
  \[ + \tau^{(1 + \alpha) / 2} (t) \int_0^t  \|
     \Delta_{\leqslant k} P_{t - s} \Phi_{} (s) \|_{L^{\infty} (\rho_s)}
     \mathd s + \tau^{(1 + \alpha) / 2} (t) \| \Delta_{> k} \phi (t)
     \|_{L^{\infty} (\rho_t)} . \]
Recall  that $P_{t}=e^{t(\Delta-\mu)}$. Hence in view of \eqref{eq:sch10}, the definition of $\tau$ and the fact that $t\in(0,1)$, $\alpha - 2 < 0$ and $\alpha > 0$,  the above is further estimated
  by
  \[ \lesssim \tau^{(1 + \alpha) / 2} (t) 2^{k} \| \phi (0)
     \|_{\CC^{- 1} (\rho_0)} + 2^{- (\alpha - 2) k} \tau (
     t)\| \tau^{(1 + \alpha) / 2} \Phi_{} \|_{C \CC^{\alpha - 2}
     (\rho)} \]
  \[ + 2^{- \alpha k} \| \tau^{(1 + \alpha) / 2} \phi \|_{C \CC^{\alpha}
     (\rho)} . \]
  Using the definition of $k$ we therefore obtain
  \[ \| \tau^{(1 + \alpha) / 2} \phi (t) \|_{L^{\infty} (\rho_t)} \lesssim
     \tau (t)^{\alpha / 2} \frac{\tau (t)^{1 / 2}}{\tau (t)^{1 / 2}} \| \phi
     (0) \|_{\CC^{ - 1} (\rho_0)} + \tau (t)^{\alpha / 2} \frac{\tau
     (t)}{ \tau (t)} \| \tau^{(1 + \alpha) / 2} \Phi_{} \|_{C
     \CC^{\alpha - 2} (\rho)} \]
  \[ + \lambda \tau^{\alpha / 2} \| \tau^{(1 + \alpha) / 2} \phi \|_{C
     \CC^{\alpha} (\rho)}. \]
Hence we may divide  by $\tau^{\alpha/2}$ to obtain
  \[ \| \tau^{1 / 2} \phi (t) \|_{L^{\infty} (\rho_t)} \lesssim_{} \| \phi (0)
     \|_{\CC^{ - 1} (\rho_0)} + \| \tau^{(1 + \alpha) / 2} \Phi_{} \|_{C
     \CC^{\alpha - 2} (\rho)} + \lambda \| \tau^{(1 + \alpha) / 2} \phi \|_{C
     \CC^{\alpha} (\rho)} . \]
  Taking supremum in time and applying (\ref{eq:down-phi}) leads to
  \[ \| \phi \|_{L^{\infty} L^{\infty} (\tau^{1 / 2} \rho)} \lesssim_{} \|
     \phi (0) \|_{\CC^{ - 1} (\rho_0)} + \lambda^{} \| \phi
     \|_{L^{\infty} L^{\infty} (\tau^{1 / 2} \rho)} + \| \Phi \|_{C \CC^{- 2 +
     \alpha} (\tau^{(1 + \alpha) / 2} \rho)}. \]
Hence we can absorb the second term on the right hand side into the left hand side by choosing $\lambda$ sufficiently small.
Therefore, according to our construction of the initial datum $\phi(0)$, namely due to \eqref{eq:phi0}, and \eqref{eq:Ph}, \eqref{eq:down-phi}   we obtain
  \[ \| \phi \|_{L^{\infty} L^{\infty} (\tau^{1 / 2} \rho)} + \| \phi \|_{C
     \CC^{\alpha} (\tau^{(1 + \alpha) / 2} \rho^{})} \lesssim 1\]
     uniformly in the initial condition.
  Next, we apply Lemma \ref{lemma:tau-weighted-schauder} again, use the above
  $L^{\infty}$-bound for $\phi$ together with estimates similar to Section
  \ref{ssec:phi2} to obtain for some $\varepsilon\in (0,1)$
  \begin{equation}\label{eq:ph-down}
   \| \phi \|_{C \CC^{1 / 2 + \alpha} (\tau^{3 / 4 + \alpha / 2} \rho^{3 / 2
     + \alpha})} \lesssim 1 + \| \Phi \|_{C \CC^{- 3 / 2 + \alpha} (\tau^{3 /
     4 + \alpha / 2} \rho^{3 / 2 + \alpha})} \lesssim 1 + \| \psi
     \|_{L^{\infty} L^{\infty} (\tau^{1 / 2} \rho)}^{\varepsilon} .
     \end{equation}


As in Section \ref{sec:43}, the next step is a paracontrolled ansatz
  for $\phi .$ Here we have to be more careful due to the weight $\tau$. More
  precisely, we introduce a modified paracontrolled ansatz according to the
  formula
  \begin{equation}\label{eq:th2}
   \phi = \breve{\vartheta} - \tau^{- \frac{1 + \nu}{2}} \left(3 [\tau^{\frac{1 +
     \nu}{2}} (- \tthreeone{X} + \phi + \psi)] \precprec \ttwoone{X}\right)
  \end{equation}
  where $\nu > 0$ will be chosen later. This leads us to
  \[ 0 = \LL \breve{\vartheta} + \LL \psi - \left( \tau^{- \frac{1 + \nu}{2}}
     \left( 3 \left[ \tau^{\frac{1 + \nu}{2}} (- \tthreeone{X} + \phi + \psi)
     \right] \precprec \llbracket X^2\rrbracket \right) - 3 (- \tthreeone{X} + \phi + \psi)
     \prec \llbracket X^2\rrbracket \right) \]
  \[ + 3 \llbracket X^2\rrbracket \preccurlyeq (- \tthreeone{X} + \phi + \psi) \]
  \[ - \left( \LL \left[ \tau^{- \frac{1 + \nu}{2}} \left( \left[ 3
     \tau^{\frac{1 + \nu}{2}} (- \tthreeone{X} + \phi + \psi) \right]
     \precprec \ttwoone{X} \right) \right] - \tau^{- \frac{1 + \nu}{2}} \left(
     \left[ 3 \tau^{\frac{1 + \nu}{2}} (- \tthreeone{X} + \phi + \psi) \right]
     \precprec\llbracket X^2\rrbracket \right) \right) \]
  \[ + 3 X (- \tthreeone{X} + \phi + \psi)^2 + (- \tthreeone{X} + \phi +
     \psi)^3 + 3 b \varphi \]
  as a replacement of \eqref{eq:rhs43a}. Similarly
  to Section \ref{sec:43}, we obtain $\LL \breve{\vartheta} +
  \breve{\Theta} = 0$ with $\breve{\vartheta}(0)=\phi(0)$ and where $\breve{\Theta}$ differs from $\Theta$ only in the
  two commutators above, which we denote by $\mathrm{com_{1}},\,\mathrm{com_{2}}$ according to their order of appearance.
  
  Recall that in Section \ref{sec:45} and Section
  \ref{sec:43}, a suitable bound for $\vartheta$ was only needed in order to
  control the term $\llbracket X^2\rrbracket  \circ \vartheta$ in the equation for $\psi$. To
  this end, the choice of the weight $\rho^{2 + \alpha}$ in Section
  \ref{ssec:theta} was rather arbitrary, which has already been observed in Section \ref{sec:43}. Indeed, in order to control $\llbracket X^2\rrbracket 
  \circ \vartheta$ in $\CC^{\gamma} (\rho^{3 + \gamma})$ (cf. Section
  \ref{ssec:psi1}) it is sufficient to bound $\vartheta$ in
  $\CC^{1 + \alpha} (\rho^{3 + \gamma'})$ for some $0 < \gamma' < \gamma$.
 Similarly to Section \ref{sec:43}, we make use of this flexibility here: our goal is to apply Lemma
  \ref{lemma:tau-weighted-schauder} in order to estimate $\breve{\vartheta}$
  in $C \CC^{1 + \alpha}_{} (\tau^{(2 + \alpha) / 2} \rho_{}^{3 + \gamma'})$. This higher weight will allow us to control the term coming from $\Psi$ in the bound for time regularity below, see \eqref{eq:tregpsi}.
  
  
  Let us first estimate the new commutators appearing in $\breve{\Theta} .$ We
  rewrite the first commutator as
  \[ \tmop{com}_1= \tau^{- \frac{1 + \nu}{2}} \left( 3 \left[
     \tau^{\frac{1 + \nu}{2}} (- \tthreeone{X} + \phi + \psi) \right]
     \precprec \llbracket X^2\rrbracket - 3 \left[ \tau^{\frac{1 + \nu}{2}} (- \tthreeone{X} +
     \phi + \psi) \right] \prec \llbracket X^2\rrbracket \right), \]
  and observe that by Lemma \ref{lem:5.1} it can be estimated in $\CC^{- 1 +
  \alpha} (\rho_{}^{3 + \gamma'})$, using the time regularity of
  $\tau^{\frac{1 + \nu}{2}} (- \tthreeone{X} + \phi + \psi)$, provided we can control the blow-up in time by a suitable power of the $\tau$ weight. Hence, in view of Lemma
  \ref{lemma:tau-weighted-schauder} (with $\beta_i = 0$), which we aim to apply in order to gain the required regularity of $\breve{\vartheta}$, we  estimate
  (provided $2 + \alpha \geq 1 + \nu$)
  \[ \| \tmop{com}_1 \|_{C \CC^{- 1 + \alpha} (\tau^{(2 + \alpha) / 2}
     \rho_{}^{3 + \gamma'})^{}} \]
  \[ \lesssim \left\| \left[ \tau^{\frac{1 + \nu}{2}} (- \tthreeone{X} + \phi
     + \psi) \right] \precprec \llbracket X^2\rrbracket- 3 \left[ \tau^{\frac{1 + \nu}{2}} (-
     \tthreeone{X} + \phi + \psi) \right] \prec\llbracket X^2\rrbracket \right\|_{C \CC^{- 1 +
     \alpha} (\rho_{}^{3 + \gamma'})^{}} \]
  \[ \lesssim 1 + \left\| \tau^{\frac{1 + \nu}{2}} (\phi + \psi)
     \right\|_{C^{(\alpha + \kappa) / 2} L^{\infty} (\rho^{3 + \gamma''})} \]
  for some $0 < \gamma'' < \gamma'$. For the second commutator, it holds
  \[ \tmop{com}_2= - \frac{1 + \nu}{2} \tau^{- \frac{1 + \nu}{2} - 1}
     (1 - \tau) \left[ \left( 3 \tau^{\frac{1 + \nu}{2}} (- \tthreeone{X} +
     \phi + \psi) \right) \precprec \ttwoone{X} \right] \]
  \[ + \tau^{- \frac{1 + \nu}{2}} \left[ \LL, \left( 3 \tau^{\frac{1 +
     \nu}{2}} (- \tthreeone{X} + \phi + \psi) \right) \precprec \right]
     \ttwoone{X} \backassign \tmop{com}_{21} + \tmop{com}_{22} . \]
  The first term can be estimated in $\CC^{1 - \kappa}$ with a suitable
  weight. Hence Lemma \ref{lemma:tau-weighted-schauder} allows us to
  compensate for the blow up in $\tau$. More precisely, for this term we apply
  Lemma \ref{lemma:tau-weighted-schauder} with $\beta_i = 2 - \kappa - \alpha$
  to obtain (provided $4 - \kappa \geq 3 + \nu$)
  \[ \| \tau^{(2 - \kappa - \alpha) / 2} \tmop{com}_{21} \|_{C \CC^{1 -
     \kappa} (\tau^{(2 + \alpha) / 2} \rho_{}^{3 + \gamma'})^{}} \]
  \[ \lesssim \left\| \left( \tau^{\frac{1 + \nu}{2}} (- \tthreeone{X} + \phi
     + \psi) \right) \precprec \ttwoone{X} \right\|_{C \CC^{1 - \kappa}
     (\rho_{}^{3 + \gamma'})} \lesssim 1 + \| \tau^{1 / 2} (\phi + \psi)
     \|_{L^{\infty} L^{\infty} (\rho)} . \]
  The second term can be estimated using Lemma \ref{lem:5.1} in $\CC^{- 1 +
  \alpha} (\rho^{3 + \gamma'})$. So we again apply Lemma~\ref{lemma:tau-weighted-schauder} with $\beta_i = 0$ to deduce (since $2 +
  \alpha \geq 1 + \nu$)
  \[ \| \tmop{com}_{22} \|_{C \CC^{- 1 + \alpha} (\tau^{(2 + \alpha) / 2}
     \rho_{}^{3 + \gamma'})} \]
  \[ \lesssim 1 + \left\| \tau^{\frac{1 + \nu}{2}} (\phi + \psi)
     \right\|_{C^{(\alpha + \kappa) / 2} L^{\infty} (\rho^{3 + \gamma''})} +
     \left\| \tau^{\frac{1 + \nu}{2}} (\phi + \psi) \right\|_{C^{} \CC^{\alpha
     + \kappa} (\rho^{3 + \gamma''})} . \]
  All the other terms in $\breve{\Theta}$ can be estimated as in Section
  \ref{sec:43}, or more precisely pointwise in time by the approach of Section \ref{ssec:theta}. To summarize, Lemma \ref{lemma:tau-weighted-schauder} gives
  \[ \| \breve{\vartheta} \|_{C \CC^{1 + \alpha} (\tau^{(2 + \alpha) / 2}
     \rho^{3+\gamma'})} \lesssim \| \breve{\vartheta} \|_{C \CC^{1 / 2 +
     \alpha} (\tau^{3 / 4 + \alpha / 2} \rho^{3 / 2 + \alpha})} + 1 \]
  \[ + \| \psi \|_{C \CC^{1 / 2 + \alpha} (\tau^{3 / 4 + \alpha / 2} \rho^{3 /
     2 + \alpha})} + \| \psi \|_{L^{\infty} L^{\infty} (\tau^{1 / 2} \rho)}^{1
     + \varepsilon} \]
  \[ + \| \tau^{\frac{1 + \nu}{2}} (\phi + \psi) \|_{C^{(\alpha + \kappa)
     / 2} L^{\infty} (\rho^{3 + \gamma''})} + \left\| \tau^{\frac{1 + \nu}{2}}
     (\phi + \psi) \right\|_{C^{} \CC^{\alpha + \kappa} (\rho^{3 +
     \gamma''})}, \]
  where the only term requiring some time regularity is the second last one. This
  is the reason we introduced the modified $\vartheta$. In order to estimate
  the first term on the right hand side, we use the definition of
  $\breve{\vartheta}$ in \eqref{eq:th2} together with \eqref{eq:ph-down} to obtain (provided $3/2 + \alpha \geq 1+ \nu$)
  \[ \| \breve{\vartheta} \|_{C \CC^{1 / 2 + \alpha} (\tau^{3 / 4 + \alpha /
     2} \rho^{3 / 2 + \alpha})} \leqslant \| \phi \|_{C \CC^{1 / 2 + \alpha}
     (\tau^{3 / 4 + \alpha / 2} \rho^{3 / 2 + \alpha})} \]
  \[ + \left\| \tau^{- \frac{1 + \nu}{2}} (3 [\tau^{\frac{1 + \nu}{2}} (-
     \tthreeone{X} + \phi + \psi)] \precprec \ttwoone{X}) \right\|_{C \CC^{1 /
     2 + \alpha} (\tau^{3 / 4 + \alpha / 2} \rho^{3 / 2 + \alpha})} \]
  \[ \lesssim 1 + \| \psi \|_{L^{\infty} L^{\infty} (\tau^{1 / 2}
     \rho)}^{\varepsilon} + \left\| \tau^{- \frac{1+\nu}{2}} ([\tau^{\frac{1 +
     \nu}{2}} (- \tthreeone{X} + \phi + \psi)] \precprec \ttwoone{X})
     \right\|_{C \CC^{1 / 2 + \alpha} (\tau^{3 / 4 + \alpha / 2} \rho^{3 / 2 +
     \alpha})} \]
  \[ \lesssim 1 + \| \psi \|_{L^{\infty} L^{\infty} (\tau^{1 / 2}
     \rho)}^{\varepsilon} + \| \phi + \psi \|_{L^{\infty} L^{\infty} (\tau^{1
     / 2} \rho)} \lesssim 1 + \| \psi \|_{L^{\infty} L^{\infty}
     (\tau^{1 / 2} \rho)}^{} . \]
  Finally, it only remains to establish the time regularity of $\tau^{(1 +
  \nu) / 2} (\phi + \psi)$. A time-interpolation as in Lemma \ref{lemma:interp2} together with the Schauder estimates from
  Lemma~\ref{lemma:time-tau-schauder} (choosing the right regularity for each contributions and gaining
  powers of $\tau$)  yields for  $\delta\in(0,1)$
  \[ \| \tau^{\frac{1 + \nu}{2}} (\phi + \psi) \|_{C^{(\alpha + \kappa) / 2}
     L^{\infty} (\rho^{3 + \gamma''})}\lesssim  \| \tau^{\frac{1 + \nu}{2}} (\phi + \psi) \|_{L^{\8}     L^{\infty} (\rho)}+\delta \| \tau^{\frac{1 + \nu}{2}} (\phi + \psi) \|_{C^{(1/2+\alpha) / 2}
     L^{\infty} (\rho^{3 + \gamma})} \]
  \[ \lesssim  \| \phi + \psi \|_{L^{\infty} L^{\infty} (\tau^{1
     / 2} \rho)} +\delta \| \tau^{\frac{1 + \nu}{2}} (\phi + \psi) \|_{C \CC^{1/2+\alpha} (\rho^{3 + \gamma})} + \delta\| \LL [\tau^{\frac{1 + \nu}{2}} (\phi +
     \psi)]\|_{C \CC^{- 3/2 + \alpha} (\rho^{3 + \gamma})} \]
  \[ \lesssim \| \phi + \psi \|_{L^{\infty} L^{\infty} (\tau^{1
     / 2} \rho)} +\delta \| \tau^{\frac{1 + \nu}{2}} (\phi + \psi) \|_{C \CC^{1/2+\alpha} (\rho^{3/2 + \alpha})} \]
  \[ +\delta \frac{(1 + \nu)}{2} \| \tau^{(3/2 - \alpha
     ) / 2} \tau^{\frac{- 1 + \nu}{2}} (1 - \tau) (\phi + \psi) \|_{C
     L^{\infty} (\rho^{3 + \gamma})}+ \delta\|\tau^{\frac{1 + \nu}{2}} \Phi \|_{C \CC^{-
     3 / 2 + \alpha} (\rho^{3/2 + \alpha})} \]
  \begin{equation}\label{eq:tregpsi}
     + \delta\| \tau^{(3/2 + \gamma 
     - \alpha) / 2} \tau^{\frac{1 + \nu}{2}} \Psi \|_{C \CC^{\gamma}
     (\rho^{3 + \gamma})}+ \rmbb{\delta\| \tau^{(3/2  
     - \alpha) / 2} \tau^{\frac{1 + \nu}{2}} \psi^{3} \|_{C L^{\infty}
     (\rho^{3 + \gamma})} }. 
\end{equation}
 In fact, the small factor $\delta$ will only be needed to control $\Psi$ as it in turn also requires time regularity of $\tau^{(1+\nu)/2}(\phi+\psi)$ which needs to be absorbed into the left hand side (cf. \eqref{eq:two43aa}). Hence taking $\nu = 1 / 2 + \alpha$ we get (for a suitable choice of the
  parameters $\alpha, \kappa$)
\[
 \| \tau^{\frac{1 + \nu}{2}} (\phi + \psi) \|_{C^{(\alpha + \kappa) / 2}
     L^{\infty} (\rho^{3 + \gamma''})} \lesssim \|
     \phi + \psi \|_{L^{\infty} L^{\infty} (\tau^{1 / 2} \rho)}+ \| \phi + \psi \|_{C^{} \CC^{1
     / 2 + \alpha} (\tau^{3 / 4 + \alpha / 2} \rho^{3 / 2 + \alpha})}  \]
  \[ + \| \Phi \|_{C \CC^{- 3 / 2 + \alpha} (\tau^{3 / 4 + \alpha / 2} \rho^{3
     / 2 + \alpha})} +\delta \| \tau^{\frac{3 + \gamma}{2}} \Psi \|_{C
     \CC^{\gamma} ( \rho^{3 + \gamma})}+\rmbb{\| \tau^{3/2  
     }  \psi^{3} \|_{C L^{\infty}
     (\rho^{3 + \gamma})} } .
\]
On the other hand, the same estimates as in Section \ref{sec:43} (with further details in Section \ref{ssec:psi1}) lead to
  \[ \| \Psi \|_{C \CC^{\gamma} (\tau^{(3 + \gamma) / 2} \rho^{3 + \gamma})}
     \lesssim \| \psi \|_{C \CC^{1 + \alpha} (\tau^{1 + \alpha / 2} \rho^{2 +
     \alpha})} + \| \breve{\vartheta} \|_{C \CC^{1 + \alpha} (\tau^{1 + \alpha
     / 2} \rho^{3 +\gamma'})} \]
  \[ + \| \psi \|_{L^{\infty} L^{\infty} (\tau^{1 / 2} \rho)}^{1 +
     \varepsilon} \| \psi \|_{C \CC^{\gamma} (\tau^{1 / 2 + \gamma / 2}
     \rho^{1 + \gamma})} + \| \psi \|_{L^{\infty} L^{\infty} (\tau^{1 / 2}
     \rho)}^{2 + \varepsilon} \]
  \[ + (1 + \| \psi \|_{L^{\infty} L^{\infty} (\tau^{1 / 2} \rho)}) \| \psi
     \|_{C \CC^{1 / 2 + \alpha} (\tau^{3 / 4 + \alpha / 2} \rho^{3 / 2 +
     \alpha})} + 1 + \| \tau^{\frac{1 + \nu}{2}} (\phi + \psi) \|_{C^{(\alpha + \kappa) / 2}
     L^{\infty} (\rho^{3 + \gamma''})},\]
  which together with the bound for $\breve{\vartheta}$, $\phi$ above and choosing $\delta$ sufficiently small allows to control the time regularity as follows
\[ \| \tau^{\frac{1 + \nu}{2}} (\phi + \psi) \|_{C^{(\alpha + \kappa) / 2}
     L^{\infty} (\rho^{3 + \gamma''})} \lesssim 1+ \|
      \psi \|^{\rmbb{3}}_{L^{\infty} L^{\infty} (\tau^{1 / 2} \rho)}  \]
  \[+ \| \psi \|_{C \CC^{1 + \alpha} (\tau^{1 + \alpha / 2} \rho^{2 +
     \alpha})}  + \| \psi \|_{L^{\infty} L^{\infty} (\tau^{1 / 2} \rho)}^{1 +
     \varepsilon} \| \psi \|_{C \CC^{\gamma} (\tau^{1 / 2 + \gamma / 2}
     \rho^{1 + \gamma})} \]
  \[ + (1 + \| \psi \|_{L^{\infty} L^{\infty} (\tau^{1 / 2} \rho)}) \| \psi
     \|_{C \CC^{1 / 2 + \alpha} (\tau^{3 / 4 + \alpha / 2} \rho^{3 / 2 +
     \alpha})}.\]
This
  can be employed again in Lemma \ref{lemma:schauder-par1} to get
  \[ \| \psi \|_{C \CC^{2 + \gamma} (\tau^{(3 + \gamma) / 2} \rho^{3 +
     \gamma})} \lesssim 1+ \| \psi \|^{3+\gamma}_{L^{\infty} L^{\infty} (\tau^{1 / 2} \rho)}
     + \| \Psi \|_{C \CC^{\gamma_{}} (\tau^{(3 + \gamma) / 2} \rho^{3 +
     \gamma})} . \]
  As a consequence using also the weighted coercive estimate in
  Lemma~\ref{lemma:weighted-coercive} we can close our estimates (exactly as
  in Section \ref{sec:43}) and deduce that
  \[ \phi \in C \CC^{\alpha} (\tau^{\frac{1}{2}} \rho) \cap C \CC^{\frac{1}{2}
     + \alpha} ((\tau^{\frac{1}{2}} \rho)^{\frac{3}{2} + \alpha}), \qquad
     \breve{\vartheta} \in C \CC^{1 + \alpha}_{} (\tau^{(2 + \alpha) / 2}
     \rho_{}^{3 + \gamma'}), \]
  \[ \psi \in C \CC^{2 + \gamma} ((\tau^{\frac{1}{2}} \rho)^{3 + \gamma})  \cap
     L^{\infty} L^{\infty} (\tau^{\frac{1}{2}} \rho) . \]
  Since all the weighted data is zero at time zero, the estimates we obtain
  are uniform in the initial conditions.
\end{proof}







\appendix







\section{Auxiliary PDE results}
\label{s:a}

Here we show an auxiliary existence results needed in the main body of the paper.



\begin{proposition}\label{prop:aux-1}
Let $\Psi\in\CC^{\gamma}(\mathbb{T}^{d})$ for some $\gamma\in (0,1)$. There exists $\psi\in \CC^{2+\gamma}(\mathbb{T}^{d})$ which is a unique classical solution
\begin{equation}\label{eq:aux-1}
\Q \psi + \psi^3 + \Psi= 0.
\end{equation}
\end{proposition}

\begin{proof}
The  energy functional associated to the  equation in \eqref{eq:aux-1} reads as
\[ I (u) = \frac{1}{2} \int_{\mathbb{T}^{d}} | \nabla u |^2\, \dd x+ \frac{\mu}{2} \int_{\mathbb{T}^{d}} | u
   |^2 \, \dd x+ \frac{1}{4} \int_{\mathbb{T}^{d}} | u |^4 \, \dd x+ \int_{\mathbb{T}^{d}} \Psi u\, \dd x. \]
   It is differentiable on $W^{1,2}(\mathbb{T}^{d})\cap L^4(\mathbb{T}^{d})$ and
   $$
   I'(u)v=\int_{\mathbb{T}^{d}} \nabla u\cdot\nabla v\, \dd x+ \mu\int_{\mathbb{T}^{d}} uv\, \dd x +  \int_{\mathbb{T}^{d}}  u^3 v\, \dd x +\int\Psi v\, \dd x.
   $$
For $u, v\in W^{1,2}(\mathbb{T}^{d})\cap L^4(\mathbb{T}^{d})$ it holds
\[ (I' (u) - I' (v)) (u - v) = \int_{\mathbb{T}^{d}} \nabla (u - v)\cdot \nabla (u - v)\, \dd x
    + \mu \int_{\mathbb{T}^{d}} (u - v) (u - v) \, \dd x + \int_{\mathbb{T}^{d}} (u^3 - v^3) (u - v)\, \dd x
    \]
\[ = \| \nabla (u - v) \|_{L^2(\mathbb{T}^{d})}^2 + \mu \| u - v \|^2_{L^2(\mathbb{T}^{d})} + \int_{\mathbb{T}^{d}} (u - v
   )^2 (u^2 + u v + v^2 ) \geq 0 \]
   since $\mu>0$
and  $u^2 + u v + v^2 \geq 0$. In addition, if $u \neq v$ the above is
strictly positive and therefore $I$ is strictly convex on $W^{1,2}(\mathbb{T}^{d})\cap L^4(\mathbb{T}^{d})$ according to \cite[Proposition 1.5.10]{BS}. Moreover, it holds
$$
I(u)\geq \frac{1}{2}\|\nabla u\|_{L^2(\mathbb{T}^{d})}^2+\frac{\mu}{2}\|u\|^2_{L^2(\mathbb{T}^{d})}+\frac{1}{4}\|u\|_{L^4(\mathbb{T}^{d})}^4-\|\Psi\|_{L^2(\mathbb{T}^{d})}\|u\|_{L^2(\mathbb{T}^{d})}
$$
$$
\geq  \frac{1}{2}\|\nabla u\|_{L^2(\mathbb{T}^{d})}^2+\frac{\mu}{2}\|u\|^2_{L^2(\mathbb{T}^{d})}+\frac{1}{4}\|u\|_{L^4(\mathbb{T}^{d})}^4-c\|u\|_{L^2(\mathbb{T}^{d})}
$$
and consequently $I$ is coercive on $W^{1,2}(\mathbb{T}^{d})\cap L^4(\mathbb{T}^{d})$. Finally, if $u_n\to u$ in $W^{1,2}(\mathbb{T}^{d})\cap L^4(\mathbb{T}^{d})$ then $I(u_n)\to I(u)$ and hence $I$ is continuous on $W^{1,2}(\mathbb{T}^{d})\cap L^4(\mathbb{T}^{d})$.
Therefore, it follows from \cite[Theorem~1.5.6, Theorem~1.5.8]{BS} that $I$ has a unique minimum and  as a consequence   \eqref{eq:aux-1} possesses a unique weak solution in $W^{1,2}(\mathbb{T}^{d})\cap L^4(\mathbb{T}^{d})$.


Next, we show that  $\|\psi\|_{L^{\infty}(\mathbb{T}^{d})}\leqslant \|\Psi\|^{1/3}_{L^{\infty}(\mathbb{T}^{d})}$. To this end, let $R>0$ be such that $R^{3}=\| \Psi \|_{L^{\infty}(\mathbb{T}^{d})}$ and test the equation by $(\psi-R)_{+}$ to obtain
$$
- \int_{\mathbb{T}^{d}} (\psi-{R})_{+} \Delta\psi \,\dd x+ \mu\int_{\mathbb{T}^{d}} (\psi-R)_{+}\psi \,\dd x +\int_{\mathbb{T}^{d}} (\psi-R)_{+}\psi^{3}\,\dd x=\int_{\mathbb{T}^{d}} (\psi-R)_{+}\Psi\,\dd x.
$$
We rewrite this equation as
\begin{align}\label{eq:max}
\begin{aligned}
&- \int_{\mathbb{T}^{d}} (\psi-{R})_{+} \Delta\psi \,\dd x+ \mu\int_{\mathbb{T}^{d}} (\psi-R)_{+}\psi\,\dd x +\int_{\mathbb{T}^{d}} (\psi-R)_{+}(\psi^{3}-R^{3})\,\dd x\\
&\qquad =\int_{\mathbb{T}^{d}} (\psi-R)_{+}(\Psi -R^{3})\,\dd x
\end{aligned}
\end{align}
and estimate all the terms. The first term on the left hand side is nonnegative due to integration by parts
$$
- \int_{\mathbb{T}^{d}} (\psi-{R})_{+} \Delta\psi\,\dd x=\int_{\mathbb{T}^{d}} |\nabla (\psi-{R})_{+}|^{2}\,\dd x\geq 0 .
$$
As $\mu>0$, the linear term on the left hand side is nonnegative and the cubic term is also nonnegative since $\psi\geq R$ implies $\psi^{3}\geq R^{3}$.
Since also $\Psi\leqslant R^{3}$ due to the definition of $R$, the first term on the right hand side  of \eqref{eq:max} is nonpositive.
Hence we deduce that
$$
 \int_{\mathbb{T}^{d}} (\psi-{R})_{+}\psi \,\dd x\leqslant 0
$$
which further implies  $\psi\leqslant R$.
Applying the same approach to $-\psi$ yields $\psi\geq -R$ and the claim is proved.



Now, we include the cubic term $\psi^3$ into the right hand side and apply the Schauder estimates from \cite[(1.7)]{T06}. We obtain
$$
\|\psi\|_{W^{2,p}(\mathbb{T}^{d})}\lesssim \|(-\Delta+\mu)\psi\|_{L^{p}(\mathbb{T}^{d})} \leqslant \|\psi^3\|_{L^{p}(\mathbb{T}^{d})} +  \|\Psi\|_{L^{p}(\mathbb{T}^{d})},
$$
which is finite for all $p\in[1,\infty)$.  It follows that also  $\psi^3\in W^{2,p}(\mathbb{T}^{d})$ and due to the embedding $W^{2,p}(\mathbb{T}^{d})=F^{2}_{p2}(\mathbb{T}^{d})\hookrightarrow B^{2}_{p\infty}(\mathbb{T}^{d})\hookrightarrow B^{2-\delta}_{\infty\infty}(\mathbb{T}^{d})$ which holds true for all $\delta>0$ by choosing $p\in [1,\infty)$ sufficiently large   (see \cite[(1.3), (1.299), (1.305)]{T06}), we obtain that   $\psi^3\in \CC^{2-\delta}(\mathbb{T}^{d})$.
Thus,  the Schauder estimates \cite[(1.6)]{T06}  imply
$$
\|\psi\|_{\CC^{2+\gamma}(\mathbb{T}^{d})}\lesssim \|(-\Delta+\mu)\varphi\|_{\CC^{2+\gamma}(\mathbb{T}^{d})} \leqslant \|\psi^3\|_{\CC^{2-\delta}(\mathbb{T}^{d})} +  \|\Psi\|_{\CC^{\gamma}(\mathbb{T}^{d})}.
$$
Therefore, $\psi$ is a classical solution to   \eqref{eq:aux-1} and belongs to  $\CC^{2+\gamma}(\mathbb{T}^{d})$.
\end{proof}

\begin{proposition}\label{prop:42d}
Let $T>0$, $a\in C^{\infty}([0,T])$, $\xi\in C^{\infty}([0,T]\times\Td)$ and $\varphi_{0}\in C^{\infty}(\Td)$. There exists $\varphi\in C^{\infty}([0,T]\times \mathbb{T}^{d})$ which is the unique classical solution to
\begin{equation}\label{eq:42d}
(\partial_{t}-\Delta) \varphi +a\varphi + \varphi^3- \xi = 0,\qquad\varphi(0)=\varphi_{0}.
\end{equation}
\end{proposition}

\begin{proof}
The existence of a unique weak solution to \eqref{eq:42d} for initial conditions in $L^{2}(\mathbb{T}^{d})$ is classical and follows from monotonicity arguments applied within the Gelfand triplet $$[W^{1,2}(\mathbb{T}^{d})\cap L^{4}(\mathbb{T}^{d}) ]\hookrightarrow L^{2}(\mathbb{T}^{d})\hookrightarrow [W^{1,2}(\mathbb{T}^{d})\cap L^{4}(\mathbb{T}^{d})]^{*}.$$
The resulting weak solution $\varphi$ satisfies $\varphi\in C_{T}L^{2}(\mathbb{T}^{d})\cap L^{2}_{T}W^{1,2}(\mathbb{T}^{d})\cap L^{4}_{T}L^{4}(\mathbb{T}^{d})$. 



We test \eqref{eq:42d} by $\varphi^{2p-1}$ and apply the weighted Young inequality  to obtain
$$
\frac{1}{2p}\partial_{t}\int_{\mathbb{T}^{d}} |\varphi|^{2p}\,\dd x+(2p-1)\int_{\mathbb{T}^{d}} |\varphi|^{2p-2}|\nabla\varphi|^2\,\dd x+\int_{\mathbb{T}^{d}} |\varphi|^{2p+2}\,\dd x
$$
$$
\leqslant \int_{\mathbb{T}^{d}} \varphi^{2p-1}\xi\,\dd x+\|a\|_{L^{\8}}\int_{\mathbb{T}^{d}} |\varphi|^{2p}\,\dd x
$$
$$
\leqslant \kappa \int_{\mathbb{T}^{d}}|\varphi |^{2p+2}\,\dd x+c_{\kappa,p}\int_{\mathbb{T}^{d}}|\xi|^{\frac{2p+2}{3}}\,\dd x+\|a\|_{L^{\8}}\int_{\mathbb{T}^{d}} |\varphi|^{2p}\,\dd x
$$
for every $\kappa\in (0,1)$. Hence the Gronwall Lemma implies
$$
\frac{1}{2p}\partial_{t}\int_{\mathbb{T}^{d}} |\varphi|^{2p}\,\dd x+(2p-1)\int_{\mathbb{T}^{d}} |\varphi|^{2p-2}|\nabla\varphi|^2\,\dd x+\int_{\mathbb{T}^{d}} |\varphi|^{2p+2}\,\dd x\leqslant c_{T,p}
$$
and $\varphi\in L^{\infty}_{T}L^{2p}(\mathbb{T}^{d})$ for every $p\in\N$.
By interpolation, we deduce that $\varphi^3\in L^{\infty}_{T}L^{p}(\mathbb{T}^{d})$ for all $p\in [1,\infty)$. Hence we may include the  term $\varphi^{3}+a\varphi$ to the right hand side and apply a classical regularity result as for instance recalled in \cite[Theorem 3.1]{DDMH} to deduce that there exists $\alpha\in (0,1)$ and $p\in[1,\infty)$ such that
$$
\|\varphi\|_{\mathcal{C}^{\alpha/2,\alpha}}\lesssim \|\varphi_{0}\|_{\CC^{\alpha}(\mathbb{T}^{d})}+\|a\varphi\|_{L^{\infty}_{T}L^{p}(\mathbb{T}^{d})} + \|\varphi^3\|_{L^{p}_{T}L^{p}(\mathbb{T}^{d})} +  \|\xi\|_{L^{\infty}_{T}L^{p}(\mathbb{T}^{d})}\leqslant c_{\mu},
$$
where $\mathcal{C}^{\alpha/2,\alpha}=\mathcal{C}^{\alpha/2,\alpha}([0,T]\times \mathbb{T}^{d})$ denotes the parabolic H\"older space, that is, the H\"older space of order $\alpha$ with respect  to the parabolic distance $d((t,x),(s,y))=\max\{|t-s|^{1/2},|x-y|\}$. It is given by the norm
$$
\|f\|_{\mathcal{C}^{\alpha/2,\alpha}}=\sup_{(t,x)}|f(t,x)|+\sup_{(t,x)\neq(s,y)}\frac{|f(t,x)-f(s,y)|}{\max\{|t-s|^{\alpha},|x-y|^{\beta}\}}.
$$
Thus, it follows that $\varphi^{3}\in  \mathcal{C}^{\alpha/2,\alpha}$ and \cite[Theorem 3.4]{DDMH} yields that
$$
\|\varphi\|_{\mathcal{C}^{(\alpha+2)/2,\alpha+2}}\lesssim_{\mu} \|\varphi_{0}\|_{\CC^{\alpha+2}(\mathbb{T}^{d})}+\|a\varphi\|_{\mathcal{C}^{\alpha/2,\alpha}}+ \|\varphi^3\|_{\mathcal{C}^{\alpha/2,\alpha}} +  \|\xi\|_{\mathcal{C}^{\alpha/2,\alpha}}\leqslant c_{\mu},
$$
where the parabolic H\"older space  $\mathcal{C}^{(\alpha+k)/2,\alpha+k}=\mathcal{C}^{(\alpha+k)/2,\alpha+k}([0,T]\times \mathbb{T}^{d})$ for $\alpha\in (0,1)$ and $k\in\N$ is given by the norm
$$
\|f\|_{\mathcal{C}^{(\alpha+k)/2,\alpha+k}}=\sum_{r\in\N_{0},\gamma\in \N_{0}^{d};2r+|\gamma|\leqslant k}\|\partial^{r}_{t}\partial^{\gamma}f\|_{\mathcal{C}^{\alpha/2,\alpha}}.
$$
Since $\xi\in C^{\infty}([0,T]\times \mathbb{T}^{d})$, we may repeat the application of  \cite[Theorem 3.4]{DDMH} or also \cite[Chapter IV, Theorem 5.2]{LSU} to finally conclude that $\varphi\in C^{\infty}([0,T]\times \mathbb{T}^{d}).$
\end{proof}

In the following result we regard  functions on $\mathbb{T}^{d}_{1/\varepsilon}$  as periodic functions defined on the full space $\R^{d}$.

\begin{corollary}\label{cor:42d}
Let $\rho$ be a  space-time weight and let $\xi\in C^{\infty}([0,\8)\times\Td)$ be such that $\xi\in C\CC^{\kappa} (\rho^{3 + \kappa})\cap L^{\8}L^{\8}(\rho^{3})$ and $a\in C^{\8}([0,\8))\cap C^{1}_{b}([0,\infty))$. Let $\varphi\in C^{\infty}([0,\8)\times\Td)$  be the corresponding unique solution to \eqref{eq:42d} constructed in Proposition \ref{prop:42d}. Then $\varphi\in C\CC^{2 + \kappa} (\rho^{3 + \kappa})\cap C^{1}L^{\infty}(\rho^{3+\kappa}) \cap L^{\infty}L^{\infty} (\rho).$
\end{corollary}

\begin{proof}
The space and time regularity follows from Lemma \ref{lemma:schauder-par}. The proof of the $L^{\8}$-bound can be obtained by the same argument as in Lemma \ref{lemma:apriori-parabolic} applied on a finite interval $[0,T]$ and then sending $T\to\8$, since the proportionality constant does not depend on $T$.
\end{proof}




\section{Refined Schauder estimates}
\label{s:sch}

Here we establish a preliminary a priori bound which is needed   in Lemma \ref{lemma:tau-weighted-schauder}. Recall that $\tau$ is the time weight given by $\tau (t) = 1 - e^{- t}$.


\begin{lemma}
  \label{lemma:basic-tau-schauder}For any $\alpha \in \mathbbm{R}$ and
  $\beta_{} \in [0, 2)$ we have
  \[ \| v \|_{C \CC^{2 + \alpha}} \lesssim \left\| \tau^{\beta_{} / 2}
     \mathcal{} \LL v \right\|_{C \CC^{\alpha + \beta_{}}} \]
  provided $v (0) = 0$.
\end{lemma}

\begin{proof}
  Let $f = (\partial_t + \mu - \Delta) v$ and recall that we denoted by $P_{t}=e^{t(\Delta-\mu)}$ the  semigroup of operators generated by $\Delta-\mu$. Since $v (0) = 0$ it holds
  \[ \| \Delta_k v (t) \|_{L^{\infty}} \lesssim \int_0^t \| \Delta_k P_{t - s}
     f (s) \|_{L^{\infty}} \mathd s. \]
  Fix $k \geqslant - 1$. If $4 \geqslant t \geqslant 2^{- 2 k}$ we proceed to
  bound this quantity as follows
  \[ \| \Delta_k v (t) \|_{L^{\infty}} \lesssim 2^{- k (2 + \alpha - 2
     \varepsilon)} \int_0^{t - 2^{- 2 k}} (t - s)^{\beta / 2 - 1 +
     \varepsilon} \tau (s)^{- \beta / 2} \| (\tau^{\beta / 2} f) (s)
     \|_{\CC^{\alpha + \beta}} \mathd s \]
  \[ + 2^{- k (\alpha + \beta)} \int_{t - 2^{- 2 k}}^t \tau (s)^{- \beta / 2}
     \| (\tau^{\beta / 2} f) (s) \|_{\CC^{\alpha + \beta}} \mathd s \]
  \[ \lesssim 2^{- k (2 + \alpha + 2 \varepsilon)} \int_0^{t - 2^{- 2 k}} (t -
     s)^{\beta / 2 - 1 - \varepsilon} \tau (s)^{- \beta / 2} \mathd s \|
     \tau^{\beta / 2} f\|_{C \CC^{\alpha + \beta}} \]
  \[ + 2^{- k (\alpha + \beta)} \int_{t - 2^{- 2 k}}^t \tau (s)^{- \beta / 2}
     \mathd s \| \tau^{\beta / 2} f\|_{C \CC^{\alpha + \beta}} \]
  \[ \lesssim 2^{- k (2 + \alpha + 2 \varepsilon)} \int_0^{t - 2^{- 2 k}} (t -
     s)^{\beta / 2 - 1 - \varepsilon} s^{- \beta / 2} \mathd s \| \tau^{\beta
     / 2} f\|_{C \CC^{\alpha + \beta}} \]
  \[ + 2^{- k (\alpha + \beta)} \int_{t - 2^{- 2 k}}^t s^{- \beta / 2} \mathd
     s \| \tau^{\beta / 2} f\|_{C \CC^{\alpha + \beta}} . \]
  For the first integral we obtain
  \[ \int_0^{t - 2^{- 2 k}} (t - s)^{\beta / 2 - 1 - \varepsilon} s^{- \beta /
     2} \mathd s \]
  \[ \lesssim \int_{t / 2}^{t - 2^{- 2 k}} (t - s)^{\beta / 2 - 1 -
     \varepsilon} s^{- \beta / 2} \mathd s + \int_0^{t / 2} (t - s)^{\beta / 2
     - 1 - \varepsilon} s^{- \beta / 2} \mathd s \]
  \[ \lesssim t^{- \beta / 2} \int_{t / 2}^{t - 2^{- 2 k}} (t - s)^{\beta / 2
     - 1 - \varepsilon} \mathd s + t^{- \varepsilon} \int_0^{1 / 2} (1 -
     s)^{\beta / 2 - 1 - \varepsilon} s^{- \beta / 2} \mathd s \]
  \[ \lesssim t^{- \beta / 2} 2^{2 k (\varepsilon - \beta / 2)} + 2^{2 k
     \varepsilon} \lesssim 2^{2 k \varepsilon}, \]
  whereas for the second one
  \[ \int_{t - 2^{- 2 k}}^t s^{- \beta / 2} \mathd s \lesssim 2^{- 2 k} t^{-
     \beta / 2} \lesssim 2^{- 2 k (1 - \beta / 2)} . \]
  Hence, this leads to the desired bound
  \[ \| \Delta_k v (t) \|_{L^{\infty}} \lesssim 2^{- k (2 + \alpha)} \|
     \tau^{\beta / 2} f\|_{C \CC^{\alpha + \beta}} . \]
  If $0 < t \leqslant 2^{- 2 k}$ then
  \[ \| \Delta_k v (t) \|_{L^{\infty}} \lesssim \int_0^t \| \Delta_k P_{t - s}
     f (s) \|_{L^{\infty}} \mathd s \lesssim 2^{- k (\alpha + \beta)} \int_0^t
     \tau (s)^{- \beta / 2} \mathd s \| \tau^{\beta / 2} f\|_{C \CC^{\alpha +
     \beta}} \]
  \[ \lesssim 2^{- k (\alpha + \beta)} t^{1 - \beta / 2} \| \tau^{\beta / 2}
     f\|_{C \CC^{\alpha + \beta}} \lesssim 2^{- k (\alpha + \beta)} 2^{- 2 k
     (1 - \beta / 2)} \| \tau^{\beta / 2} f\|_{C \CC^{\alpha + \beta}} \]
  \[ \lesssim 2^{- k (2 + \alpha)} \| \tau^{\beta / 2} f\|_{C \CC^{\alpha +
     \beta}} . \]
  Finally, for $t > 4$ we write
  \[ \| \Delta_k v (t) \|_{L^{\infty}} \lesssim \int_0^{1 / 2} \| \Delta_k
     P_{t - s} f (s) \|_{L^{\infty}} \mathd s + \int_{1 / 2}^t \| \Delta_k
     P_{t - s} f (s) \|_{L^{\infty}} \mathd s = I_1 + I_2, \]
  and estimate
  \[ I_1 \lesssim 2^{- k (2 + \alpha)} \int_0^{1 / 2} (t - s)^{\beta / 2 - 1}
     \tau (s)^{- \beta / 2} \mathd s \| \tau^{\beta / 2} f \|_{C \CC^{\alpha +
     \beta}} \]
  \[ \lesssim 2^{- k (2 + \alpha)} \int_0^{1 / 2} s^{- \beta / 2} \mathd s \|
     \tau^{\beta / 2} f \|_{C \CC^{\alpha + \beta}} \lesssim 2^{- k (2 +
     \alpha)} \| \tau^{\beta / 2} f \|_{C \CC^{\alpha + \beta}}, \]
  \[ I_2 \lesssim \int_{1 / 2}^t \| \tau^{\beta / 2} \Delta_k P_{t - s} f (s)
     \|_{L^{\infty}} \mathd s \lesssim 2^{- k (2 + \alpha)} \| \tau^{\beta /
     2} f \|_{C \CC^{\alpha + \beta}}, \]
  where the last term was estimated as in the standard Schauder estimates. The
  conclusion follows.
\end{proof}



%\bibliographystyle{alpha}
%\bibliography{elliptic}

\begin{thebibliography}{DdMH15}

\bibitem[AK17]{albeverio_invariant_2017}
Sergio Albeverio and Seiichiro Kusuoka.
\newblock The invariant measure and the flow associated to the
  {$\Phi^4_3$}-quantum field model.
\newblock {\em arXiv:1711.07108}, November 2017.


\bibitem[AB06]{AB06}
C. D. Aliprantis and Kim C. Border.
\newblock Infinite dimensional analysis. A hitchhiker?s guide.
\newblock Springer, Berlin, Third edition, 2006. 

\bibitem[BCD11]{BCD}
Hajer Bahouri, Jean-Yves Chemin, and Rapha{\"e}l Danchin.
\newblock {\em Fourier {Analysis} and {Nonlinear} {Partial} {Differential}
  {Equations}}.
\newblock Springer, January 2011.

\bibitem[BS11]{BS}
Marino Badiale and Enrico Serra.
\newblock {\em Semilinear elliptic equations for beginners}.
\newblock Universitext. Springer, London, 2011.
\newblock Existence results via the variational approach.

\bibitem[CC18]{CC}
R.~Catellier and K.~Chouk.
\newblock Paracontrolled distributions and the 3-dimensional stochastic
  quantization equation.
\newblock {\em to appear in The Annals of Probability}, 2018.

\bibitem[DdMH15]{DDMH}
Arnaud Debussche, Sylvain de~Moor, and Martina Hofmanov\'a.
\newblock A regularity result for quasilinear stochastic partial differential
  equations of parabolic type.
\newblock {\em SIAM J. Math. Anal.}, 47(2):1590--1614, 2015.

\bibitem[Eva10]{E98}
Lawrence~C. Evans.
\newblock {\em Partial {Differential} {Equations}: {Second} {Edition}}.
\newblock American Mathematical Society, 2 edition, March 2010.

\bibitem[FG17]{furlan_weak_2017}
Marco Furlan and Massimiliano Gubinelli.
\newblock Weak universality for a class of 3d stochastic reaction-diffusion
  models.
\newblock {\em arXiv:1708.03118 [math]}, August 2017.
\newblock arXiv: 1708.03118.

\bibitem[GIP15]{GIP}
Massimiliano Gubinelli, Peter Imkeller, and Nicolas Perkowski.
\newblock Paracontrolled distributions and singular {PDEs}.
\newblock {\em Forum of Mathematics. Pi}, 3:e6, 75, 2015.

\bibitem[GP17]{gubinelli_kpz_2017}
Massimiliano Gubinelli and Nicolas Perkowski.
\newblock {KPZ} {Reloaded}.
\newblock {\em Communications in Mathematical Physics}, 349(1):165--269,
  January 2017.

\bibitem[Hai14]{hairer_theory_2014}
M.~Hairer.
\newblock A theory of regularity structures.
\newblock {\em Inventiones mathematicae}, 198(2):269--504, March 2014.

\bibitem[HL15]{MR3358965}
Martin Hairer and Cyril Labb\'e.
\newblock A simple construction of the continuum parabolic {A}nderson model on
  {${\bf R}^2$}.
\newblock {\em Electron. Commun. Probab.}, 20:no. 43, 11, 2015.

\bibitem[HL18]{MR3779690}
Martin Hairer and Cyril Labb\'e.
\newblock Multiplicative stochastic heat equations on the whole space.
\newblock {\em J. Eur. Math. Soc. (JEMS)}, 20(4):1005--1054, 2018.

\bibitem[KFP83]{klein_supersymmetry_1983}
Abel Klein and J.~Fernando~Perez.
\newblock Supersymmetry and dimensional reduction: {A} non-perturbative proof.
\newblock {\em Physics Letters B}, 125(6):473--475, June 1983.

\bibitem[KLP84]{klein_supersymmetry_1984}
Abel Klein, Lawrence~J. Landau, and J.~Fernando Perez.
\newblock Supersymmetry and the {Parisi}-{Sourlas} dimensional reduction: {A}
  rigorous proof.
\newblock {\em Communications in Mathematical Physics}, 94(4):459--482,
  December 1984.

\bibitem[Kup16]{kupiainen_renormalization_2016}
Antti Kupiainen.
\newblock Renormalization {Group} and {Stochastic} {PDEs}.
\newblock {\em Annales Henri Poincar{\'e}}, 17(3):497--535, March 2016.

\bibitem[LSU68]{LSU}
O.~A. Ladyzhenskaya, V.~A. Solonnikov, and N.~N. Ural'tseva.
\newblock {\em Linear and quasilinear equations of parabolic type}.
\newblock Translated from the Russian by S. Smith. Translations of Mathematical
  Monographs, Vol. 23. American Mathematical Society, Providence, R.I., 1968.

\bibitem[MW17a]{MWcomedown}
Jean-Christophe Mourrat and Hendrik Weber.
\newblock The dynamic {$\Phi^4_3$} model comes down from infinity.
\newblock {\em Comm. Math. Phys.}, 356(3):673--753, 2017.

\bibitem[MW17b]{MW17}
Jean-Christophe Mourrat and Hendrik Weber.
\newblock Global well-posedness of the dynamic {$\Phi^4$} model in the plane.
\newblock {\em The Annals of Probability}, 45(4):2398--2476, July 2017.

\bibitem[MWX16]{mourrat_construction_2016}
Jean-Christophe Mourrat, Hendrik Weber, and Weijun Xu.
\newblock Construction of \${\textbackslash}{Phi}{\textasciicircum}4\_3\$
  diagrams for pedestrians.
\newblock {\em arXiv:1610.08897 [math-ph]}, October 2016.
\newblock arXiv: 1610.08897.

\bibitem[OW16]{otto_quasilinear_2016}
Felix Otto and Hendrik Weber.
\newblock Quasilinear {SPDEs} via rough paths.
\newblock {\em arXiv:1605.09744 [math]}, May 2016.
\newblock arXiv: 1605.09744.

\bibitem[PS79]{parisi_random_1979}
G.~Parisi and N.~Sourlas.
\newblock Random {Magnetic} {Fields}, {Supersymmetry}, and {Negative}
  {Dimensions}.
\newblock {\em Physical Review Letters}, 43(11):744--745, September 1979.

\bibitem[PW81]{parisi_perturbation_1981}
G.~Parisi and Yong~Shi Wu.
\newblock Perturbation theory without gauge fixing.
\newblock {\em Scientia Sinica. Zhongguo Kexue}, 24(4):483--496, 1981.

\bibitem[Rod93]{MR1249275}
Luigi Rodino.
\newblock {\em Linear partial differential operators in {G}evrey spaces}.
\newblock World Scientific Publishing Co., Inc., River Edge, NJ, 1993.

\bibitem[Sim87]{S87}
Jacques Simon.
\newblock Compact sets in the space {$L^p(0,T;B)$}.
\newblock {\em Ann. Mat. Pura Appl. (4)}, 146:65--96, 1987.

\bibitem[ST87]{MR891189}
Hans-J\"urgen Schmeisser and Hans Triebel.
\newblock {\em Topics in {F}ourier analysis and function spaces}.
\newblock A Wiley-Interscience Publication. John Wiley \& Sons, Ltd.,
  Chichester, 1987.

\bibitem[Tri92]{T92}
Hans Triebel.
\newblock {\em Theory of {Function} {Spaces} {II}}.
\newblock Springer, August 1992.

\bibitem[Tri06]{T06}
Hans Triebel.
\newblock {\em Theory of {Function} {Spaces} {III}}.
\newblock Springer, August 2006.

\end{thebibliography}

\end{document}













